% Options for packages loaded elsewhere
\PassOptionsToPackage{unicode}{hyperref}
\PassOptionsToPackage{hyphens}{url}
%
\documentclass[
  openany]{book}
\usepackage{lmodern}
\usepackage{amsmath}
\usepackage{ifxetex,ifluatex}
\ifnum 0\ifxetex 1\fi\ifluatex 1\fi=0 % if pdftex
  \usepackage[T1]{fontenc}
  \usepackage[utf8]{inputenc}
  \usepackage{textcomp} % provide euro and other symbols
  \usepackage{amssymb}
\else % if luatex or xetex
  \usepackage{unicode-math}
  \defaultfontfeatures{Scale=MatchLowercase}
  \defaultfontfeatures[\rmfamily]{Ligatures=TeX,Scale=1}
\fi
% Use upquote if available, for straight quotes in verbatim environments
\IfFileExists{upquote.sty}{\usepackage{upquote}}{}
\IfFileExists{microtype.sty}{% use microtype if available
  \usepackage[]{microtype}
  \UseMicrotypeSet[protrusion]{basicmath} % disable protrusion for tt fonts
}{}
\makeatletter
\@ifundefined{KOMAClassName}{% if non-KOMA class
  \IfFileExists{parskip.sty}{%
    \usepackage{parskip}
  }{% else
    \setlength{\parindent}{0pt}
    \setlength{\parskip}{6pt plus 2pt minus 1pt}}
}{% if KOMA class
  \KOMAoptions{parskip=half}}
\makeatother
\usepackage{xcolor}
\IfFileExists{xurl.sty}{\usepackage{xurl}}{} % add URL line breaks if available
\IfFileExists{bookmark.sty}{\usepackage{bookmark}}{\usepackage{hyperref}}
\hypersetup{
  pdftitle={Glossário de Termos Técnicos em Avaliação de Tecnologias em Saúde},
  pdfauthor={Carlos Lederman3 e Wanderley Marques Bernardo4},
  hidelinks,
  pdfcreator={LaTeX via pandoc}}
\urlstyle{same} % disable monospaced font for URLs
\usepackage{longtable,booktabs}
% Correct order of tables after \paragraph or \subparagraph
\usepackage{etoolbox}
\makeatletter
\patchcmd\longtable{\par}{\if@noskipsec\mbox{}\fi\par}{}{}
\makeatother
% Allow footnotes in longtable head/foot
\IfFileExists{footnotehyper.sty}{\usepackage{footnotehyper}}{\usepackage{footnote}}
\makesavenoteenv{longtable}
\usepackage{graphicx}
\makeatletter
\def\maxwidth{\ifdim\Gin@nat@width>\linewidth\linewidth\else\Gin@nat@width\fi}
\def\maxheight{\ifdim\Gin@nat@height>\textheight\textheight\else\Gin@nat@height\fi}
\makeatother
% Scale images if necessary, so that they will not overflow the page
% margins by default, and it is still possible to overwrite the defaults
% using explicit options in \includegraphics[width, height, ...]{}
\setkeys{Gin}{width=\maxwidth,height=\maxheight,keepaspectratio}
% Set default figure placement to htbp
\makeatletter
\def\fps@figure{htbp}
\makeatother
\setlength{\emergencystretch}{3em} % prevent overfull lines
\providecommand{\tightlist}{%
  \setlength{\itemsep}{0pt}\setlength{\parskip}{0pt}}
\setcounter{secnumdepth}{5}
\ifluatex
  \usepackage{selnolig}  % disable illegal ligatures
\fi
\usepackage[]{natbib}
\bibliographystyle{apalike}

\title{Glossário de Termos Técnicos em Avaliação de Tecnologias em Saúde}
\author{Carlos Lederman\textsuperscript{3} e Wanderley Marques Bernardo\textsuperscript{4}}
\date{Novembro/2020}

\begin{document}
\maketitle

\hypertarget{introduuxe7uxe3o}{%
\chapter*{Introdução}\label{introduuxe7uxe3o}}
\addcontentsline{toc}{chapter}{Introdução}

\begin{center}\includegraphics[width=height=4in,width=3in]{capa11} \end{center}

\begin{center}\rule{0.5\linewidth}{0.5pt}\end{center}

A Rede Internacional de Agências para a Avaliação de Tecnologias em Saúde (INAHTA\textsuperscript{1}, The International Network of Agencies for Health Technology Assessment) é uma organização sem fins lucrativos formada por instituições
da comunidade internacional, fundada em 1993. Todos os seus membros são organizações públicas, também sem fins lucrativos, que produzem ATS (Avaliações em Tecnologias em Saúde) sendo vinculadas a governos regionais ou nacionais.

A INAHTA mantém um Glossário\textsuperscript{2} de termos técnicos em ATS e os apresenta em quatro idiomas: Inglês, Espanhol, Francês e Alemão.

Além da tradução dos termos, o glossário incorpora uma breve explicação sobre cada termo, quando necessária, que acreditamos ser o grande valor deste glossário, pois apóia o leitor no entendimento do termo.

Dado o expressivo tamanho consolidado das populações cujo principal idioma é o Português, criamos este glossário e o apresentamos à INAHTA. As regras da INAHTA para inclusão de um novo idioma em seu Glossário exigem o suporte de uma instituição com determinadas características, que a ela se filie e que garanta a manutenção do novo glossário.

Depois de algumas tentativas em localizar esse tipo de instituição, consultamos a INAHTA com relação à sua eventual permissão no uso dos termos em Inglês de seu Glossário, para sua publicação em Português, respeitadas as seguintes condições: o Glossário seria publicado e ofertado ao público sem custo e sem a responsabilidade da INAHTA com relação ao conteúdo.

A permissão ocorreu em Setembro de 2020, após a qual publicamos este conteúdo em dois formatos: HTML, e PDF, dentro da plataforma Github, de acordo com as regras de licenciamento da Creative Commons.

A tradução foi feita por Carlos Lederman\textsuperscript{3} e Wanderley Marques Bernardo\textsuperscript{4}, e revista pelos seguintes alunos de Medicina do Centro Universitário Lusíada (UNILUS):

Guilherme Tavares e

Luca Schiliró Tristão

\begin{center}\rule{0.5\linewidth}{0.5pt}\end{center}

Os autores agradecem antecipadamente pelo envio de sugestões de aperfeiçoamentos e de novos termos (estes acompanhados de uma breve descrição de seu significado). O e-mail para recebimento dessas sugestões é \href{mailto:carloslederman01@gmail.com}{\nolinkurl{carloslederman01@gmail.com}}, colocando no campo de referência da mensagem as palavras ``Glossário ATS''.

\begin{center}\rule{0.5\linewidth}{0.5pt}\end{center}

O conteúdo deste não expressa opiniões das entidades onde os autores colaboraram ou colaboram. Todas as fontes utilizadas foram citadas. Todos os arquivos de dados tiveram seu acesso e uso autorizado para a finalidade deste livro.

\begin{center}\rule{0.5\linewidth}{0.5pt}\end{center}

Esta publicação foi realizada sem suporte econômico.

\begin{center}\rule{0.5\linewidth}{0.5pt}\end{center}

Ilustração da capa: realizada pelos autores utilizando o software R.

\begin{center}\rule{0.5\linewidth}{0.5pt}\end{center}

Glossário de Termos Técnicos em Avaliação de Tecnologias em Saúde

Novembro/2020

Carlos Lederman\textsuperscript{3}
BSc em Ciência da Computação/Unicamp, MBA em Gestão Estratégica e Econômica de Projetos/FGV, aluno de pós-graduação na Faculdade de Medicina da USP

Wanderley Marques Bernardo\textsuperscript{4}
PhD em Cirurgia Torácica/FMUSP, livre docente pelo Departamento de Gastroenterologia da Faculdade de Medicina da USP

\hypertarget{agradecimentos}{%
\chapter*{Agradecimentos}\label{agradecimentos}}
\addcontentsline{toc}{chapter}{Agradecimentos}

Os autores agradecem à INAHTA pela autorização no uso dos termos em Inglês de seu Glossário, o que possibilitou trazer este trabalho ao público, acreditando que será útil a todos os que atuam e/ou tem interesse na área da Saúde.

\hypertarget{objetivos}{%
\chapter*{Objetivos}\label{objetivos}}
\addcontentsline{toc}{chapter}{Objetivos}

Promover a sintonia dos países de língua portuguesa com os termos técnicos em Avaliação de Tecnologias em Saúde, utilizados nos países das línguas inglesa, alemã, espanhola e francesa.

Permitir que todos os países de língua portuguesa sugiram a inclusão de sinônimos, termos novos ou termos com significado distinto e que sejam característicos de seu país de aplicação.

\hypertarget{glossuxe1rio-portuguuxeas-ingluxeas}{%
\chapter*{Glossário Português / Inglês}\label{glossuxe1rio-portuguuxeas-ingluxeas}}
\addcontentsline{toc}{chapter}{Glossário Português / Inglês}

Para pesquisar um termo:

\begin{enumerate}
\def\labelenumi{\alph{enumi})}
\item
  clique na ``lupa'' (ferramenta de busca) e coloque a palavra de interesse, em Português ou em Inglês ou
\item
  no menu à esquerda, clique na letra que corresponde à primeira da palavra sendo pesquisada, aparecerão todas as existentes em ordem alfabética do Português
\end{enumerate}

\hypertarget{a}{%
\chapter*{A}\label{a}}
\addcontentsline{toc}{chapter}{A}

abordagem em capital humano

Um método para estimar os custos relacionados à perda de produção devida à doença. Nota: a estimativa é baseada no valor econômico para a sociedade da duração da vida saudável de uma pessoa de mesma idade e o valor econômico é medido pela estimativa da renda que pessoa poderia estar apta a obter.

human capital approach (HCA)

A method for estimating the costs related to loss of production due to illness. Note: The estimate is based on the economic value to society of the length of healthy life of a person the same age, and the economic value is measured by estimating the income that the person would have been able to earn.

\begin{center}\rule{0.5\linewidth}{0.5pt}\end{center}

ação de pesquisa

na pesquisa clínica: um processo no qual um investigador identifica, analisa e corrige um problema, e então o reacessa, se necessário, iniciando o processo novamente, até o problema ser resolvido completamente. Na pesquisa social: um método de pesquisa para a resolução de um problema num dado contexto, através de uma investigação democrática na qual os investigadores trabalham com os atores locais na busca e aplicação de soluções para um problema que seja de importância significativa para a comunidade.

action research

In clinical research: a process in which an investigator identifies, analyses and corrects a problem, then reassesses it and, if necessary, begins the process over again, until the problem is completely resolved. In social research: a research method for resolving a problem in a given context through a democratic investigation in which the investigators work with the local actors to seek and apply solutions to a problem that is of major importance to the community.

\begin{center}\rule{0.5\linewidth}{0.5pt}\end{center}

aceitabilidade do custo-efetividade

Curva que ilustra a probabilidade de uma dada opção ser eficiente com base no valor atribuido a um ano de vida adicional de qualidade ajustada (QALY).

cost-effectiveness acceptability (CEA)

A curve illustrating the probability that a given option is efficient on the basis of the value assigned to an additional quality-adjusted life year (QALY).

\begin{center}\rule{0.5\linewidth}{0.5pt}\end{center}

acesso

é o grau ao qual um paciente está apto a obter cuidados ou serviços, considerando-se restrições financeiras e organizacionais dos sistemas de saúde.

access

The degree to which a patient or group is able to obtain care or services, taking into account the health system's financial and organizational constraints.

\begin{center}\rule{0.5\linewidth}{0.5pt}\end{center}

acurácia

\begin{enumerate}
\def\labelenumi{\arabic{enumi})}
\tightlist
\item
  no contexto de um estudo, é a qualidade de uma medida (por exemplo, a estimativa média do efeito de um tratamento) que é correta ou que reflete a real efetividade do tratamento; não confundir com precisão: uma estimativa pode ser acurada, embora não precisa, se for baseada num método não tendencioso que mostre grandes variações aleatórias (severidade da doença, condições coexistentes). 2) no contexto de um teste diagnóstico, é a proporção na qual os resultados correspondem àqueles do teste de referência adotado (i.e.~a soma do verdadeiros positivos com a soma dos verdadeiros negativos, dividida pelo tamanho da amostra estudada).
\end{enumerate}

accuracy

\begin{enumerate}
\def\labelenumi{\arabic{enumi})}
\tightlist
\item
  In the context of a study, the quality of a measurement (e.g.~the mean estimate of a treatment effect) that is correct or that reflects the actual effectiveness of the treatment.Note: Not to be confused with preciseness: an estimate can be accurate, yet not be precise, if it is based on an unbiased method that shows great random variations (severity of the disease, coexisting conditions). 2) In the context of a diagnostic test, the proportion in which the results correspond to those of the chosen refererence test, i.e.~the sum of the true positives and the true negatives, divided by the size of the sample of the population studied.
\end{enumerate}

\begin{center}\rule{0.5\linewidth}{0.5pt}\end{center}

acurácia diagnóstica

vide acurácia

diagnostic accuracy

see accuracy

\begin{center}\rule{0.5\linewidth}{0.5pt}\end{center}

aderência

\begin{enumerate}
\def\labelenumi{\arabic{enumi})}
\tightlist
\item
  em pesquisa, é a medida da extensão à qual os elementos passam numa dada intervenção. Nota: pode ser a extensão à qual os elementos aderem com relação à terapia medicamentosa, passam pela prescrição médica ou procedimento cirúrgico, executam um determinado regime de exercícios, se abstém de fumar, etc. 2) para investigadores, o fato de seguir um protocolo de pesquisa 3) na definição de parâmetros clínicos, os elementos que se supõe seguir um tratamento ou aqueles registrados num programa, pelo fato de aderirem com os requisitos do tratamento ou do programa. Nota: como esta palavra tem esses 3 significados, para evitar qualquer confusão, é melhor referir-se a ``aderência a um protocolo de pesquisa'' quando se quiser indicar a aderência do investigador.
\end{enumerate}

compliance

\begin{enumerate}
\def\labelenumi{\arabic{enumi})}
\tightlist
\item
  In research, a measure of the extent to which the subjects undergo an assigned intervention. Note: It may be the extent to which the subjects comply with the drug therapy, undergo the prescribed medical or surgical procedure, do the prescribed exercise regimen, abstain from smoking, etc. 2) For investigators, the fact of following a research protocol. 3) In a clinical setting, for subjects who are supposed to undergo a treatment or who are registered in a program, the fact of complying with the requirements of that treatment or program. Note: Since this word has these three meanings, to avoid any confusion, it might be better to refer to ``compliance with the research protocol'' when what is meant is the investigator's compliance.
\end{enumerate}

\begin{center}\rule{0.5\linewidth}{0.5pt}\end{center}

advogado do consumidor

Pessoa ou organização que está ativamente envolvida com terceiros e que apresenta as perspectivas e preocupações de um grupo de pacientes.

consumer advocate

A person or organisation who/that is actively involved with others and presents the perspectives and concerns of a group of patients.

\begin{center}\rule{0.5\linewidth}{0.5pt}\end{center}

alfa

em teste de hipóteses, a probabilidade de erro do Tipo I, i.e.~a probabilidade de se concluir incorretamente que a hipótese nula é falsa. Nota: por exemplo, alfa poderia ser a probabilidade de se concluir que uma intervenção é efetiva quando ela não tem um efeito verdadeiro. Se os investigadores definirem alfa como 0.05, isso significa que eles aceitam 5\% de risco de um erro do Tipo I. O valor de alfa é comumente definido como 0.01, 0.05 ou 0.10.

alfa

In hypothesis testing, the probability of a Type I error, i.e.~the probability of concluding incorrectly that a null hypothesis is false. Note: For example, α could be the probability of concluding that an intervention is effective when it has no true effect. If the investigators set α at 0.05, this means that they accept that there is a 5\% risk of a Type I error. The α-level is commonly set at 0.01, 0.05 or 0.10.

\begin{center}\rule{0.5\linewidth}{0.5pt}\end{center}

alocação do tratamento

veja randomização

treatment allocation

See randomization

\begin{center}\rule{0.5\linewidth}{0.5pt}\end{center}

alocação não estritamente randômica

veja ocultação da alocação.

quasi random allocation

See concealment of allocation. Alternate spellings: quasi-random allocation, quasirandom allocation.

\begin{center}\rule{0.5\linewidth}{0.5pt}\end{center}

alvo móvel

Um assunto de avaliação (uma tecnologia, por exemplo) que é difícil de definir porque está mudando rapidamente.

moving target

A subject of evaluation (a technology, for example) that is difficult to define because it is changing rapidly.

\begin{center}\rule{0.5\linewidth}{0.5pt}\end{center}

amostragem randomizada

Um método de obtenção de uma mostra representativa de uma população.

random sampling

A method of obtaining a representative sample of a population.

\begin{center}\rule{0.5\linewidth}{0.5pt}\end{center}

análise Bayesiana

método estatístico que explicitamente inclui uma distribuição de probabilidades anterior, baseada numa opinião subjetiva ou numa evidência objetiva, tal como os resultados de uma pesquisa anterior.

Bayesian analysis

A statistical method that explicitly includes a prior probability distribution based on a subjective opinion or objective evidence, such as the results of previous research.

\begin{center}\rule{0.5\linewidth}{0.5pt}\end{center}

análise da intenção-para-tratar
Uma análise na qual os efeitos observados nos indivíduos de um ensaio são avaliados de acordo com a intervenção para a qual foram randomizados, se eles a receberam ou não e se eles concordaram ou não com o protocolo do estudo. Nota: este tipo de análise é favorecido, porque mantém a equivalência dos grupos randomizados. Além disso, na avaliação da efetividade prática de uma intervenção, ela espelha a não aderência e as alterações de tratamento que normalmente ocorrem na prática. Conceito relacionado: por protocolo.

intention-to-treat (ITT) analysis

An analysis in which the effects observed in the subjects of a trial are evaluated according to the intervention to which they were randomised, whether or not they received it, or whether or not they complied with the study protocol. Note: This type of analysis is favoured because it maintains the equivalence of the randomised groups. Also, in the evaluation of the practical effectiveness of an intervention, it mirrors the non-compliance and the treatment changes likely to occur in practice. (Related concept): per protocol

\begin{center}\rule{0.5\linewidth}{0.5pt}\end{center}

análise da minimização de custo

Avaliação econômica consistindo na comparação dos custos de várias opções presumidas para produzir desfechos equivalentes e para se determinar, destas, a opção de menor custo.

cost-minimization analysis (CMA)

An economic evaluation consisting of comparing the costs of various options presumed to produce equivalent outcomes and determining the least costly of those options.

\begin{center}\rule{0.5\linewidth}{0.5pt}\end{center}

análise de custo-beneficio

Avaliação econômica que consiste na comparação de várias opções, nas quais os custos e desfechos são quantificados em unidades monetárias comuns.

cost-benefit analysis

An economic evaluation consisting of comparing various options, in which costs and outcomes are quantified in common monetary units.

\begin{center}\rule{0.5\linewidth}{0.5pt}\end{center}

análise de custo-consequência

Avaliação econômica consistindo na comparação de várias opções, nas quais o custo dos componentes e os vários desfechos de uma doença ou intervenção são calculados e apresentados separadamente.

cost-consequent analysis (CCA)

An economic evaluation consisting of comparing various options, in which the cost components and the various outcomes of a disease or intervention are calculated and presented separately.

\begin{center}\rule{0.5\linewidth}{0.5pt}\end{center}

análise de custo-utilidade

Avaliação econômica consistindo na comparação de várias opções, nas quais os custos são medidos em unidades monetárias e os desefechos são medidos em unidades de medida de utilidade, usualmente em termos da utilidade para o paciente (usando anos de vida com a qualidade ajustada, por exemplo). Nota: esta é uma forma de análise de custo-efetividade na qual uma opção é ajustada com base na qualidade de vida.

cost-utility analysis (CUA)

An economic evaluation consisting of comparing various options, in which costs are measured in monetary units and outcomes are measured in utility units, usually in terms of utility to the patient (using quality-adjusted life years, for example). Note: This is a form of cost-effectiveness analysis in which the effectiveness of an option is adjusted on the basis of quality of life.

\begin{center}\rule{0.5\linewidth}{0.5pt}\end{center}

análise de impacto

Uma avaliação de efeitos de um relatório de avaliação de tecnologia em saúde ou de um serviço ou programa de saúde.

impact analysis

An evaluation of the effects of a health technology assessment report or of a health service or programme.

\begin{center}\rule{0.5\linewidth}{0.5pt}\end{center}

análise de regressão

Um método consistindo na seleção de um modelo matemático (regressão linear ou regressão logística, por exemplo) para descrever ou predizer o efeito de uma ou mais variáveis independentes X numa variável dependente Y. Nota: os termos são assim usados: regressão liner quando a variável Y é continua, regressão logística quando Y é uma variável dicotômica, regressão simples quando existe uma variável única X, e regressão múltipla quando existem várias variáveis X.

regression analysis

A method consisting of selecting a mathematical model (linear or logistic regression, for example) that explains the data to describe or predict the effect of one or more independent variables X on a dependent variable Y. Note: The terms are used as follows: linear regression when the variable Y is continuous, logistic regression when Y is a dichotomous variable, simple regression when there is a single variable X, and multiple regression when there are several variables X.

\begin{center}\rule{0.5\linewidth}{0.5pt}\end{center}

análise de sensibilidade

Um meio para se avaliar a robustez de um modelo matemático pelo teste da variação plausível das estimativas das varíáveis-chave independentes, para determinar se tais variações resultam em alterações significativas nos resultados dos modelos. Nota: a análise de sensibilidade também pode ser usada para outros estudos, tais como análises de ensaios clínicos, para determinar se a inclusão ou exclusão de certos dados altera os resultados, e em meta análise, para determinar se a inclusão ou exclusão de certos estudos alteram os resultados.

sensitivity analysis

A means for evaluating the robustness of a mathematical model by testing a plausible range of estimates of key independent variables to determine whether such variations result in meaningful changes in the model's results. Note: Sensitivity analysis can also be used for other study types, such as clinical trials analysis, to determine whether inclusion or exclusion of certain data changes the results, and meta-analysis, to determine whether inclusion or exclusion of certain studies changes the results.

\begin{center}\rule{0.5\linewidth}{0.5pt}\end{center}

análise do custo da doença

Avaliação do impacto financeiro geral na sociedade de uma dada doença ou problema de saúde. Nota: a soma, expressa em unidades monetárias, inclui custos de tratamento, mas não inclui os benefícios dos tratamentos ou dos desfechos.

cost-of-illness analysis

An evaluation of the overall economic impact on society of a given disease or health problem. Note: The sum, expressed in monetary units, includes treatment costs, but not treatment benefits or outcomes.

\begin{center}\rule{0.5\linewidth}{0.5pt}\end{center}

análise do impacto do orçamento

avaliação do impacto financeiro pela introdução de uma tecnologia ou de um serviço nos orçamentos de capital e operacional de um governo ou agência.

budget impact analysis

An evaluation of the financial impact of the introduction of a technology or service on the capital and operating budgets of a government or agency.

\begin{center}\rule{0.5\linewidth}{0.5pt}\end{center}

análise do subgrupo

Uma técnica consistindo na análise de dados de um subgrupo separadamente daqueles da população geral sendo estudada. Nota: esse tipo de análise pode ser planejado antecipadamente no protocolo de pesquisa, e os resultados que ele provê podem ser considerados apenas exploratórios. Qualquer efeito significativo observado num subgrupo deverá ser verificado em outros estudos. Conceito relacionado: teste múltiplo.

subgroup analysis

A technique consisting of analysing the data from a subgroup separately from those from the overall population studied. Note: This type of analysis should be planned beforehand in the research protocol, and the results it provides are to be considered exploratory only. Any significant effect observed in a subgroup must be verified in other studies. Related concept: multiple testing.

\begin{center}\rule{0.5\linewidth}{0.5pt}\end{center}

análise interna

Uma análise realizada durante um estudo para comparar os grupos experimentais com relação à efetividade ou segurança de uma intervenção, quando a situação requer. Nota: análises internas devem ser planejadas no protocolo (seu númerom o tempo de sua aplicação, os critérios subjacentes, etc.). Elas diferem da análise padrão que é relaizada ao final do período de seguimento provido pelo protocolo do estudo. O número de indivíduos num estudo deve refletir o número planejado de análises internas.

interim analysis

An analysis performed during a study to compare the experimental groups with regard to the effectiveness or safety of an intervention, when the situation requires.Note: Interim analyses must be planned in the protocol (their number, the time of their application, the criteria underlying them, etc.). They differ from the standard analysis, which is performed at the end of the follow-up period provided for in the study protocol. The number of subjects in a study must reflect the planned number of interim analyses.

\begin{center}\rule{0.5\linewidth}{0.5pt}\end{center}

ano de vida ajustado pela qualidade (QALY)

Uma unidade de desfecho de uma intervenção onde são ganhos (ou perdidos) anos de vida posteriores a essa intervenção, ajustados com base na qualidade de vida durante esses anos. Nota: esse parâmetro pode prover uma unidade comum para comparação de custo-utilidade através de diferentes intervenções e problemas de saúde. O DALY (ano de vida ajustado pela incapacidade) e o HYE (anos saudáveis equivalentes) são unidades análogas ao QALY.

quality-adjusted life year (QALY)

A unit of outcome of an intervention where gains (or losses) of years of life subsequent to this intervention are adjusted on the basis of the quality of life during those years. Note: This parameter can provide a common unit for comparing cost-utility across different interventions and health problems. Disability-adjusted life year (DALY) and healthy-years equivalent (HYE) are QALY-analogous units.

\begin{center}\rule{0.5\linewidth}{0.5pt}\end{center}

ano de vida com deficiência ajustada (DALY)

Uma unidade de estado de saúde onde a expectativa de vida, de acordo com a idade, é ajustada pela perda de saúde e anos de vida devidos à deficiência ou à lesão. Nota: esta medida é frequentemente utilizada para estabelecer a incidência global de uma moléstia (para uma grande população).

disability-adjusted life year (DALY)

A unit of health status where life expectancy according to age is adjusted by the loss of health and years of life due to disability from disease or injury. Note: This measurement is often used to assess the overall burden of disease (for a large population).

\begin{center}\rule{0.5\linewidth}{0.5pt}\end{center}

ano fiscal

Um período, geralmente de um ano, após o qual uma entidade prossegue com o encerramento das suas contas e a definição de suas demonstrações financeiras anuais.

FY

A period, generally one year, at the end of which an entity proceeds with the closing of its accounts and the establishment of its annual financial statements.

\begin{center}\rule{0.5\linewidth}{0.5pt}\end{center}

anualização

método para converter o custo do capital no custo equivalente anual, calculado com base na vida útil, custo inicial (custo de aquisição) e valor no final da vida útil de um equipamento, e as taxas de juros sobre o capital investido. Nota: este método também pode ser aplicado no cálculo da renda.

annualisation

A method for converting capital cost into an equivalent annual cost, calculated on the basis of the lifespan, initial cost (acquisition cost) and end-of-lifespan value of the equipment, and the interest rates on the invested capital.Note: This method can also be applied to income.

\begin{center}\rule{0.5\linewidth}{0.5pt}\end{center}

Aplicabilidade

a possibilidade de aplicar em outros contextos os resultados de uma observação ou estudo.

applicability

The possibility of applying in other settings the results of an observation or study.

\begin{center}\rule{0.5\linewidth}{0.5pt}\end{center}

aplicação com a aprovação de pre-marketing

Uma aplicação feita pelo fabricante ao FDA (Food and Drug Administration) para aprovação de fornecer ao mercado um produto médico, nos Estados Unidos. Nota: esta aplicação inclui a informação pertinente documentando a segurança e a efetividade do produto. Sinônimo: aplicação com a aprovação de pré-mercado

premarketing approval application (PMA)

An application made by the manufacturer to the Food and Drug Administration (FDA) for approval to market a medical product in the United States. Note: This application includes the pertinent information documenting the safety and effectiveness of the product. Alternative spellings: pre-market approval application, premarket approval application

\begin{center}\rule{0.5\linewidth}{0.5pt}\end{center}

aplicação de nova droga sob análise

Uma aplicação para a aprovação de uma droga não aprovada ou de uma droga previamente aprovada para outro uso, submetida pelo fabricante ao FDA (Unites States Food and Drug Administration - Administração de Alimentos e Drogas norte-americano) antes do teste dessa droga em humanos . Outra definição: um arquivo fornecido às autoridades competentes por um fabricante para obter a notificação de aprovação para uma nova droga e que contém as fases I, I e II dos ensaios clínicos. Nota: o termo ``submissão de nova droga'' é usado no Canadá e o termo ``aplicação de nova droga'' é usado nos Estados Unidos.

new drug application (NDA)

An application submitted by a manufacturer to the Food and Drug Administration (FDA) for approval to market a new drug (a new, non-biological molecular entity) for human use. Other definition: A file that is submitted to the competent authorities by a manufacturer to obtain a notice of compliance for a new drug and that contains phases I, II and III clinical trials. Note: The term new drug submission is used in Canada and the term new drug application is used in the United States.

\begin{center}\rule{0.5\linewidth}{0.5pt}\end{center}

Auditoria

ferramenta para a medida da aderência a padrões pré-estabelecidos, para a identificação de boas e más práticas. Nota: em alguns casos, a auditoria pode incluir a definição de padrões.

Audit

A tool for measuring compliance with pre-determined standards, to identify good and bad practices. Note: In some cases, the audit may include standard setting.

\begin{center}\rule{0.5\linewidth}{0.5pt}\end{center}

Autonomia

O princípio do direito à auto-determinação, i.e.~o direito dos pacientes de tomar suas próprias decisões à respeito de cuidados médicos. Nota: autonomia implica que as pessoas são livres e aptas a tomas decisões que as afetem com base em seu próprio sistema de valores e crenças

autonomy

The principle of the right of self-determination, i.e.~patients' right to make their own decisions regarding medical care. Note: Autonomy implies that people are free and able to make the decisions that affect them on the basis of their own system of values and beliefs.

\begin{center}\rule{0.5\linewidth}{0.5pt}\end{center}

Avaliação

processo científico usado para descrever e analisar as propriedades de uma tecnologia da saúde - sua segurança, eficácia, aplicabilidade e indicações de uso, custo e custo - efetividade, bem como as consequências sociais, econômicas e éticas

assessment

A scientific process used to describe and analyse the properties of a health technology---its safety, efficacy, feasibility and indications for use, cost and cost-effectiveness, as well as social, economic and ethical consequences.

\begin{center}\rule{0.5\linewidth}{0.5pt}\end{center}

avaliação crítica

Processo de se estabelecer e interpretar resultados de pesquisas científicas, sistematicamente analisando sua validade, sua significância estatística e sua relevância clínica

critical appraisal

The process of assessing and interpreting scientific research results by systematically analysing their validity, clinical and statistical significance, and clinical relevance.

\begin{center}\rule{0.5\linewidth}{0.5pt}\end{center}

avaliação da qualidade

Atividades de medida e monitoramento para se determinar o quão bem o cuidado com a saúde é entregue, com referência aos padrões aplicáveis ou limites aceitáveis de cuidados. (Conceitos relacionados: garantia da qualidade e cuidado com a qualidade)

quality assessment

The measurement and monitoring activities for determining how well health care is delivered with reference to the applicable standards or acceptable bounds of care.(Related concepts: quality assurance and quality of care.)

\begin{center}\rule{0.5\linewidth}{0.5pt}\end{center}

avaliação de campo

Um estudo projetado expressamente para coletar dados originais de uma tecnologia promissora, para verificar sua efetitividade prática, quando não há evidência suficiente para se tomar uma decisão baseada em informações.

field evaluation

A study designed expressly to collect original data on a promising technology to verify its practical effectiveness when there is insufficient evidence to make an informed decision.

\begin{center}\rule{0.5\linewidth}{0.5pt}\end{center}

avaliação de necessidades de saúde

Uma avaliação das necessidades de saúde da comunidade e o uso dessa informação para estabelecer prioridades e planejar a entrega e os serviços de cuidado com a saúde.

health needs assessment

An assessment of a community's health needs and the use of this information to set priorities and plan health care delivery and services.

\begin{center}\rule{0.5\linewidth}{0.5pt}\end{center}

avaliação de pesquisa

Um campo de questões que avalia a operação, processos, desfechos e utilidade de estruturas organizacionais (programas, agências, políticas, etc..).

evaluation research
A field of inquiry that assesses the operation, processes, outcomes and utility of organizational structures (programme, agency, policy, etc.).

\begin{center}\rule{0.5\linewidth}{0.5pt}\end{center}

avaliação de risco

A estimativa qualitativa ou quantitativa da verosimilhança de um evento adverso que possa resultar da exposição a um dado fator de risco ou da ausência de condições benéficas.

risk assessment

The qualitative or quantitative estimation of the likelihood of an adverse event that may result from exposure to a given risk factor or the absence of beneficial conditions.

\begin{center}\rule{0.5\linewidth}{0.5pt}\end{center}

avaliação de tecnologia

Veja avaliação de tecnologia em saúde

technology assessment

See health technology assessment.

\begin{center}\rule{0.5\linewidth}{0.5pt}\end{center}

Avaliação de tecnologia da saúde (ATS)

A avaliação sistemática das propriedades e efeitos da tecnologia na saúde, endereçando os efeitos diretos e intencionais dessa tecnologias, bem como as consequências indiretas e não intencionais, com o objetivo principal de informar a tomada de decisão relacionada a essa tecnologia. Nota: ATS é conduzida por grupos interdisciplinares que usam gabaritos analíticos explícitos numa variedade de métodos.

healt technology assessment (HTA)

The systematic evaluation of the properties and effects of a health technology, addressing the direct and intended effects of this technology, as well as its indirect and unintended consequences, and aimed mainly at informing decision making regarding health technologies. Note: HTA is conducted by interdisciplinary groups that use explicit analytical frameworks drawing on a variety of methods.

\begin{center}\rule{0.5\linewidth}{0.5pt}\end{center}

Avaliação de Tecnologia em Saúde Internacional

HTAi (Avaliação de Tecnologia em Saúde Internacional) - uma sociedade internacional focada na avaliação de tecnologia em saúde que congrega instituições acadêmicas, instituições de cuidados com a saúde, indústria, negócios, o setor de voluntariado e o setor público.

HTAi

Health Technology Assessment International - an international society focussed on health technology assessment that brings together academic institutions, health care institutions, industry, business, the voluntary sector and the public sector.

\begin{center}\rule{0.5\linewidth}{0.5pt}\end{center}

avaliação econômica

A análise comparativa dos custos e consequências de duas ou mais opções possíveis. Nota: dependendo das consequências serem expressas como variáveis monetárias, físicas ou qualitativas, a análise poderá ser, respectivamente, de custo-benefício, de custo-efetividade ou de custo-utilidade.

economic evaluation

The comparative analysis of the costs and consequences of two or more possible options. Note: Depending on whether the consequences are expressed as monetary, physical or qualitative variables, the analysis may be a cost-benefit, cost-effectiveness or cost-utility analysis.

\hypertarget{b}{%
\chapter*{B}\label{b}}
\addcontentsline{toc}{chapter}{B}

banco de dados bibliográfico

fonte computadorizada e indexada de citações descrevendo documentos tais como artigos de jornais e relatórios científicos. Nota: as citações tipicamente incluem o(s) nome(s) do autor(es), o titulo, a fonte, um resumo e, quando aplicavel, informações relacionadas (incluindo o texto completo em alguns casos). Exemplos: MEDLINE e EMBASE.

bibliographic database

An indexed computerized or printed source of citations, describing documents such as journal articles and scientific reports. Note: The citations typically include the name(s) of the author(s), the title, the source, an abstract and, where applicable, related information (including the full text in some cases). Example: MEDLINE and EMBASE.

\begin{center}\rule{0.5\linewidth}{0.5pt}\end{center}

base de dados de ATS

Uma base de dados contendo relatórios sobre avaliações de tecnologia na saúde em desenvolvimento e completadas. Nota: essa base de dados é gerenciada pelo CRD (Centre for Reviews and Dissemination - centro de revisões e disseminação) em colaboração com o secretariado da INAHTA. Está disponível na Livraria Cochrane no site do CRD (\url{http://www.york.ac.uk/inst/crd/}).

HTA database

A database containing the reports on ongoing and completed health technology assessments. Note: This database is managed by the Centre for Reviews and Dissemination (CRD) in collaboration with the INAHTA Secretariat. It is available in the Cochrane Library and on the CRD site (\url{http://www.york.ac.uk/inst/crd/}).

\begin{center}\rule{0.5\linewidth}{0.5pt}\end{center}

base de dados factual

Uma fonte indexada e computadorizada ou impressa que provê acesso direto a informações tais como guias para diagnóstico e tratamento, estatísticas e diretórios de instituições. Outra definição: uma coleção extensa e considerada completa, de fatos e dados sobre materiais num campo especializado, disponíveis para análise e aplicação. Nota: a coleção pode ser computadorizada através de vários métodos modernos de pesquisa bibliográfica. Este conceito difere do banco de dados bibliográfico, que se limita a coleções de citações.

factual database

An indexed computerized or printed source that provides direct access to information such as guidelines for diagnosis and treatment, statistics, and directories of institutions. Other definition: An extensive collection, deemed complete, of facts and data from materials in a specialized field, made available for analysis and application. Note: The collection may be computerized through various modern literature search methods. This concept differs from that of a bibliographic database, which is limited to collections of citations.

\begin{center}\rule{0.5\linewidth}{0.5pt}\end{center}

base do estudo

Uma população estudada numa particular localização por um período de tempo especificado.

study base

A population studied in a particular location over a specified period of time.

\begin{center}\rule{0.5\linewidth}{0.5pt}\end{center}

benefício líquido

O benefício de uma opção em unidades monetárias, menos o seu custo total em unidades monetárias. Nota: o benefício líquido é um critério básico de decisão na análise de custo-benefício (ACB).

net benefit

The benefit of an option in monetary units, minus its total cost in monetary units. Note: Net benefit is a basic decision criterion in cost-benefit analysis (CBA).

\begin{center}\rule{0.5\linewidth}{0.5pt}\end{center}

benefício marginal

O benefício adicional produzido pelo uso de recursos adicionais, expressos, por exemplo, em unidades de desfecho em saúde.

marginal benefit

The additional benefit produced by the use of additional resources, expressed, for example, in units of health outcome.

\begin{center}\rule{0.5\linewidth}{0.5pt}\end{center}

beta

em teste de hipóteses, a probabilidade de erro do Tipo II, i.e.~a probabilidade de se concluir incorretamente que a hipótese nula é verdadeira. Nota: por exemplo, beta poderia ser a probabilidade de se concluir que uma intervenção não é efetiva quando ela tem um efeito verdadeiro. (1-beta) é o poder estatístico de um teste permitindo a rejeição de uma hipótese nula que seja verdadeiramente falsa (i.e.~detectando o efeito de uma intervenção que relamente existe)

beta

In hypothesis testing, the probability of a Type II error, i.e.~the probability of concluding incorrectly that a null hypothesis is true. Note: For example, β could be the probability of concluding that an intervention is not effective when it has a true effect. (1-β) is the statistical power of a test allowing for rejection of a null hypothesis that is truly false (e.g.~detecting the effect of an intervention that truly exists).

\begin{center}\rule{0.5\linewidth}{0.5pt}\end{center}

Biblioteca Cochrane

conjunto de bases de dados atualizadas trimestralmente, projetadas para prover informação científica e evidências para suportar a tomada de decisões no campo da saúde, publicadas em CD-ROM e na Internet. Nota: as bases de dados da Biblioteca Cochrane são: o banco de dados Cochrane de revisões sistemáticas; a central de registro de ensaios controlados da Cochrane; o banco de dados de resumos de revisões de efeitos da Cochrane; o banco de dados de Avaliação de Tecnologias em Saúde e o banco de dados NHS de avaliação econômica.

Cochrane library

A collection of databases updated quarterly, designed to provide scientific information and evidence to support decision making in the health field, published on CD-ROM and on the Internet. Note: The Cochrane Library databases are: the Cochrane Database of Systematic Reviews, the Cochrane Central Register of Controlled Trials, the Database of Abstracts of Reviews of Effects, the Cochrane Methodology Register, the Health Technology Assessment database and the NHS Economic Evaluation database.

\begin{center}\rule{0.5\linewidth}{0.5pt}\end{center}

Bioética

campo de investigação compreendendo a análise de questões éticas e dilemas morais pertencentes ao estudo, classificação e manipulação de vida biológica, especialmente na saúde humana.

bioethics

A field of inquiry comprising the analysis of ethical issues and moral dilemmas pertaining to the study, classification, and manipulation of biological life, especially human health.

\hypertarget{c}{%
\chapter*{C}\label{c}}
\addcontentsline{toc}{chapter}{C}

cadeia de Monte Carlo / Markov (MCMC)

Um método de análise probabilística combinado com o modelo de Markov, que captura as probabilidades transacionais entre vários estados do estudo da interação entre as distribuições de probabilidade relacionadas a cada variável. Nota: este tipo de simulação proporciona uma estimativa do grau de incerteza de um modelo. É largamente utilizado na análise Bayesiana, na construção ou na atualização de modelos de distribuição conjunta de variáveis.

Markov chain Monte Carlo (MCMC)

A probabilistic analysis method combined with the Markov model, which captures the transitional probabilities between various states to study the interaction between the probability distributions attached to each variable. Note: This type of simulation provides an estimate of the degree of uncertainty in a model. It is widely used in Bayesian analysis to build or update models of the joint distribution of variables.

\begin{center}\rule{0.5\linewidth}{0.5pt}\end{center}

Caso

em epidemiologia, a pessoa com a doença ou característica de interesse para um ensaio clínico.

Case

In epidemiology, a person with the disease or characteristic of interest in a clinical trial.

\begin{center}\rule{0.5\linewidth}{0.5pt}\end{center}

causalidade

causas relacionadas e os efeitos por elas produzido. Nota: os critériosde Bradford Hill para associação causal são: consistência; força; especificidade; relação dose/resposta; relação temporal (a exposição sempre antecede o desfecho; este é o único critério essencial); plausibilidade biológica; coerência e experimento.

causality

The relating of causes to the effects they produce. Note: The Bradford Hill criteria for causal association are: consistency; strength; specificity; dose--response relationship; temporal relationship (exposure always precedes the outcome; it is the only essential criterion); biological plausibility; coherence; and experiment.

\begin{center}\rule{0.5\linewidth}{0.5pt}\end{center}

CDMR - base de dados de revisões metodológicas

veja a metodologia de registro Cochrane.

Cochrane database of methodology reviews (CDMR)

See Cochrane Methodology Register.

\begin{center}\rule{0.5\linewidth}{0.5pt}\end{center}

CDSR - base de dados de revisões sistemáticas da Cochrane

base de dados composta de textos completos de todas as revisões sistemáticas e de todos os projetos de revisões sistemáticas (protocolos) da ``Cochrane Collaboration'' (cooperação Cochrane).

Cochrane database od systematic reviews (CDSR)

A database composed of the full text of all Cochrane Collaboration systematic reviews and systematic review projects (protocols).

\begin{center}\rule{0.5\linewidth}{0.5pt}\end{center}

cegamento

num ensaio clínico, a ocultação dos elementos e/ou investigadores da identidade das pessoas recebendo as intervenções em estudo e as intervenções de controle (ou cuidados-padrão). Nota: a intenção deste método é eliminar a possibilidade de que o conhecimento da intervenção sendo ministrada aos pacientes possa afetar o desfecho ouque o comportamento do investigador possa afetar o desfecho. Nem sempre o cegamento é prático (Por exemplo, ao comparar uma cirurgia com um tratamento medicamentoso), mas deveria ser usado sempre que possível quando compatível com a otimização do cuidado ao paciente.

blinding

In a clinical trial, the concealing from the subjects and/or investigators of the identity of the persons receiving the interventions being studied and the control interventions (or standard care). Note: This method is intended to eliminate the possibility that knowledge of which intervention is being received will affect patient outcomes or investigator behaviours that may affect outcomes. Blinding is not always practical (e.g.~when comparing surgery to drug treatment), but it should be used whenever it is possible and compatible with optimal patient care.

\begin{center}\rule{0.5\linewidth}{0.5pt}\end{center}

cegamento triplo

Um ensaio no qual os indivíduos, os investigadores (assessores de desfecho) e os estatísticos realizando a análise dos desfechos não sabem quem está em cada grupo sendo comparado. (Conceitos relacionados: cegamento, cegamento único e cegamento duplo). Nota: os termos segamento e mascaramento são sinônimos e o termo mascaramento duplo é preferido sobre o cegamento duplo em ensaios oftalmológicos.

triple-blind

A trial in which the subjects, the investigators (outcome assessors) and the statistician doing the outcomes analysis do not know who is in which comparison group. (Related concepts: blinding, single-blind and double-blind). Note: The terms blinding and masking are synonymous and the term double-masked is preferred over double-blind for ophthalmology trials.

\begin{center}\rule{0.5\linewidth}{0.5pt}\end{center}

cegamento único

Um ensaio onde apenas os investigadores ou os indivíduos sabiam qual intervenção ou tratamento seria ministrado. (Conceitos relacionados: cegamento, cegamento duplo e cegamento triplo). Nota: os termos cegamento e mascaramento são sinônimos e o termo mascaramento duplo é preferido ao invés de cegamento duplo em função de ensaios oftalmológicos.

single-blind

A trial in which only the investigators or subjects know which intervention or treatment is being administered. (Related concepts: blinding, double-blind and triple-blind). Note: The terms blinding and masking are synonymous and the term double-masked is preferred over double-blind for ophthalmology trials.

\begin{center}\rule{0.5\linewidth}{0.5pt}\end{center}

CENTRAL -central de registro de ensaios controlados da Cochrane

base de dados de ensaios clínicos no campo da saúde, compilado das listas de grupos especializados do ``Cochrane Collaboration'' (colaboração Cochrane) e de outras organizações e de pesquisas nas bases de dado MEDLINE, EMBASE e outras.

Cochrane central register of controlled trials (CENTRAL)

A database of clinical trials in the health field, compiled from the lists of the specialized groups in the Cochrane Collaboration and other organisations and from searches of MEDLINE, EMBASE and other databases.

\begin{center}\rule{0.5\linewidth}{0.5pt}\end{center}

chance

A proporção (o quociente) da probabilidade de que um evento ocorrerá com relação à probabilidade de que esse evento não ocorrerá.

odds

The ratio (quotient) of the probability that an event will occur to the probability that this event will not occur.

\begin{center}\rule{0.5\linewidth}{0.5pt}\end{center}

CINAHL (índice cumulativo de literatura sobre enfermagem e assuntos aliados à saúde)

base de dados eletrônica descrevendo literatura sobre enfermagem e ciências paramédicas publicadas desde 1982.

CINAHL (cumulative index to nursing and allied health literature)

An electronic database describing the nursing and paramedical sciences literature published since 1982.

\begin{center}\rule{0.5\linewidth}{0.5pt}\end{center}

Citação

registro numa base de dados bibliográfica descrevendo documentos tais como artigos, livros e relatórios. Nota: a descrição inclui o(s) nome(s) do(s) autor(es), o título, a fonte, o ano de publicação, e, às vezes, as palavras-chave e resumo.

Citatiton

A record in a bibliographic database describing documents such as articles, books and reports.Note: The description includes the name(s) of the author(s), the title, the source, the year of publication and, sometimes, the keywords and abstract.

\begin{center}\rule{0.5\linewidth}{0.5pt}\end{center}

classificação indevida

A classificação de um indivíduo ou uma característica como sendo de uma categoria diversa daquela para a qual ela foi atribuida, levando a um viés de medida.

misclassification

The classification of a person or a characteristic as being in a category other than that to which it should have been assigned, leading to a measurement bias.

\begin{center}\rule{0.5\linewidth}{0.5pt}\end{center}

CMR - metodologia de registro da Cochrane

banco de dados de estudos metodológicos, publicado pela Biblioteca Cochrane

Cochrane metodology register (CMR)

A database of methodological studies, published by the Cochrane Library.

\begin{center}\rule{0.5\linewidth}{0.5pt}\end{center}

cobertura condicional

financiamento ou restituição (de valores monetários) em tecnologia da saúde ou num serviço condicionado a um critério específico. Nota 1: o financiamento é frequentemente proporcionado por uma agência de cuidado da saúde ou por um pagador, tal como: o governo, organização de manutenção da saúde ou organização de seguro-saúde. Nota 2: os critérios podem ser: - impostos pelo pagador (por exemplo, uma restrição na população elegível para a tecnologia, um acordo de preço/volume ou um finaciamento atrelado ao atingimento de um desfecho clínico) ou - acordado entre o pagador e o provedor da tecnologia como medida intermediária durante o desenvolvimento da evidência (por exemplo, entrada gerenciada ou cobertura com evidência de arranjos de desenvolvimento)

conditional coverage

The funding or reimbursement of a health technology or service that is conditional on specific criteria. Note 1: The funding is often provided by a health care agency or payer, such as: the government, health maintenance organization or health insurance organization. Note 2: The criteria could be:
- imposed by the payer (for example, a restriction on the population eligible for the technology, a price-volume agreement or funding tied to the achievement of clinical outcomes), or
- agreed between the payer and the technology provider as an interim measure during evidence development (for example, managed entry or coverage with evidence development arrangements).

\begin{center}\rule{0.5\linewidth}{0.5pt}\end{center}

coeficiente de correlação

Um número entre -1 e +1 que expressa a força da associação linear entre duas variáveis numéricas. Nota: Numa amostra, a estimativa é representada por r. O coeficiente de correlação igual a zero indica que não há relação linear etre as duas variáveis. Um coeficiente de correlação igual a +1 indica que existe uma relação linear perfeita, e um coeficiente de correlação igual a -1 idica que existe uma relação linear perfeitamente negativa. Nota: Não se pode concluir pela correlação de que existe uma relação de causa e efeito.

correlation coeficiente

A number between -1 and +1 that expresses the strength of the linear association between two numerical variables. Note: In a sample, the estimate is noted as r. A correlation coefficient of 0 indicates that there is no linear relationship between the two variables. A correlation coefficient of +1 indicates that there is a perfect positive linear relationship, and a correlation coefficient of -1 indicates that there is a perfect negative linear relationship. Note: It cannot be concluded from the correlation that there is a cause-and-effect relationship.

\begin{center}\rule{0.5\linewidth}{0.5pt}\end{center}

co-intervenção

num ensaio clínico controlado e randomizado, a aplicação de diagnóstico adicional ou procedimentos terapêuticos aos elementos de um grupo experimental ou ao grupo de controle, ou a ambos.

cointervention

In a randomized controlled trial, the application of additional diagnostic or therapeutic procedures to the members of the experimental group or the control group, or both.

\begin{center}\rule{0.5\linewidth}{0.5pt}\end{center}

co-morbidade

é a presença de doenças coexistentes com a doença sendo estudada

comorbidity

The presence of diseases co-existing with the disease being studied.

\begin{center}\rule{0.5\linewidth}{0.5pt}\end{center}

comparação

processo de gerenciamento da qualidade que uma organização utiliza para estabelecer metas e medir sua performance com relação a organizações reconhecidas por sua excelência.

benchmarking

A quality management process that an organization uses to set goals and measure its performance against that of organizations recognized for their excellence.

\begin{center}\rule{0.5\linewidth}{0.5pt}\end{center}

comparador

em pesquisa clínica, é uma tecnologia ou intervenção que serve com referência em estudos comparativos. Na economia em saúde, é uma opção com a qual a tecnologia ou intervenção é comparada.

comparator

In clinical research, a technology or intervention that serves as a reference in comparative studies. In health economics, an option with which the technology or intervention evaluated is compared.

\begin{center}\rule{0.5\linewidth}{0.5pt}\end{center}

confiabilidade

\begin{enumerate}
\def\labelenumi{\arabic{enumi})}
\tightlist
\item
  A habilidade de um exame ou teste, feito repetidamente na mesma população, em produzir o mesmo resultado (confiabilidade no teste-reteste). 2) A habilidade de um único exame ou teste em sempre fazer a mesma distinção entre os indivíduos de uma dada população.
\end{enumerate}

reliability

\begin{enumerate}
\def\labelenumi{\arabic{enumi})}
\tightlist
\item
  The ability of an examination or test, done repeatedly and in the same population, to produce the same result (test-retest reliability). 2) The ability of a single examination or test to always make the same distinction among individuals in a given population.
\end{enumerate}

\begin{center}\rule{0.5\linewidth}{0.5pt}\end{center}

confiabilidade entre avaliadores distintos (confiabilidade inter-avaliadores)

O grau de estabilidade de uma medida quando é repetida sob condições idênticas por avaliadores distintos. Nota: confiabilidade entre avaliadores distintos torna possível determinar o grau ao qual os resultados obtidos podem ser replicados. A falta dessa confiabilidade pode acontecer por divergências entre avaliadores ou pela instabilidade do atributo sendo medido. Conceito relacionado: confiabilidade no mesmo avaliador (confiabilidade intra-avaliador).

inter-rater reliability

The degree of stability of a measurement when it is repeated under identical conditions by different raters.Note: Inter-rater reliability makes it possible to determine the degree to which the results obtained can be replicated. Lack of this reliability may arise from divergences between raters or instability of the attribute measured. Related concept: intra-rater reliability

\begin{center}\rule{0.5\linewidth}{0.5pt}\end{center}

confiabilidade entre o mesmo avaliador

O grau de estabilidade observado quando uma medida é repetida sob condições idênticas pelo mesmo avaliador. Nota: a confiabillidade entre o mesmo avaliador torna possível determinar o grau ao qual os resultados obtidos por um procedimento de medição pode ser replicado. A falta de confiabilidade intra-avaliador pode conduzir a divergências entre instrumentos de medida ou a instabilidade do atributo sendo medido. Conceito relacionado: confiabilidade enter-avaliadores.

intra-rater reliability

The degree of stability observed when a measurement is repeated under identical conditions by the same rater. Note: Intra-rater reliability makes it possible to determine the degree to which the results obtained by a measurement procedure can be replicated. Lack of intra-rater reliability may arise from divergences between measurement instruments or instability of the attribute measured. Related concept: inter-rater reliability

\begin{center}\rule{0.5\linewidth}{0.5pt}\end{center}

confidencialidade

promessa de um profissional ao seu cliente em não revelar informações sem o consentimento do cliente. Outra definição: um requisito ético e regra baseada em privacidade que implica na obrigação dos profissionais da saúde em manter para si o que aprenderam, viram e ouviram na prática de sua profissão. Nota: a confidencialidade é vital para manter o sucesso de uma relação terapêutica. Pela lei, a confidencialidade é considerada explícita ou implícita no contrato entre o paciente e o profissional do qual se requer manter a confidencialidade.

confidentiality

A professional's promise to his or her client not to reveal information without the client's consent. Other definition: An ethical requirement and rule based on privacy that implies an obligation for health professionals to keep to themselves what they have learned, seen and heard in the practising of their profession. Note: Confidentiality is vital for maintaining a successful therapeutic relationship. In law, confidentiality is considered an explicit or implicit contract between the patient and the professional that requires the latter to maintain confidentiality.

\begin{center}\rule{0.5\linewidth}{0.5pt}\end{center}

conflito de escolha de tempo

Um método para determinar a preferência entre dois estados de saúde em diferentes durações de tempo, para estimar quantos anos de vida uma pessoa está preparada para sacrificar para melhorar seu estado de saúde. Nota: para estados crônicos, as opções são a referência do estado de saúde no tempo t seguida pela morte, ou a saúde perfeita por um curto período de tempo seguida pela morte. Para estados temporários, as opções dão a referência do estado de saúde no tempo t seguida por um desfecho específico (usualmente a morte), ou um pior estado de saúde por um curto período de tempo seguida pelo mesmo desfecho.

time trade-off (TTO)

A method for determining preference between two health states for different lengths of time, to estimate how many years of life a person is prepared to sacrifice to improve his/her health status. Note: For chronic states, the options are the reference health state for time t followed by death, or perfect health for a shorter time followed by death. For temporary states, the options are the reference health state for time t followed by an explicitly specified outcome (usually health), or a worse health state for a shorter time followed by the same outcome.

\begin{center}\rule{0.5\linewidth}{0.5pt}\end{center}

conflito de interesse

situação na qual o interesse privado de uma pessoa contribui para o exercício de sua influência sobre a qualidade ou sobre os resultados da precisão na avaliação dos dados.

conflict of interest

A situation in which the private interests of a person contributing to an assessment influence the quality or the results of the assessment or the accuracy of the data.

\begin{center}\rule{0.5\linewidth}{0.5pt}\end{center}

confundidor

fator que distorce o efeito aparente de uma intervenção num desfecho de interesse, ou que distorce a aparente associação entre a exposição e o desfecho. Exemplo: se as pessoas do grupo experimental de um ensaio controlado são mais novas do que as do grupo controle, será difícil decidir se um menor risco de morte num grupo é devido á intervenção ou á diferença de idades. A idade é então denominada confundidor. Nota 1: A identificação de confundidores requer conhecimento especializado ou substantivo sobre a rede causal que formam a intervenção (ou exposição) e o desfecho (por exemplo, conhecimento patopsicológico e conhecimento clínico). Tentativas de se selecionar confundidores somente com base nas associações estatísticas observadas podem conduzor a viéses. Nota 2: Variável confundidora e fator confundidor são sinônimos. Nota 3: termos relacionados incluem viés de confusão.

confounder

A factor that distorts the apparent effect of an intervention on the outcome of interest, or that distorts the apparent association between the exposure and the outcome. Example: If people in the experimental group of a controlled trial are younger than those in the control group, it will be difficult to decide whether a lower risk of death in one group is due to the intervention or to the difference in ages. Age is then said to be a confounder. Note 1: The identification of confounders requires expert or substantive knowledge about the causal network of which the intervention (or exposure) and the outcome are part (e.g.~pathophysiological and clinical knowledge). Attempts to select confounders solely on the basis of observed statistical associations may lead to bias. Note 2: Synonyms include confounding variable and confounding factor. Note 3: Related terms include confounding bias.

\begin{center}\rule{0.5\linewidth}{0.5pt}\end{center}

consentimento informado

A obrigação ética e legal em não realizar algum procedimento médico significativo até que o competente paciente tenha sido informado da natureza e riscos do procedimento e as alternativas ao mesmo, bem como o prognóstico se o procedimento não for realizado. Nota: se o paciente concordar com o procedimento, ele ou ela devem fazê-lo com liberdade e voluntariamente.

informed consent

The legal and ethical obligation not to perform any significant medical procedure until a competent patient has been informed of the nature and risks of the procedure and the alternatives to it, as well as of the prognosis if the procedure is not done. Note: If the patient agrees to the procedure, he or she must do so freely and voluntarily.

\begin{center}\rule{0.5\linewidth}{0.5pt}\end{center}

consequência

Desfecho positivo ou negativo de uma doença ou intervenção, expressa qualitativa ou quantitativamente. Nota: por exemplo, acidente vascular cerebral, morte, efeitos colaterais ou doença evitada

consensus report

The positive or negative outcome of a disease or intervention, expressed qualitatively or quantitatively. Note: For example, stroke, death, side effects or avoided disease.

\begin{center}\rule{0.5\linewidth}{0.5pt}\end{center}

construtivismo social

Uma doutrina filosófica em ciência e conhecimento que postula não haver objetivo ou realidade absolutamente social, mas apenas várias formas de interpretação dessa relaidade. Sinônimo: construcionismo social.

social constructivism

A philosophical doctrine on science and knowledge that postulates that there is no objective or absolute social reality, but only various ways of interpreting that reality. Syn.: social constructionism

\begin{center}\rule{0.5\linewidth}{0.5pt}\end{center}

consumidor

Pessoa que usa, é afetada, intitulada ou compelida a usar um serviço de saúde. Nota: no campo da saúde, o termo consumidor é usado principalmente nos Estados Unidos, onde não existe um sistema de saúde universal. Em outros países, o termo paciente é frequentemente utilizado, mas como essa palavra deveria ser aplicada somente a pessoas que consultam um médico, termos como usuário, recipiente e cliente são utilizados no sistema de saúde canadense.

consumer

A person who uses, is affected by, is entitled to or is compelled to use a health service. Note: In the health care field, the term consumer is used mainly in the United States, where there is no universal health insurance system. Elsewhere, patient is often used, but since this word should be applied only to persons who consult a physician, terms such as user, recipient and client are used in the Canadian health system.

\begin{center}\rule{0.5\linewidth}{0.5pt}\end{center}

contaminação

Em ensaios clínicos, a situação que ocorre quando membros do grupo de controle recebem a intervenção experimental ou quando os elementos do grupo de intervenção não a recebem.

contamination

In clinical trials, the situation that occurs when the members of the control group receive the experimental intervention or when the subjects in the experimental group do not receive the intervention.

\begin{center}\rule{0.5\linewidth}{0.5pt}\end{center}

conteúdo atual

Base de dados eletrônica que provê acesso a tabelas de conteúdo e dados bibliográficos de publicações recentes dos principais jornais científicos e dos principais livros em ciências, ciências sociais, artes e humanidades.

current contents

An electronic database that provides access to the tables of contents and bibliographic data from recent issues of the major science journals and the leading books in the sciences, social sciences, arts and humanities.

\begin{center}\rule{0.5\linewidth}{0.5pt}\end{center}

contexto

Condições e circunstâncias nas quais a intervenção é aplicada. Nota: por exemplo, a preparação do ambiente (no hospital, em casa ou numa ambulância aérea), o tempo (dia útil, feriado, período noturno), o tipo de prática (cuidados primários, secundários ou terciários, prática de seguro, caridade).

context

The conditions and circumstances in which an intervention is applied. Note: For example, the setting (in hospital, at home, in an air ambulance), the time (working day, holiday, night time), the type of practice, (primary, secondary or tertiary care, private practice, insurance practice, charity).

\begin{center}\rule{0.5\linewidth}{0.5pt}\end{center}

contra-indicação

Qualquer circunstância indicando que a intervenção poderia ser inadequada

contraindication

Any circumstance indicating that an intervention would be inappropriate.

\begin{center}\rule{0.5\linewidth}{0.5pt}\end{center}

controle

\begin{enumerate}
\def\labelenumi{\arabic{enumi}.}
\tightlist
\item
  Em ensaios clínicos, ao se comparar duas ou mais intervenções um elemento do grupo de controle ao qual é oferecido um placebo, não recebe tratamento ou recebe o cuidado padrão. 2. Em estudos de caso-controle, uma pessoa do grupo de comparação que não tem a doença ou o estado clínico de interesse. 3. Em controle estatístico, um procedimento que ajusta os resultados para levar em consideração fatores estranhos que poderiam distorcê-los. 4. Controle de infecção. Em saúde pública, um programa que deseja reduzir a extensão ou as consequências de uma doença infecciosa.
\end{enumerate}

control

\begin{enumerate}
\def\labelenumi{\arabic{enumi}.}
\tightlist
\item
  In clinical trials comparing two or more interventions, a subject in the comparison group who is offered a placebo, receives no treatment or receives standard care. 2. In case-control studies, a person in the comparison group who is without the disease or clinical state of interest. 3. statistical control. In statistics control, a procedure to adjust the results to take into account extraneous factors that could distort them. 4. infection control
  In public health, a programme aimed at reducing the extent or the consequences of an infectious disease.
\end{enumerate}

\begin{center}\rule{0.5\linewidth}{0.5pt}\end{center}

controle concorrente

um grupo de controle que é observado pelos investigadores ao mesmo tempo do grupo de intervenção, onde os dois grupos não foram definidos por randomização. Nota: podem ocorrer diferenças nas características dos elementos dos dois grupos.

concurrent control

A control group that is observed by the investigators at the same time as the treatment group, where the two groups were not established using randomisation. Note: Differences in the characteristics of the subjects in the two groups may result.

\begin{center}\rule{0.5\linewidth}{0.5pt}\end{center}

controle histórico

Um grupo de controle composto de pessoas que foram observadas numa ocasião anterior à presente, envolvendo a observação do grupo de interesse. Nota: a validade das comparações com esse tipo de grupo de controle é suspeita, porque o grupo histórico tende a diferir do grupo atual em sua composição, diagnóstico, severidade da doença, medidas dos desfechos, ou outros meios importantes que poderiam influenciar as descobertas da pesquisa. Pode ser factível usar controles históricos em circunstâncias especiais onde os desfechos de um tratamento padrão (ou a ausência de tratamento) forem bem conhecidos e variem pouco para uma dada população.

historical control

A control group made up of persons who have been observed on an occasion previous to the present one involving observation of the group of interest. Note: The validity of comparisons with this type of control group is suspect because the historical group is likely to differ from the current group in its composition, diagnosis, disease severity, measurement of outcomes, or other important ways that could influence the research findings. It may be feasible to use historical controls in special instances where the outcomes of a standard treatment (or no treatment) are well known and vary little for a given population.

\begin{center}\rule{0.5\linewidth}{0.5pt}\end{center}

cuidador

\begin{enumerate}
\def\labelenumi{\arabic{enumi})}
\tightlist
\item
  pessoa devidamente treinada e paga que fornece cuidados a uma pessoa com uma doença ou com uma incapacidade. 2) pessoa (frequentemente familiar ou amigo), remunerado ou não, que fornece cuidado de algum tipo, regularmente a uma pessoa com uma doença ou incapacidade
\end{enumerate}

carer

\begin{enumerate}
\def\labelenumi{\arabic{enumi})}
\tightlist
\item
  A duly trained and paid person who provides a person with a disease or disability with care. 2) A person (often a family member or friend), paid or unpaid, who regularly provides a person with a disease or disability with any form of care.
\end{enumerate}

\begin{center}\rule{0.5\linewidth}{0.5pt}\end{center}

cuidados com a saúde baseados em evidências

Os cuidados e serviços resultantes da aplicação de princípios de medicina baseada em evidências para todas as profissões associadas com cuidados à saúde, incluindo o gerenciamento e as compras de bens e serviços.

evidence-based health care

The care and services flowing from the application of the principles of evidence-based medicine to all professions associated with health care, including management and the purchasing of goods and services.

\begin{center}\rule{0.5\linewidth}{0.5pt}\end{center}

curva ROC - característica de operação do receptor

veja curva ROC

receiver operating characteristic (ROC) curve

See ROC curve.

\begin{center}\rule{0.5\linewidth}{0.5pt}\end{center}

curva ROC (Característica de Operação do Receptor)

Na avaliação de um diagnóstico ou de um teste de triagem, uma representação gráfica da relação entre a proporção de verdadeiros positivos (sensitividade) e a proporção de falsos positivos (1-especificidade) para vários pontos positivos de corte do teste. Nota: a área sob a curva ROC é uma expressão da performance do teste independentemente da população paciente e pode ser utilizada par comparar vários testes.

ROC (Receiver Operating Characteristic) curve

In the evaluation of a diagnostic or screening test, a graphical depiction of the relationship between the true positive ratio (sensitivity) and the false positive ratio (1 - specificity) for various positivity cut-off points of the test. Note: The area under the ROC curve is an expression of a test's performance independent of the patient population and can be used to compare several tests.

\begin{center}\rule{0.5\linewidth}{0.5pt}\end{center}

custo

Produto da quantidade de um recurso pelo preço unitário deste.

cost

The product of the quantity of a resource and the unit price of that resource.

\begin{center}\rule{0.5\linewidth}{0.5pt}\end{center}

custo de aquisição

Preço pago por uma instituição, agência ou pessoa na compra de uma droga, produto ou serviço de saúde.

acquisition cost

The price paid by an institution, agency or person for the purchase of a drug, device, product or health service.

\begin{center}\rule{0.5\linewidth}{0.5pt}\end{center}

custo de oportunidade

O valor dos benefícios de uma melhor opção que é oferecida quando outra opção é escolhida

opportunity cost

The value of the benefits of a better option that is given up when another option is chosen.

\begin{center}\rule{0.5\linewidth}{0.5pt}\end{center}

custo direto

Todos os custos fixos ou variáveis dos recursos (bens, serviços, et.) monopolizados na implementação de uma intervenção e gerenciando suas consequências. Nota: Drummond e seus colaboradores não suportam o uso deste conceito.

direct cost

All the fixed or variable costs of the resources (goods, services, etc.) monopolized in implementing an intervention and managing its consequences. Note: Drummond and his collaborators do not support the use of this concept.

\begin{center}\rule{0.5\linewidth}{0.5pt}\end{center}

custo direto não-médico

Todas as despesas não-médicas associadas à entrega de cuidados com a saúde e com serviços de saúde. Nota: por exemplo, despesas com o cuidado de crianças e despesas realcionadas a transporte para o recebimento de cuidados. Nota: Drummond e seus colaboradores não suportam o uso deste conceito.

direct non-medical cost

All non-medical expenses associated with the delivery of health care and health services. Note: For example, childcare expenses and expenses connected with travel to receive care. Note: Drummond and his collaborators do not support the use of this concept.

\begin{center}\rule{0.5\linewidth}{0.5pt}\end{center}

custo do capital

custo de um item que permanece útil além do período de aquisição. Nota: custos de capital incluem a depreciação da propriedade utilizada pelo item (infraestrutura física e equipamentos, por exemplo), bem como os ativos não depreciáveis, como o terreno.

capital cost

The cost of an item that remains useful beyond the acquisition period. Note: Capital costs include depreciable property, such as tangible capital assets (physical infrastructure and equipment, for example), as well as non-depreciable assets, such as land.

\begin{center}\rule{0.5\linewidth}{0.5pt}\end{center}

custo do tempo perdido

Custo das horas perdidas num trabalho ou a diminuição da produtividade devido a uma doença, incapacidade ou morte. Nota: o valor do tempo é geralmente calculado com base na média da renda.

cost of lost time

The cost of lost hours of work and decreased productivity due to disease, disability or death. Note: The value of the time is generally calculated on the basis of the average earnings.

\begin{center}\rule{0.5\linewidth}{0.5pt}\end{center}

custo futuro de cuidado com a saúde

Todos os custos resultantes do consumo de recursos adicionais após uma dada intervenção (porque isso aumenta a expectativa de vida, por exemplo).

future health care cost

All costs resulting from the consumption of additional resources following a given intervention (because it increases life expectancy, for example).

\begin{center}\rule{0.5\linewidth}{0.5pt}\end{center}

custo incremental

A diferença entre o custo da adoção de uma opção e o custo da adoção de outra opção, com a qual estiver sendo comparada.

incremental cost

The difference between the cost of an option and the cost of another option with which it is compared.

\begin{center}\rule{0.5\linewidth}{0.5pt}\end{center}

custo indireto

Todos os custos atribuidos ao valor de uma produção econômica perdida, devido a uma doença, injuria incapacitante.

indirect cost

All costs attributed to the value of economic production lost due to disease, a disabling injury or premature death.* •Not to be confused with indirect costs in accounting, where the term refers to overhead costs. Note: Drummond and his collaborators do not support the use of this concept.

\begin{center}\rule{0.5\linewidth}{0.5pt}\end{center}

custo intangível

O custo não quantificável da dor e sofrimento causados por uma doença, problema de de saúde ou intervenção. Nota: Drummond e seus colaboradores não endossam o uso deste conceito.

intangible cost

The non-quantifiable cost of pain and suffering caused by a disease, health problem or intervention. Note: Drummond and his collaborators do not support the use of this concept.

\begin{center}\rule{0.5\linewidth}{0.5pt}\end{center}

custo marginal

O custo adicional requerido para produzir uma unidade adicional de benefício. Nota: por exemplo, o custo adicional requerido para produzir uma unidade adicional de desfecho em saúde.

marginal cost

The additional cost required to produce an additional unit of benefit. Note: For example, the additional cost required to produce an additional unit of health outcome.

\begin{center}\rule{0.5\linewidth}{0.5pt}\end{center}

custo médico direto

Todos os custos fixos diretamente relacionados à entrega de cuidados com a saúde e com serviços de saúde. Nota: por exemplo, a compensação (remuneração) do médico, taxas hospitalares e medicamentos. Nota: Drummond e seus colaboradores não suportam o uso deste conceito.

direct medical cost

All the fixed or variable costs directly related to the delivery of health care and health services. Note: For example, physicians' compensation, hospital fees and drugs. Note: Drummond and his collaborators do not support the use of this concept.

\begin{center}\rule{0.5\linewidth}{0.5pt}\end{center}

custo médio

custo total de um tratamento ou programa dividido pelo número total de unidades de tratamento proporcionadas

average cost

The total cost of a treatment or programme, divided by the total number of treatment units provided.

\begin{center}\rule{0.5\linewidth}{0.5pt}\end{center}

custo não relacionado

Os custos não especificamente atribuidos a um caminho terapêutico e suas consequências.

unrelated cost

The costs not specifically attributable to a therapeutic pathway and its consequences.

\begin{center}\rule{0.5\linewidth}{0.5pt}\end{center}

custo por QALY

Medida utilizada na análise de custo-utilidade para comparar programas, intervenções ou tratamentos, expressos como um custo unitário por unidade de desfecho (efetividade), ajustada com base na qualidade de vida.

cost per QALY

A measure used in cost-utility analysis to compare programmes, interventions or treatments, expressed as a monetary cost per unit of outcome (effectiveness), adjusted on the basis of quality of life.

\begin{center}\rule{0.5\linewidth}{0.5pt}\end{center}

custo recorrente

Todos os custos que ocorrem periodicamente - anualmente, por exemplo.

recurrent cost

All costs that occur periodically---each year, for example.

\hypertarget{d}{%
\chapter*{D}\label{d}}
\addcontentsline{toc}{chapter}{D}

dados binários

dados que podem ter apenas dois valores, tais como morto ou vivo, fumante ou não fumante, presente ou ausente, homem ou mulher. Este é um caso específico de dados categóricos, onde o número de categorias é igual a 2.

binary data

Data that can have only two values, such as dead or alive, smoker or non-smoker, present or absent, or man or woman. Note: This is a specific case of a categorical data, where the number of categories is equal to 2.

\begin{center}\rule{0.5\linewidth}{0.5pt}\end{center}

dados categóricos

dados que podem ser classificados em duas ou mais categorias. Nota: quando existe uma ordem para as categorias, o termo usado é ``dados ordenados'' (i.e.~estágios de cancer ou nível educacional). Quando não existe uma ordem, o termo usado é ``dados nominais'' (Exemplo: grupo sanguíneo, estado civil ou grupo étnico). Não há relação matemática entre os valores (Exemplo: pessoa com câncer no estágio IV não está duas vezes mais doente do que uma pessoa com câncer no estágio II).

categorical data

Data that can be classified into two or more categories. Note: When there is an order to the categories, the term used is ordinal data (e.g.~stages of cancer or level of education). When there is no order, the term used is nominal data (e.g.~blood group, civil status or ethnic origin). There is no mathematical relationship between the values (e.g.~a person with Stage IV cancer is not twice as sick as a person with Stage II cancer).

\begin{center}\rule{0.5\linewidth}{0.5pt}\end{center}

dados contínuos

Dado numérico com um número potencialmente infinito de valores possíveis numa faixa contínua. Nota: altura, peso e pressão sanguínea são exemplos de dados contínuos

continuous data

Numerical data with a potentially infinite number of possible values on a continuum. Note: Height, weight and blood pressure are examples of continuous data.

\begin{center}\rule{0.5\linewidth}{0.5pt}\end{center}

dados de paciente individual

Numa revisão sistemática, dados brutos do participante de cada estudo em cada ensaio incluído em oposição aos dados agregados nos grupos comparados em cada estudo.

individual patient data

In a systematic review, raw data for each study participant in each included trial, as opposed to aggregate data on the groups compared in each study.

\begin{center}\rule{0.5\linewidth}{0.5pt}\end{center}

dados dicotômicos

Veja dados binários

dichotomous data

See binary data.

\begin{center}\rule{0.5\linewidth}{0.5pt}\end{center}

dados discretos

Um conjunto de valores é chamado de discreto se os valores/observações pertencentes a esse conjunto são distintos e separados, isto é, podem ser contados (1,2,3\ldots). Nota 1: exemplos incluem o número de paciente das cirurgias de um médico; sexo (masculino, feminino); grupo sanguíneo (O,A,B,AB). Nota 2: dados discretos são a contrapartida de dados contínuos.

discrete data

A set of data is said to be discrete if the values/observations belonging to it are distinct and separate, i.e.~they can be counted (1,2,3,\ldots.). Note 1: Examples include the number of patients in a doctor's surgery; sex (male, female); blood group (O, A, B, AB). Note 2: Discrete data is the counterfactual to continuous data.

\begin{center}\rule{0.5\linewidth}{0.5pt}\end{center}

dados ordinais

Dados classificados em duas ou mais categorias, quando existe uma ordem natural das mesmas, tais como não fumantes, ex-fumantes, fumantes leves e fumantes pesados. Sinônimo: dados categóricos ordenados

ordinal data

Data classified into two or more categories where there is a natural order to the categories, such as non-smokers, ex-smokers, light smokers and heavy smokers. Syn.: ordered categorical data

\begin{center}\rule{0.5\linewidth}{0.5pt}\end{center}

declaração de conflito de interesse

declaração de uma pessoa contribuindo para um relatório ou revisão literária de qualquer relação financeira ou pessoal que pudesse influenciar os resultados ou a interpretação dos mesmos nesse relatório ou revisão.

conflict of interest declaration

A statement by a person contributing to an evaluation report or a literature review of any financial or personal relationship that could influence the results or the interpretation of the results of this report or review.

\begin{center}\rule{0.5\linewidth}{0.5pt}\end{center}

Desconto

Um processo matemático usado para trazer custos e benefícios futuros ao seu valor presente. Nota: Estes ajustes implicam que os custos futuros e benefícios tenham um valor menor do que os mesmos custos e benefícios no presente.

discounting

A mathematical process used to bring future costs and benefits to their present value. Note: These adjustments imply that future costs and benefits have less value than the same costs and benefits in the present.

\begin{center}\rule{0.5\linewidth}{0.5pt}\end{center}

desenvolvimento de consenso

Recomendações ou diretrizes práticas produzidas através de consenso especializado

consensus development

The recommendations or practice guideline produced through expert consensus.

\begin{center}\rule{0.5\linewidth}{0.5pt}\end{center}

desfecho

Um componente mensurável observado após a aplicação de uma intervenção.

outcome

A measurable component observed after an intervention has been applied.

\begin{center}\rule{0.5\linewidth}{0.5pt}\end{center}

desfecho clínico

desfecho da maior importância para o paciente, determinado com base no problema de saúde sendo estudado (Exemplo, a frequencia de cura de fraturas osteoporéticas ou de úlceras gátricas, e suas taxas de recaída). Nota: os desfechos clínicos diferem dos desfechos na saúde pois o primeiro está relacionado com o efeito da doença enquanto que o segundo esta relacionado com o efeito na saúde (saúde pública, por exemplo).

clinical outcome

An outcome of major importance for the patient, determined on the basis of the health problem being studied (e.g.~the frequency of healing of osteoporotic fractures or gastric ulcers, and their relapse rates). Note: The clinical outcome differs from the health outcome in that the former has to do with the effect on the disease, while the latter has to do with the effect on health (public health, for example).

\begin{center}\rule{0.5\linewidth}{0.5pt}\end{center}

desfecho de pesquisa

Um campo de questionamento com foco em problemas de saúde ou doenças e que visa essencialmente a avaliação de cuidados entregues por equipes multidisciplinares estabelecidas de acordo com as definições gerais do mundo real. Nota: este tipo de pesquisa avalia uma ampla faixa de indicadores, incluindo mortalidade, morbidade, estado funcional, bem estar mental e outros aspectos da saude relacionados à qualidade de vida. Pode também incluir uma avaliação do impacto econômico de desfechos clínicos - por exemplo, uma análise de custo-efetividade ou custo-utilidade. Pode ser baseada em qualquer método de coleta de dados originais ou para sumarizar dados de estudos originais.

outcomes research

A field of inquiry that focusses on health problems or diseases and that is essentially aimed at evaluating the care delivered by multidisciplinary teams in general, real-world settings. Note: This type of research looks at a wide range of indicators, including mortality, morbidity, functional status, mental well-being and other aspects of health-related quality of life. It may also include an evaluation of the economic impact of clinical outcomes---a cost-effectiveness or cost-utility analysis, for example. It may be based on any method for collecting original data or for summarizing data from original studies.

\begin{center}\rule{0.5\linewidth}{0.5pt}\end{center}

desfecho intermediário

Veja desfecho substituto.

intermediate outcome

See surrogate endpoint.

\begin{center}\rule{0.5\linewidth}{0.5pt}\end{center}

desfecho na saúde

O resultado de uma intervenção ( ou da ausência de uma intervenção) na saúde ou no bem-estar de um indivíduo ou de uma população. Nota: a intervenção pode ser um procedimento clínico, intervenção social, política de saúde, programa público, etc. Nota: o desfecho na saúde difere do desfecho clínico no qual o primeiro está relacionado com o efeito na saúde (saúde pública, por exemplo), enquanto que o segundo está relacionado com o efeito da doença.

health outcome

The result of an intervention (or the absence of an intervention) on the health or well-being of a subject or a population. Note: The intervention may be a clinical procedure, social intervention, health policy, public programme, etc. Note: The health outcome differs from the clinical outcome in that the former has to do with the effect on health (public health, for example), while the latter has to do with the effect on the disease.

\begin{center}\rule{0.5\linewidth}{0.5pt}\end{center}

desfecho substituto

Um indicador que, quando não sendo de interesse direto para o paciente, pode refletir desfechos importantes. Nota: por exemplo, a pressão sanguínea não é do interesse clínico direto para o paciente, mas é frequentemente usada como critério de avaliação em ensaios clínicos porque é um fator de risco para acidentes vasculares e ataques cardíacos. Um desfecho intermediário é frequentemente um marcador psicológico ou bioquímico e pode ser medido rápido e facilmente, e isso é considerado de grande valor preditivo. É frequentemente usado quando a observação do desfecho clínico requer um longo acompanhamento. Nota: desfecho intermediário não é sinônimo de desfecho substituto. No entanto, um desfecho intermediário pode tornar-se um desfecho substituto se for mais fácil de ser medido do que um critério clínico ou se houver uma relação estatística entre a ocorrência do indicador clínico e a ocorrência do desfecho substituto, ou se houver uma relação permitindo predizer o efeito do fator estudado no indicador clínico, com base no efeito observado no desfecho substituto.

surrogate endpoint

An indicator that, while not being of direct interest for the patient, may reflect important outcomes. Note: For example, blood pressure is not of direct clinical interest to the patient, but is often used as an evaluation criterion in clinical trials because it is a risk factor for stroke and heart attacks. An intermediate outcome is often a physiological or biochemical marker that can be quickly and easily measured, and that is considered to have great predictive value. It is often used when observation of clinical outcomes requires long follow-up. Note: Intermediate outcome is not a synonym for surrogate endpoint. However, an intermediate outcome can become a surrogate endpoint if it is easier to measure than a clinical criterion or if there is a statistical relationship between the occurrence of the clinical outcome indicator and the occurrence of the surrogate endpoint, or if there is a relationship allowing for prediction of the effect of the factor studied on the clinical indicator, on the basis of the observed effect on the surrogate endpoint.

\begin{center}\rule{0.5\linewidth}{0.5pt}\end{center}

desinvestimento

A redução deliberada e sistemática de investimento numa tecnologia de saúde questionável ou, comparativamente, de menor valor. Nota: O investimento é frequentemente governamental, ou seja, susidiando os custos do cuidado com a saúde.

desinvestment

The deliberate and systematic reduction of funding for a health technology of questionable or comparatively low value. Note: The funding is often government funding, i.e.~subsidisation of the costs of health care.

\begin{center}\rule{0.5\linewidth}{0.5pt}\end{center}

diagnóstico

A determinação da natureza de uma doença através de seus sinais e sintomas, e pelos resultados de investigações. Ao se estabelecer um diagnóstico também se excluem outras doenças e fatores causais das manifestações clínicas.

diagnosis

The determination of the nature of a disease by means of its signs and symptoms and the results of investigations.Note: The establishment of a diagnosis also means ruling out other diseases and causal factors for the clinical manifestations.

\begin{center}\rule{0.5\linewidth}{0.5pt}\end{center}

diferença da média ponderada

Em meta-análise, quando os resultados do estudo medidos usando a mesma escala são combinados,e a diferença entre as duas médias é pesada pela precisão do estudo. Nota: a precisão da estimativa do estudo do efeito pode, por exemplo, corresponder ao inverso da variância. (conceito relacionado: diferença da média padrão).

weighted mean difference (WMD)

In a meta-analysis, when study results measured using the same scale are being combined, the difference between two means, weighted by the precision of the study. Note: The precision of the study's estimate of effect may, for example, correspond to the inverse of the variance. (Related concept: standardised mean difference)

\begin{center}\rule{0.5\linewidth}{0.5pt}\end{center}

diferença de risco

O valor da diferença entre a probabilidade de que um evento ocorrerá no grupo exposto a um dado fator e a probabilidade de que esse evento ocorrerá no grupo não exposto a esse fator. Nota: por exemplo, se os resultados de um ensaio forem que a probabilidade de morte foi de 25\% no grupo controle e de 10\% no grupo experimental, a redução de risco absoluto será de 0,25-0,10=0,15. Sinônimo: redução de risco absoluto (conceitos relacionados: número necessário para tratamento, razão de chances e redução do risco relativo).

risk difference

The value of the difference between the probability that an event will occur in the group exposed to a given factor and the probability that this event will occur in the group not exposed to this factor. Note: For example, if the results of a trial were that the probability of death was 25\% in the control group and 10\% in the experimental group, the absolute risk reduction would be 0.25 -- 0.10 = 0.15. Syn.: absolute risk reduction. (Related concepts: number needed to treat, odds ratio and relative risk reduction)

\begin{center}\rule{0.5\linewidth}{0.5pt}\end{center}

diferença padronizada da média

A diferença entre duas médias, dividida por uma estimativa do desvio padrão intra-grupo. Nota: quando uma variável contínua (tal como a dor) é medida numa variedade de estudos transversais (usando diferentes escalas), pode ocorrer a impossibilidade de se comparar ou combinar os resultados do estudo numa meta análise. Se os efeitos são expressos como um valor padronizado, os resultados podem ser combinados, desde que não sejam mais expressos com unidades de medida. Diferenças de médias padronizadas são às vezes referidas como um ``índice d''. * *Conceito introduzido por J. Cohen (1988). Outra forma de referência: diferença padronizada da média (SMD).

standardised mean difference (SMD)

The difference between two means, divided by an estimate of the within-group standard deviation. Note: When a continuous variable (such as pain) is measured in a variety of ways across studies (using different scales), it may prove impossible to compare or combine the study results in a meta-analysis. If the effects are expressed as a standardised value, the results can be combined, since they are no longer expressed as units of measurement. Standardised mean differences are sometimes referred to as a ``d index.''* *Concept introduced by J. Cohen (1988). Alternative spelling: standardized mean difference (SMD).

\begin{center}\rule{0.5\linewidth}{0.5pt}\end{center}

difusão de inovação

Um processo onde uma inovação é comunicada ao longo do tempo através de certos canais dos membros de um sistema social. Nota: no caso de tecnologias médicas, são os fatores que influenciam a adoção de uma nova tecnologia dentro de um sistema de cuidados de saúde.

diffusion of inovation

A process whereby an innovation is communicated over time through certain channels among the members of a social system. Note: In the case of medical technologies, factors influencing the adoption of a new technology within the health care system.

\begin{center}\rule{0.5\linewidth}{0.5pt}\end{center}

difusão de tecnologia

Veja difusão de inovação.

technology diffusion

See diffusion of innovation.

\begin{center}\rule{0.5\linewidth}{0.5pt}\end{center}

dilema

Uma escolha forçada entre dois cursos de ações, ambas com desvantagens.

dilemma

A forced choice between two courses of actions, both of which have disadvantages.

\begin{center}\rule{0.5\linewidth}{0.5pt}\end{center}

diretriz clínica

diretriz desenvolvida usando-se um método sistemático para ajudar os praticantes e os pacientes a tomar decisões sobre um cuidado de saúde adequado. Nota: o desenvolvimento de diretrizes clínicas pode usar métodos de avaliação ou ser baseados em relatórios de avaliação tecnológica.

clinical guideline

A guideline developed using a systematic method to help practitioners and patients make decisions about appropriate health care. Note: The development of clinical guidelines may use technology assessment methods or be based on technology assessment reports.

\begin{center}\rule{0.5\linewidth}{0.5pt}\end{center}

disposição em pagar

O montante máximo que uma pessoa está disposta a pagar: (a) para obter um bom estado de saúde ou um desfecho em particular, ou para aumentar sua probabilidade de ocorrência, ou (b) para evitar um mau estado de saúde ou desfecho, ou para reduzir sua probabilidade de ocorrência.

willingness to pay (WTP)

The maximum amount that a person is willing to pay: (a) to achieve a good health state or particular outcome, or to increase its probability of occurrence; or (b) to avoid a bad health state or outcome, or to decrease its probability of occurrence.

\begin{center}\rule{0.5\linewidth}{0.5pt}\end{center}

disseminação

Um processo de transmissão de informação a audiências-alvo.

dissemination

A process for transmission of information to target audiences.

\begin{center}\rule{0.5\linewidth}{0.5pt}\end{center}

distribuição

Em estatística, uma função matemática descrevendo a frequência dos valoes de uma variável.

distribution

In statistics, a mathematical function describing the frequency of each value of a variable.

\begin{center}\rule{0.5\linewidth}{0.5pt}\end{center}

distribuição de probabilidades

Uma função matemática descrevendo as probabilidades associadas com cada valor. Nota: por exemplo, as distribuições normal, qui-quadrado, binomial ou Poisson.

probability distribution

A mathematical function describing the probabilities associated with each value. Note: For example, normal, chi square, binomial or Poisson distribution.

\begin{center}\rule{0.5\linewidth}{0.5pt}\end{center}

dominância

A superioridade de uma opção que proporciona menores custos do que outra opção e que tem benefícios iguais ou maiores do que a outra opção, ou que proporciona custos iguais ao da outra opção, porém com maiores benefícios.

dominance

The superiority of an option that entails lower costs than another option and has benefits equal to or greater than the other option, or that entails costs equal to those of another option and has greater benefits than the other option.

\begin{center}\rule{0.5\linewidth}{0.5pt}\end{center}

duplo-cego

Um ensaio no qual as duas partes do ensaio (objetos e investigadores) não sabem como os tratamentos são alocados. Nota: este método reduz o risco de viés ao se medir os desfechos. (coceitos relacionados: cegamento, cegamento único, cegamento triplo). Sinônimo: mascaramento. Nota: os termos cegamento e mascaramento são sinônimos e o termo mascaramento duplo é preferido em ensaios oftálmicos, ao invés de duplamento cego.

double-blind

A trial in which the two parties in the trial (the subjects and the investigators) do not know how the treatments are allocated. Note: This method reduces the risk of bias when measuring outcomes. (Related concepts: blinding, single-blind and triple-blind). Syn.: masking. Note: The terms blinding and masking are synonymous and the term double-masked is preferred over double-blind for ophthalmology trials.

\hypertarget{e}{%
\chapter*{E}\label{e}}
\addcontentsline{toc}{chapter}{E}

economia na saúde

Uma disciplina científica que busca aplicar os princípios e regras da Econômia no campo do cuidado à saúde. Nota: a econômia na área da saúde inclui a avaliação do sistema de saúde e de políticas da saúde de um ponto de vista econômico, da demanda e do fornecimento de cuidados com a saúde, e de tecnologias e de intervenções médicas; também inclui o planejamento e a organização de recursos do sistema de saúde, a consideração e a avaliação de determinantes da saúde, e a análise da performance do sistema de saúde em termos da eficiência da equidade e da alocação.

health economics

A scientific discipline that seeks to apply the principles and rules of economics in the health care field. Note: Health economics includes evaluation of the health system and of health policies from an economic perspective, of the demand for and supply of health care, and of medical technologies and interventions; it also includes health system resources planning and organization, consideration and evaluation of the determinants of health, and analysis of health system performance in terms of equity and allocative efficiency.

\begin{center}\rule{0.5\linewidth}{0.5pt}\end{center}

efeito adverso

um efeito nocivo ou não intencionado de uma intervenção .

adverse effect

A noxious or unintended effect of an intervention.

\begin{center}\rule{0.5\linewidth}{0.5pt}\end{center}

efeito colateral

Veja efeito adverso e evento adverso. Nota: o termo efeito colateral, frequentemente utilizado como efeito adverso, é impreciso, pois um efeito pode ser secundário (um efeito colateral)de uma intervenção, sem ser nocivo.

side effect

See adverse effect and adverse event. Note: The term side effect, often used for adverse effect, is imprecise, in that an effect can be secondary (a side effect) to an intervention, without being noxious.

\begin{center}\rule{0.5\linewidth}{0.5pt}\end{center}

efeito do tratamento

O efeito no estado de saúde ou no bem estar dos indivíduos atribuído apenas a um tratamento ou intervenção. Nota: investigadores buscam estimar o verdadeiro efeito de um tratamento ou intervenção calculando a diferença entre o desfecho obtido no grupo experimental e no grupo controle. Conceito relacionado: tamanho do efeito.

treatment effect

The effect on the subjects' health status or well-being attributable only to a treatment or intervention. Note: Investigators seek to estimate the true effect of a treatment or intervention by calculating the difference between the outcome obtained in the experimental group and the control group.
Related concept: effect size

\begin{center}\rule{0.5\linewidth}{0.5pt}\end{center}

efeito placebo

O efeito no estado de saúde do paciente (melhora ou piora) que pode ocorrer em função da expectativa do paciente com relação ao efeito que uma dada intervenção traria. Nota: o efeito placebo é independente do efeito sendo estudado (farmacológico, cirúrgico, etc.). Quando um placebo é dado ao grupo controle, ambos os grupos experimentarão esse efeito, tornando possível isolar o efeito específico da intervenção.

placebo effect

The effect on a patient's health status (improvement or worsening) that may occur due to the patient's expectation that a particular intervention will have an effect. Note: The placebo effect is independent of the effect being studied (pharmacological, surgical, etc.). When a placebo is given to the control group, both groups will experience this effect, making it possible to isolate the specific effect of the intervention.

\begin{center}\rule{0.5\linewidth}{0.5pt}\end{center}

efetividade

O benefício de se usar uma tecnologia, programa ou intervenção para endereçar um problema específico em condições gerais ou de rotina, ao invés de condições controladas, por exemplo, por um médico num hospital ou na casa de um paciente. Sinônimo: efetividade clínica.

effectiveness

The benefit of using a technology, programme or intervention to address a specific problem under general or routine conditions, rather than under controlled conditions, for example, by a physician in a hospital or by a patient at home. Syn.: clinical effectiveness.

\begin{center}\rule{0.5\linewidth}{0.5pt}\end{center}

efetividade clínica

veja efetividade

clinical effectiveness

See effectiveness.

\begin{center}\rule{0.5\linewidth}{0.5pt}\end{center}

eficácia

O benefício de se usar uma tecnologia, programa ou intervenção para tratar um problema particular em condições ideais - por exemplo, no contexto da pesquisa em laboratório, ou num protocolo rigoroso para um ensaio clinico randomizado.

efficacy

The benefit of using a technology, programme or intervention to treat a particular problem under ideal conditions---for example, in the context of research in a laboratory or a rigourous protocol for a randomized clinical trial.

\begin{center}\rule{0.5\linewidth}{0.5pt}\end{center}

eficácia diagnóstica de um teste

O grau para o qual um teste, possuindo suas propriedades técnicas e acurácia, pode ser usado para se fazer um diagnóstico; a utilidade de um teste é também expressa em termos da sua influência na escolha do tratamento e na saúde de um indivíduo ou de uma população.

diagnostic test efficacy

The degree to which a test, owing to its technical properties and accuracy, can be used to make a diagnosis; the usefulness of a test is also expressed in terms of its influence on the choice of treatment and on the health of an individual or a population.

\begin{center}\rule{0.5\linewidth}{0.5pt}\end{center}

eficiência

A habilidade de obter o máximo benefício possível, à partir dos recursos disponíveis. Nota: é frequentemente definida como a razão entre o desfecho obtido e os recursos utilizados. (conceitos relacionados: eficiência alocativa e eficiência técnica).

efficiency

The ability to obtain the maximum possible benefit, given the resources available. Note: It is often the simple ratio between the outcome obtained and the resources used. (Related concepts: allocative efficiency and technical efficiency)

\begin{center}\rule{0.5\linewidth}{0.5pt}\end{center}

eficiência alocativa

A eficiência alocativa (um conceito econômico) refere-se à forma como as diferentes entradas de recursos são combinadas para produzir uma mistura de saídas diferentes.

allocative efficiency

The optimal distribution of healthcare resources required to produce a given level of health outcomes, taking into account consumer preferences.

Note 1: Related terms include efficiency and technical efficiency.

\begin{center}\rule{0.5\linewidth}{0.5pt}\end{center}

eficiência técnica

A habilidade em produzir a maioria ou os melhores efeitos com um nível especificado de recursos

technical efficiency

The ability to produce the most or the best effects with a specified level of resources.

\begin{center}\rule{0.5\linewidth}{0.5pt}\end{center}

EMBASE (Base de Dados da Excerpta Medica)

Uma base de dados biomédica e farmacológica cobrindo 3.500 jornais de 110 países, publicados desce 1970. Nota: esta base de dados contém um grande número de publicações européias.

EMBASE (Excerpta Medica Database)

A biomedical and pharmacological database covering 3,500 journals from 110 countries published since 1970. Note: This database contains a large number of European publications.

\begin{center}\rule{0.5\linewidth}{0.5pt}\end{center}

empírico

Resultados baseados na experiência ou observação, ao invés de somente por raciocínio.

empirical

Results based on experience or observation, rather than on reasoning alone.

\begin{center}\rule{0.5\linewidth}{0.5pt}\end{center}

EMTREE (Árvore da Excerpta Medica)

Vocabulário controlado organizado em um dicionário para a base de dados EMBASE.

EMTREE (Excerpta Medica Tree)

Controlled vocabulary organized into a thesaurus for the EMBASE database.

\begin{center}\rule{0.5\linewidth}{0.5pt}\end{center}

ensaio aberto

Um ensaio clínico no qual o investigador e o indivíduo sabem quem está recebendo e qual é a intervenção. Nota: este termo é também isado quando apenas uma intervenção está sendo estudada, em estudos não cegados de longo termo de monitoração ou quando a cegagem é difícil ou não ética (comparção de um tratamento cirúrgico com um tratamento médico, por exemplo).

open-label trial

A clinical trial in which the investigator and the subject know who is receiving which intervention. Note: This term is also used when only one intervention is being studied, for unblinded long-term follow-up studies or when blinding is difficult or unethical (comparison of a surgical treatment and a medical treatment, for example).

\begin{center}\rule{0.5\linewidth}{0.5pt}\end{center}

ensaio clínicamente controlado (ECC)

Um estudo experimental no qual os investigadores comparam os desfechos obtidos num grupo de elementos que passaram pela intervenção com os desfechos dos elemento que tem características semelhantes, mas que não passaram pela intervenção. Nota: nem todos os ensaios clínicos são randomizados, mas todos os ensaios randomizados são ensaios clínicos.

controlled clinical trial (CCT)

An experimental study in which the investigators compare the outcomes obtained in a group of subjects who have undergone an intervention with the outcomes in subjects who have similar characteristics, but have not undergone the intervention. Note: Not all clinical trials are randomized, but all randomized trials are clinical trials.

\begin{center}\rule{0.5\linewidth}{0.5pt}\end{center}

ensaio clínico

um estudo experimental rigoroso e controlado, que inclui elementos sadios e doentes, em que se deseja avaliar os desfechos de uma ou mais intervenções profiláticas terapêuticas. Nota: cada ensaio é projetado para responder questões científicas específicas. Em estudos experimentais (ensaios) a exposição a drogas, terapêuticas ou manobras preventivas são impostas pelo investigador. O investigador decide, através de algum método, incluir através de uma seleção randômica, quais elentos serão expostos; em estudos observacionais, a exposição é determinada pelo próprio elemento, naturalmente. Nota: a diferença entre um ensaio e um estudo é relativamente complexa, mas podem ser diferenciados da seguinte maneira; um ensaio suporta uma intervenção que poderá alterar o curso dos fatos e que se aplica a humanos, enquanto que um estudo examina os desfechos de uma intervenção ou da exposição a um dado fator, ou pode ter outros objetivos, tal como a observação do curso de uma doença. Um ensaio pode ser parte de um estudo, mas a recíproca não é verdadeira.

clinical trial

A rigourous and monitored experimental study* that includes healthy or ill human subjects and that is aimed at evaluating the outcomes of one or more prophylactic or therapeutic interventions. Note: Each trial is designed to answer specific scientific questions. In experimental studies (trials), exposure to a drug or to a therapeutic or preventive manoeuvre is imposed by the investigator. The investigator decides, through some method, including random selection, who will be exposed; in observational studies, whether there is exposure is determined by the subject himself/herself or naturally.
Note: The difference between a trial and a study is quite fuzzy, but they can be distinguished as follows: a trial bears on an intervention that will change the course of things and that applies to humans, while a study examines the outcomes of an intervention or of exposure to a given factor, or may have other aims, such as observing the course of a disease. A trial may be part of a study, but the reverse is not true.

\begin{center}\rule{0.5\linewidth}{0.5pt}\end{center}

ensaio controlado não randomizado

Um ensaio clínico no qual os indivíduos são divididos entre um grupo experimental e um grupo de controle, usando um método que não envolve randomização - com base nas restrições práticas dos investigadores, por exemplo - ou outros métodos tal como estabelecer tratamentos alternados. Outra forma de referência: ensaio controlado não-randomizado.

nonrandomised controlled trial

A clinical trial in which the subjects are divided between an experimental group and a control group using a method that does not involve randomisation---on the basis of the investigators' practical constraints, for example---or other methods, such as alternate assignment of treatments. Alternate spelling: non-randomized controlled trial

\begin{center}\rule{0.5\linewidth}{0.5pt}\end{center}

ensaio de cruzamento

Ensaio no qual os elementos recebem, em sequência, o tratamento (ou a intervenção de controle) e então, após um tempo especificado, a intervenção de controle (ou tratamento). Nota: neste tipo de estudo, os elementos fazem o controle de si mesmos e a randomização é usada para determinar a ordem pela qual um elemento recebe o tratamento ou a intervenção de controle. Períodos de lavagem são usualmente incluidos, desde que efeitos do primeiro tratamento possam permanecer ativos quando o segundo período for executado. (refresh?)

crossover trial

A trial in which the subjects receive, in sequence, the treatment (or control intervention), then, after a specified time, the control intervention (or treatment). Note: In this type of study, the subjects are their own controls, and randomisation is used to determine the order in which a subject receives the treatment and the control intervention. Washout periods are usually included, since otherwise the effects of the first treatment could carry over into the period when the second is given.

\begin{center}\rule{0.5\linewidth}{0.5pt}\end{center}

ensaio de grupo paralelo

Um ensaio que compara dois grupos contemporâneos, um recebendo o tratamento de interesse e outro sendo o grupo controle (num ensaio clínico randomizado, por exemplo). Nota: alguns desses ensaios tem mais do que um grupo de tratamento; outros comparam dois grupos de tratamento, cada um sendo o controle do outro.

parallel group trial

A trial that compares two contemporaneous groups, one receiving the treatment of interest, the other being the control group (in a randomised clinical trial, for example). Note: Some of these trials have more than one treatment group; others compare two treatment groups, each being a control for the other.

\begin{center}\rule{0.5\linewidth}{0.5pt}\end{center}

ensaio equivalente

Um ensaio com o objetivo primário de mostrar que a diferença entre respostas a duas ou mais intervenções é clinicamente insignificante. Nota: este objetivo é usualmente alcançado ao se demonstrar que a diferença real entre os tratamentos tende a ficar numa zona pré-determinada na qual os desfechos são considerados equivalentes.

equivalence trial

A trial with the primary objective of showing that the difference between the responses to two or more interventions is clinically insignificant. Note: This objective is usually achieved by demonstrating that the actual difference between the treatments is likely to fall within a pre-determined zone in which the outcomes are deemed clinically equivalent.

\begin{center}\rule{0.5\linewidth}{0.5pt}\end{center}

ensaio explanatório

Um ensaio que mede, em condições ideais, os efeitos de intervenções, usando cuidadosamente os sujeitos selecionados no estabelecimento de uma pesquisa clínica. Nota: este tipo de ensaio é frequentemente usado para avaliar a eficácia teórica de intervenções. Antônimo: ensaio pragmático.

explanatory trial

A trial that measures under ideal conditions the effects of interventions, using carefully selected subjects in a clinical research setting. Note: This type of trial is often used to evaluate the theoretical efficacy of interventions. Ant.: pragmatic trial

\begin{center}\rule{0.5\linewidth}{0.5pt}\end{center}

ensaio n de 1

Um ensaio clínico onde um único indivíduo é o total da população em estudo. Nota: a alocação randômica pode ser usada para determinar a ordem pela qual o indivíduo receberá a intervenção experimental e a experimentação de controle.

n of 1 trial

A clinical trial in which a single subject is the total population for the study. Note: Random allocation may be used to determine the order in which the subject will receive the experimental intervention and the control intervention.

\begin{center}\rule{0.5\linewidth}{0.5pt}\end{center}

ensaio pragmático

Um ensaio que mede o efeito de uma intervenção na prática da rotina clínica, para avaliar a efetividade real da intervenção. Antônimo: ensaio explanatório.

pragmatic trial

A trial that measures the effects of an intervention in routine clinical practice, to evaluate the intervention's actual effectiveness. Ant.: explanatory trial.

\begin{center}\rule{0.5\linewidth}{0.5pt}\end{center}

ensaio sequencial

Um método estatístico que permite o encerramento de um experimento assim que for obtida uma resposta com a precisão desejada, com cada desfecho sendo examinado assim que ocorrer e sendo adicionado aos desfechos anteriores. Nota: as principais vantagens de um ensaio sequencial são: 1) é mais curtos do que um ensaio de duração fixa quando existe uma grande diferença na efetividade das intervenções sendo comparadas ou 2) um ensaio pode ser encerrado rapidamente quando existe um risco à saúde. Assim, com este enfoque, a atribuição dos tratamentos geralmente será conhecida antes do estudo, sendo requeridos métodos especiais para a análise e disseminação dos resultados internos. Conceito relacionado: meta análise.

sequential trial

A statistical method that allows for ending an experiment as soon as a response with the desired precision is obtained, with each outcome being examined as soon as it occurs and being added to all the previous outcomes. Note: The main advantage of a sequential trial is as follows: 1) it is shorter than a fixed-length trial when there is a large difference in the effectiveness of the interventions being compared; or 2) a trial can be terminated quickly when there is a health risk. Since, with this approach, the assignment of the treatments generally must be known before the study, special methods for the analysis and dissemination of interim results are required. Related concept: meta-analysis

\begin{center}\rule{0.5\linewidth}{0.5pt}\end{center}

ensaio simples em grande escala

Um ensaio randomizado, pragmático, caracterizado pelo recrutamento de um grande número de indivíduos, com amplos critérios para a inclusão de pacientes, múltiplos locais de estudo, um processo simplificado de coleta de dados e uso de registro eletrônico. Nota: os objetivos desse tipo de ensaio incluem a detecção de efeitos do tratamento pequenos ou moderados, pela obtenção de dados na efetividade prática do tratamento e pelo aumento da validade externa das descobertas.

large simple trial

A randomised pragmatic trial characterized by the recruitment of a large number of subjects, broad patient inclusion criteria, multiple study sites, a simplified data collection process and use of electronic registers. Note: The purposes of this type of trial include detecting small and moderate treatment effects, obtaining data on the practical effectiveness of the treatment and improving the external validity of the findings.

\begin{center}\rule{0.5\linewidth}{0.5pt}\end{center}

ensaio terapêutico

veja ensaio clínico.

therapeutic trial

See clinical trial.

\begin{center}\rule{0.5\linewidth}{0.5pt}\end{center}

ensaios clinícos randomizados

Um estudo comparando pelo menos duas intervenções, nas quais os participantes elegíveis são alocados aleatoriamente para o grupo de intervenção, ou grupos, e o grupo controle. Sigla: ECR. Nota 1: O controle pode ser uma prática padrão, um placebo, outra intervenção ativa ou nenhuma intervenção. Nota 2: Os participantes podem ser indivíduos ou grupos (por exemplo, unidade de randomização em um ensaio controlado randomizado em cluster). Ortografia alternativa: ensaio controlado randomizado

randomized controlled trials

A study comparing at least two interventions, in which the eligible participants are allocated randomly to the intervention group, or groups, and the control group. Acronym: RCT. Note 1: The control may be a standard practice, a placebo, other active intervention, or no intervention. Note 2: Participants may be individuals or groups (eg unit of randomization in a cluster randomized controlled trial). Alternate spelling: randomised controlled trial

\begin{center}\rule{0.5\linewidth}{0.5pt}\end{center}

epidemiologia

O estudo da distribuição e determinantes de estados relacionados à saúde ou eventos em populações especificadas.

epidemiology

The study of the distribution and determinants of health-related states or events in specified populations.

\begin{center}\rule{0.5\linewidth}{0.5pt}\end{center}

equidade

A alocação equitativa de recursos ou tratamentos entre distintos indivíduos ou grupos, de tal forma que cada um possui ou recebe o que deveria. Outra definição: o uso da alocação de recursos baseado no princípio da justiça distributiva. No entanto, a segunda definição envolve vários conceitos que definem essas regras com maior precisão (méerito, valor social, a maior contribuição, oportunidade equalitária, tratamento equalitário, oportunidade equalitária, prioridade ao que estiver em maior desvantagem, aposta padrão, etc.). Nota: equidade vertical significa que as pessoas com maior necessidade de serviços recebem a maioria desses serviços, enquanto que a equidade horizontal significa que pessoas com necessidades semelhantes recebem serviços semelhantes.

equity

The fair allocation of resources or treatments among different individuals or groups, such that they each get what they are owed or what they are entitled to. Other definition: The use of resource allocation rules based on the distributive justice principle. However, the latter involves various concepts that define these rules more precisely (merit, social value, the greater good, equal opportunity, equal treatment, priority to the most disadvantaged, standard gamble, etc.). Note: Vertical equity means that the people in the greatest need of services receive the most services, and horizontal equity means that people who have similar needs receive similar services.

\begin{center}\rule{0.5\linewidth}{0.5pt}\end{center}

equilíbrio

Um estado de genuina incerteza tal como saber qual de duas ou mais alternativas em tecnologia de saúde devem constituir o padrão de cuidados, considerando os benefícios e malefícios. Nota 1: Equipoise é largamente vista como necessária para um ensaio a ser justificado éticamente, especialmente quando os participantes são atribuidos a braços de ensaios randômicamente. Esta visão é sustentada pela idéia de que os pesquisadores médicos não podem

equipoise

A state of genuine uncertainty as to which of two or more alternative health technologies should constitute the standard of care, considering benefits and harms. Note 1: Equipoise is widely viewed as necessary for a trial to be ethically justified, especially when participants are assigned to trial arms randomly. This view is underpinned by the idea that physician-researchers cannot waive their duty to act in the best interests of their patients. Note 2: As specified by the definition, equipoise usually refers to the state of uncertainty inhabited by the scientific and medical community, but it can also refer to the state of uncertainty inhabited by the researchers conducting a trial or by patients and the public.

\begin{center}\rule{0.5\linewidth}{0.5pt}\end{center}

equivalente em anos de saúde

O número hipotético de anos em saúde perfeita que podem ser considerados equivalentes ao número real de anos num estado de saúde imperfeito.

healthy-years equivalent (HYE)

The hypothetical number of years spent in perfect health which could be considered equivalent to the actual number of years spent in a defined imperfect state of health.

\begin{center}\rule{0.5\linewidth}{0.5pt}\end{center}

erro falso negativo

Em teste de hipóteses, um erro que ocorre quando se conclui incorretamente que a hipótese nula é verdadeira. Nota. Por exemplo, um erro do Tipo II é cometido se não se detectar diferença entre os desfechos de um grupo experimental e os do grupo de controle, quando de fato uma diferença existe. A probabilidade de se cometer esse tipo de erro é designada como beta (β). Sinônimo: erro do tipo II.

false negative error

Em teste de hipóteses, um erro que ocorre quando se conclui incorretamente que a hipótese nula é verdadeira. Nota. Por exemplo, um erro do Tipo II é cometido se não se detectar diferença entre os desfechos de um grupo experimental e os do grupo de controle, quando de fato uma diferença existe. A probabilidade de se cometer esse tipo de erro é designada como beta (β). Sinônimo: erro do tipo II.

\begin{center}\rule{0.5\linewidth}{0.5pt}\end{center}

erro falso positivo

Em teste de hipóteses, um erro que ocorre quando se conclui incorretamente que a hipótese nula é falsa. Nota. Por exemplo, um erro do Tipo I é cometido se for detectada uma diferença entre os desfechos de um grupo experimental e os do grupo de controle, quando de fato a diferença não existe. A probabilidade de se cometer esse tipo de erro é designada como alfa (α). Sinônimo: erro do tipo I.

false positive error

In hypothesis testing, an error that occurs when it is incorrectly concluded that a null hypothesis is false. Note: For example, a Type I error is made if a difference is detected between the outcomes in an experimental group and those in a control group when in fact such a difference does not exist. The probability of making this type of error is designed as alpha (α). Syn.: Type I error.

\begin{center}\rule{0.5\linewidth}{0.5pt}\end{center}

erro na amostragem

Veja erro randômico.

sampling error

See random error.

\begin{center}\rule{0.5\linewidth}{0.5pt}\end{center}

erro randômico

Um desvio devido ao acaso entre a estimativa de um ponto e o valor do parâmetro estimado. Nota: erro randômico conduz a aproximadamente o mesmo número de resultados maiores do que o valor verdadeiro.

random error

A numerical value obtained in a sample and considered the best estimate of the population from which the sample is taken.

\begin{center}\rule{0.5\linewidth}{0.5pt}\end{center}

erro sistemático

Veja viés.

systematic error

See bias.

\begin{center}\rule{0.5\linewidth}{0.5pt}\end{center}

erro Tipo I

Veja alfa.

Type I error

See alpha.

\begin{center}\rule{0.5\linewidth}{0.5pt}\end{center}

erro Tipo II

Veja beta.

Type II error

See beta.

\begin{center}\rule{0.5\linewidth}{0.5pt}\end{center}

escolha informada

Uma escolha feita por pessoa competente com base na informação compreensiva relativa a uma intervenção, procedimento ou disponibilidade e qualidade de serviços de saúde.

informed choice

A choice made by a competent person on the basis of comprehensive information regarding an intervention, procedure, or the availability and quality of health services.

\begin{center}\rule{0.5\linewidth}{0.5pt}\end{center}

escolha pela chance

Um método de medida direta da utilidade de um estado da saúde, encontrado nos fundamentos do axioma de Neumann-Morgensten sobre a teoria da utilidade esperada. Nota: Um escore de utilidade é obtido estimando-se as probabilidades de que uma pessoa que foi solicitada a comparar duas opções, uma incerta (a de risco), e outra certa, não preferirá uma sobre a outra.

standard gamble (SG)

A method of directly measuring the utility of a health status, founded on the fundamental von Neumann-Morgenstern axioms of expected utility theory. Note: A utility score is obtained by estimating the probabilities that a person who is asked to compare two options, one uncertain (the gamble), the other certain, will not prefer one over the over.

\begin{center}\rule{0.5\linewidth}{0.5pt}\end{center}

escore da qualidade

Um valor atribuido à validade de um estudo, baseado tanto numa avaliação geral do mesmo ou em um ou mais critérios específicos (por exemplo, a ocultação da alocação). Nota: o escore da qualidade pode ser expresso em letras (A,B,C) ou números.

quality score

Um valor atribuido à validade de um estudo, baseado tanto numa avaliação geral do mesmo ou em um ou mais critérios específicos (por exemplo, a ocultação da alocação). Nota: o escore da qualidade pode ser expresso em letras (A,B,C) ou números.

\begin{center}\rule{0.5\linewidth}{0.5pt}\end{center}

especificidade

Uma característica de performance de um teste diagnóstico, definido como a proporção de pessoas que tiveram um resultado negativo no teste e as pessoas que não tiveram a doença; é calculada assim: verdadeiros negativos ÷ (verdadeiros negativos + falsos positivos). Conceitos relacionados: acurácia, curva ROC e sensibilidade.

specificity

A characteristic of the performance of a diagnostic test, defined as the proportion of persons with a negative test result among persons who do not have the disease; it is calculated as follows: true negatives ÷ (true negatives + false positives). Related concepts: accuracy, ROC curve and sensitivity

\begin{center}\rule{0.5\linewidth}{0.5pt}\end{center}

estabelecimento de prioridade

A atribuição de uma ordem de prioridade baseada num critério explícito ou implícito para a seleção de tecnologias em saúde para avaliação.

priority setting

The assignment of an order of priority based on explicit or implicit criteria for selection of health technologies for assessment.

\begin{center}\rule{0.5\linewidth}{0.5pt}\end{center}

estado de saúde

O nível de saúde de uma pessoa, grupo ou população conforme a avaliação da própria pessoa ou através de medidas objetivas.

health status

The level of health of a person, group or population, as assessed by the person himself/herself or through objective measures.

\begin{center}\rule{0.5\linewidth}{0.5pt}\end{center}

estagiamento

A classificação da severidade de uma doença em estágios com base nos sinais, sintomas e testes paraclínicos.

staging

The classification of the severity of a disease in stages on the basis of the signs and symptoms and the paraclinical tests.

\begin{center}\rule{0.5\linewidth}{0.5pt}\end{center}

Estatística Kappa

Uma medida do grau de acordo entre duas medidas da mesma variável categórica sobre e acima do acordo que foi devido apenas à probabilidade. Nota: o acordo entre dois avaliadores ou entre duas medidas de tempo para um avaliador podem ser medidas.

Kappa statistic

A measure of the degree of agreement between two measures of the same categorical variable over and above the agreement that is due to chance alone. Note: The agreement between two raters or between two measurement times for one rater can be measured.

\begin{center}\rule{0.5\linewidth}{0.5pt}\end{center}

estimativa do efeito

Num estudo, a relação entre o desfecho observado quando a intervenção é aplicada e o desfecho observado num grupo de controle. Nota: esta relação pode ser expressa como um número que precisa ser tratado, relação de probabilidade, diferença de risco, risco relativo, média padronizada ou diferença ponderada da média ponderada. (conceito relacionado: efeito do tratamento).

estimate of effect

In a study, the relationship between the outcome observed when an intervention is applied and the outcome observed in a control group. Note: This relationship may be expressed as a number needed to treat, odds ratio, risk difference, relative risk, standardised mean difference or weighted mean difference. (Related concept: treatment effect)

\begin{center}\rule{0.5\linewidth}{0.5pt}\end{center}

estimativa pontual

Um valor numérico obtido numa amostra e considerado a melhor estimativa da população da qual a amostra foi tomada. Nota: por exemplo, x e s são estimativas pontuais de µ e σ, que representam a média e o desvio padrão. Um intervalo de confiança é tipicamente construído em torno da estimativa pontual. (conceito relacionado: intervalo de confiança).

point estimate

A numerical value obtained in a sample and considered the best estimate of the population from which the sample is taken. Note: For example, x and s are point estimates of µ and σ, which represent the mean and the standard deviation. A confidence interval is typically constructed around the point estimate. (Related concept: confidence interval)

\begin{center}\rule{0.5\linewidth}{0.5pt}\end{center}

estratégia de busca

A combinação de fontes, termos e limites usados numa pesquisa de literatura para identificar informação pertinente para revisão sistemática ou avaliação de tecnologia em saúde.

search strategy

The combination of sources, terms and limits used in a literature search to identify the pertinent information for a systematic review or health technology assessment.

estudo antes e depois

projeto de pesquisa onde um grupo de elementos é observado antes e depois da intervenção ou exposição a um dado fator.

before-and-after study

A research design where a group of subjects is observed before and after an intervention or exposure to a given factor.

\begin{center}\rule{0.5\linewidth}{0.5pt}\end{center}

estudo cross-seccional

Estudo observacional no qual os investigadores escolhem um grupo de elementos de uma dada população e medem presença simultânea ou ausência do fator de risco e do desfecho de interesse. Nota: a seleção do grupo de elementos é algumas vezes feita através de sorteio de amostragem aleatória.

cross-sectional study

An observational study in which the investigators choose a group of subjects in a given population and measure the simultaneous presence or absence of the risk factor and the outcome of interest. Note: The selection of the group of subjects is sometimes done through random sampling.

\begin{center}\rule{0.5\linewidth}{0.5pt}\end{center}

estudo de braço único

Numa análise ou avaliação de um estudo com apenas uma ramificação, isto é, um ensaio no qual não houve grupo para comparação paralela e todos os indivíduos receberam a mesma intervenção.

single-arm study

An analysis or evaluation of a study with only one branch, i.e.~a trial in which there was no parallel comparison group and all the subjects received the same intervention.

\begin{center}\rule{0.5\linewidth}{0.5pt}\end{center}

estudo de caso-controle

estudo observacional retrospectivo planejado para determinar se existe uma relação entre um dado parâmetro (Por xemplo: uma doença ou um determinante para a saúde) e uma causa especificada. Nota: investigadores selecionam um grupo de elementos apresentando o parâmetro de interesse (casos) e um grupo não apresentando o parâmetro (controle). A seguir eles comparam os históricos dos casos e controles para determinar se existe uma associação entre o parâmetro e a causa sendo estudada. Nota: o termo ``estudo de caso-controle'' é frequentemente considerado, incorretamente, como sinônimo de ``estudo retrospectivo''. Embora, em ambos os casos, seja necessário olhar para o passado pela cadeia de eventos, o termo ``estudo retrospectivo'' não necessariamente inclui o elemento de comparação, que caracteriza os estudos de caso-controle. O qualificador ``retrospectivo'' simplesmente indica a direção temporal do estudo.

case-control study

A retrospective observational study designed to determine whether there is a link between a given parameter (e.g.~a disease or a health determinant) and a specified cause. Note: Investigators select a group of subjects presenting the parameter of interest (cases) and a group not presenting the parameter (controls). Then they compare the histories of the cases and controls to determine whether there is an association between the parameter and the cause being studied. Note: The term ``case-control study'' is often, incorrectly, considered a synonym of ``retrospective study.'' Although, in both cases, it is necessary to look into the past for the chain of events, the term ``retrospective study'' does not necessarily include the element of comparison of two or more groups defined at the start of the study, an element that characterizes case-control studies. The qualifier ``retrospective'' simply indicates the temporal direction of the study.

\begin{center}\rule{0.5\linewidth}{0.5pt}\end{center}

estudo de coorte

estudo observacional no qual os parâmetros de um grupo de elementos expostos a um dado fator são comparados com aqueles de um grupo similar de elementos não expostos a esse fator. Nota: para que os grupos sejam comparáveis, os investigadores podem parear elementos exposto a um dado fator com outros não expostos ao fator com base em características comuns (Por exemplo: idade, sexo, severidade da doença) ou fazer os necessários ajustes estatísticos para obter, nos elementos pareados, características que sejam as mais similares possíveis. Outra definição: um estudo observacional no qual os investigadores selecionam e monitoram um grupo (coorte) de elementos expostos a um dado fator, para verificar os efeitos da exposição ao fator na saúde dos membros do coorte. Nota: este tipo de estudo é usualmente prospectivo, mas pode ser retrospectivo quando os dados já tiverem sido coletados. O importante é o ponto inicial: identificar os elementos expostos e não expostos a um dado fator e então determinar os efeitos da exposição na saúde.

cohort study

An observational study in which the parameters of a group of subjects exposed to a given factor are compared with those of a similar group of subjects not exposed to the factor. Note: To make the groups comparable, the investigators may match subjects exposed to a given factor with subjects not exposed to the factor on the basis of common characteristics (e.g.~age, gender and disease severity) or may make the necessary statistical adjustments to obtain, in the matched subjects, characteristics that are as similar as possible. Other definition: An observational study in which the investigators select and monitor a group (cohort) of subjects exposed and non exposed to a given factor, to verify the effects of the exposure to the factor on the health of the cohort's members. Note: This type of study is usually prospective, but may also be retrospective when the data have already been gathered. What is important is the starting point: identify the subjects exposed and not exposed to a given factor and then determine the health effects of the exposure.

\begin{center}\rule{0.5\linewidth}{0.5pt}\end{center}

estudo de intervenção

Veja ensaio clínico.

intervention study

See clinical trial.

\begin{center}\rule{0.5\linewidth}{0.5pt}\end{center}

estudo de prevalência

veja estudo transversal

prevalence study

See cross-sectional study.

\begin{center}\rule{0.5\linewidth}{0.5pt}\end{center}

estudo integrativo

Veja pesquisa secundária.

integrative study

See secondary research.

\begin{center}\rule{0.5\linewidth}{0.5pt}\end{center}

estudo não experimental

Um estudo baseado nas condições existentes de exposição a um dado fator, sem a intervenção do investigador. Nota: neste tipo de estudo, comumente usado, a evidência de uma relação causal entre a intervenção e o(s) desfecho(s) é relativamente fraca. (conceito relacionado: estudo observacional).

non-experimental study

A study based on the existing conditions of exposure to a given factor, without investigator intervention. Note: In this commonly used type of study, evidence of a causal link between the intervention and the outcome(s) of interest is relatively weak. (Related concept: observational study).

\begin{center}\rule{0.5\linewidth}{0.5pt}\end{center}

estudo negativo

Um estudo que não apresenta achados estatisticamente significativos, indicando que a intervenção sendo estudada provavelmente não tem um efeito com benefícios. Nota: este termo pode gerar confusão porque se refere tanto a significância etatística quanto à direção do efeito (positivo ou negativo). Como estudos analisam desfechos múltiplos, simultâneamente, o critério para a classificação de um estudo como ``negativo'' nem sempre é claro e, no caso de estudos de riscos ou de efeitos adversos, os estudos ``negativos'' são aqueles que não apresentam um efeito prejudicial. Antônimo: estudo positivo.

negative study

A study that does not have statistically significant findings, indicating that the intervention being studied probably does not have a beneficial effect. Note: This term can generate confusion because it refers both to statistical significance and the direction of effect (positive or negative). Since studies analyse multiple outcomes simultaneously, the criteria for classifying a study as ``negative'' are not always clear and, in the case of studies of risk or adverse effects, ``negative'' studies are ones that do not show a harmful effect. Ant.: positive study.

\begin{center}\rule{0.5\linewidth}{0.5pt}\end{center}

estudo observacional

Um estudo no qual os investigadores não intervém, somente observam os indivíduos que estão (e algumas vezes não estão, para a finalidade de comparação) expostos a um dado fator, e interpretam os desfechos. Nota: este tipo de estudo está mais sujeito a viéses do que um estudo experimental, tal como um ensaio controlado e randomizado. (conceito relacionado: estudo não experimental).

observational study

A study in which the investigators do not intervene, but only observe subjects who are (and sometimes who are not, for comparison purposes) exposed to a given factor, and interpret the outcomes. Note: This type of study is more subject to bias than is an experimental study such as a randomised controlled trial. (Related concept: non-experimental study.)

\begin{center}\rule{0.5\linewidth}{0.5pt}\end{center}

estudo positivo

Um estudo no qual foram obtidos resultados significantes, indicando que a intervenção sendo estudada provavelmente tem um efeito benéfico. Nota: este termo pode gerar confusão porque se refere tanto para a significância estatística quanto para a direção ou efeito. Como os estudos tem muitos indicadores de desfecho, o critério definindo um estudo ``positivo'' nem sempre fica claro e, no caso de estudos de um risco ou de efeitos adversos, por exemplo, estudos positivos são os que mostram um efeito danoso. Antônimo: estudo negativo.

positive study

A study that has obtained statistically significant results, indicating that the intervention being studied probably has a beneficial effect. Note: This term can generate confusion because it refers both to statistical significance and the direction of effect. Since studies often have many outcome indicators, the criteria defining a ``positive'' study are not always clear and, in the case of studies on a risk or given adverse effects, for example, positive studies are those that show a harmful effect. Ant.: negative study

\begin{center}\rule{0.5\linewidth}{0.5pt}\end{center}

estudo pré-teste/pós-teste

veja estudo pré e pós

pre-test/post-test study

See before-and-after study.

\begin{center}\rule{0.5\linewidth}{0.5pt}\end{center}

estudo prospectivo

Um estudo para avaliar os efeitos da exposição a uma dada intervenção ou fator, no qual os indivíduos são divididos em grupos que são expostos ou não expostos à intervenção ou fator de interesse, antes que os desfechos tenham ocorrido. Nota: ensaios randomizados controlados são sempre prospectivos, mas os estudos de caso-controle nunca são. Estudo de cohorte podem ser prospectivos ou retrospectivos. Conceito relacionado estudo retrospectivo. † Um estudo de cohorte pode ser um estudo prospectivo, mas um estudo prospectivo não necessariamente é um estudo de cohorte: um estudo de coorte envolve uma metodologia, isto é, uma comparação que não é necessariamente comum a todos os estudos prospectivos. O qualificador ``prospectivo'', sem maiores detalhes, indica apenas a direção temporal do estudo.

prospective study

A study to evaluate the effects of exposure to a given intervention or factor, in which the subjects are divided into groups that are exposed or not exposed to the intervention or factor of interest before the outcomes have occurred. Note: Randomized controlled trials are always prospective, but case-control studies never are. Cohort studies may be prospective or retrospective. Related concept: cohort study; Ant.: retrospective study. † A cohort study may be a prospective study, but a prospective study is not necessarily a cohort study : a cohort study involves a methodology, i.e.~comparison, that is not necessarily common to all prospective studies. The qualifier ``prospective,'' without further details, indicates only the study's temporal direction.

\begin{center}\rule{0.5\linewidth}{0.5pt}\end{center}

estudo retrospectivo

Um estudo no qual os investigadores do desfecho pós-fato num grupo de indivíduos selecionados com base na sua exposição a uma dada intervenção ou fator. Nota: este tipo de estudo está mais sujeito a viés do que um estudo prospectivo. Estudos de caso-controle são sempre retrospectivos, estudos de coorte às vezes são e ensaios randomizados controlados nunca são. Antônimo: estudo prospectivo.

retrospective study

A study in which the investigators do ex-post analysis of the outcomes in a group of subjects selected on the basis of their exposure to a given intervention or factor. Note: This type of study is more subject to bias than are prospective studies. Case-control studies are always retrospective, cohort studies sometimes are, and randomized controlled trials never are. Ant.: prospective study.

\begin{center}\rule{0.5\linewidth}{0.5pt}\end{center}

estudo sintético

Veja estudo secundário.

synthetic study

See secondary research.

\begin{center}\rule{0.5\linewidth}{0.5pt}\end{center}

estudos fase I, II II ou IV

Um sistema de codificação para as várias fases de ensaios clínicos de novas tecnologias. Nota: o sistema se aplica principalmente a drogas e dispositivos. Esses ensaios são parte do desenvolvimento e do processo de aprovação exigido por agências regulatórias. Ensaio Fase I: um ensaio envolvendo usualmente 20 a 80 voluntários para determinar a segurança de uma droga, a faixa segura de dosagem, absorção, atividade metabólica, excreção e duração da atividade. Ensaio Fase II: um ensaio controlado envolvendo 100 a 300 indivíduos para determinar a eficacia da droga e reações adversas à mesma (estes ensaios são às vezes divididos em ensaios piloto Fase IIa e ensaios bem controlados Fase IIb). Ensaio Fase III: um ensaio controlado grande, para determinar eficácia da droga e monitorar eventos adversos durante o uso de longo prazo (estes ensaios são às vezes divididos em ensaios Fase IIIa, conduzidos antes da aplicação autorizada por marketing, e ensaios Fase IIIb, conduidos após a aplicação autorizada por marketing, porém antes de sua aprovação). Ensaio Fase IV: um estudo de farmaco-vigilância, pós-marketing, para monitorar os efeitos de longo termo da droa e prover informação adicional sobre a segurança e eficácia da droga, incluindo grupos de pacientes com diferentes regimes de dosagem.

phase I, II, III or IV study

A codification system for the various phases of clinical trials of new technologies. Note: The system applies mainly to drugs and devices. These trials are part of the development and approval process required by regulatory agencies. Phase I trial: A trial usually involving 20 to 80 healthy volunteers to determine a drug's safety, safe dosage range, absorption, metabolic activity, excretion and duration of activity. Phase II trial:
A controlled trial involving 100 to 300 subjects to determine a drug's efficacy and adverse reactions to the drug (these trials are sometimes divided into Phase IIa pilot trials and Phase IIb well-controlled trials). Phase III trial:
A large controlled trial to determine a drug's efficacy and monitor adverse events during longer-term use (these trials are sometimes divided into Phase IIIa trials, conducted before the marketing authorization application, and Phase IIIb trials, conducted after the marketing authorization application, but before approval). Phase IV trial:
A post-marketing pharmacovigilance study to monitor a drug's long-term effects and provide additional information on the drug's safety and efficacy, including for patient groups with different dosing regimens.

\begin{center}\rule{0.5\linewidth}{0.5pt}\end{center}

ética

Uma teoria pretendendo determinar o que é bom, tanto para o indivíduo como para a sociedade como um todo. Outra definição: uma reflexão crítica sobre os desejos e fundamentos da ação humana (incluindo valores expressos) e um guia da conduta humana através do questionamento, discussão e a proposição de regras, ao invés da imposição de regras. Nota: à partir desta perspectiva, a ética se difere da moralidade, a qual é um sistema e a codificação de requisitos (padrão ou regras) governando a conduta humana, sendo os requisitos caracterizados pela sua universalidade e por um efeito restritivo. No entanto, a ética, moralidade e lei são proximamente interdependentes.

ethics

A theory aimed at determining what is good, both for the individual and for society as a whole. Other definition: A critical reflection on the aims and foundations of human action (including expressed values) and a guiding of human conduct through questioning, discussion and the proposing of guidelines, rather than the imposing of rules. Note: From this perspective, ethics differ from morality, which is a system and codification of requirements (standards and rules) governing human conduct, requirements characterised by their universality and by a restraining effect. Nonetheless, ethics, morality and law are closely interdependent.

\begin{center}\rule{0.5\linewidth}{0.5pt}\end{center}

EuroScan

A Rede Européia de Informação sobre Mudanças em Tecnologia da Saúde: \url{http://www.euroscan.bham.ac.uk/}. Uma rede de agências em avaliação de tecnologia da saúde que identifica e avalia tecnologias novas e emergentes.

EuroScan

The European Information Network on New and Changing Health Technologies \url{http://www.euroscan.bham.ac.uk/}. A network of health technology assessment agencies that identify and assess new and emerging technologies.

\begin{center}\rule{0.5\linewidth}{0.5pt}\end{center}

evento adverso

qualquer alteração nociva, patológica ou não intencional numa função físca ou metabólica, revelada por sinais ou sintomas ou alteração de resultados de testes de laboratório, em qualquer fase de um estudo clínico, se ou não a alteração for considerada relacionada ao tratamento. Nota: ela pode envolver a exacerbação de uma condição pré-existente, distúrbios intercorrentes, um acidente, uma interação medicamentosa, ou ou significativa piora da doença.

adverse event

Any noxious, pathological or unintended change in a physical or metabolic function, revealed by signs or symptoms or a change in the results of laboratory tests, in any phase of a clinical study, whether or not the change is considered treatment-related. Note: It may involve the exacerbation of a pre-existing condition, intercurrent diseases, an accident, a drug interaction or a significant worsening of the disease.

\begin{center}\rule{0.5\linewidth}{0.5pt}\end{center}

exceção para dispositivo sob análise

Um processo regulatório e categoriação através do qual o FDA (Food and Drug Administration - Administração de Alimentos e Drogas, orgão norte-americano) autoriza o uso especificado de um dispositivo não autorizado, em configurações controladas, para a coleta de dados sobre a segurança do dispositivo e sobre a efetividade teórica e prática. Nota: os dados assim coletados podem ser utilizados em aplicações subsequentes e aprovadas de pré-marketing.

investigational device exemption (IDE)

A regulatory process and categorization whereby the United States Food and Drug Administration (FDA) authorizes specified use of an unapproved device in controlled settings for collecting data on the device's safety and theoretical or practical effectiveness. Note: The data thus collected may be used subsequently in a premarketing approval application.

\begin{center}\rule{0.5\linewidth}{0.5pt}\end{center}

exploração do horizonte

A identificação sistemática de tecnologias na saúde que são novas, emergentes ou tornando-se obsoletas, e que tem potencial de afetar a saúde, serviços de saúde e/ou a sociedade. Nota: 1: termos relacionados incluem aviso precoce e sistema de alerta.

horizon scanning

The systematic identification of health technologies that are new, emerging or becoming obsolete and that have the potential to effect health, health services and/or society. Note 1: Related terms include early awareness and alert system

\begin{center}\rule{0.5\linewidth}{0.5pt}\end{center}

extrapolação

Um cálculo, para valores de uma variável caindo fora das séries de variáveis observadas, os valores de uma função empiricamente conhecida e, por extensão, a predição resultante, dedução, previsão, transposição ou projeção.

extrapolation

A calculation, for values of the variable lying outside the series of observed variables, the values of an empirically known function and, by extension, the resulting prediction, deduction, forecasting, transposition or projection.

\hypertarget{f}{%
\chapter*{F}\label{f}}
\addcontentsline{toc}{chapter}{F}

Farmacoeconomia

Uma disciplina para avaliação econômica de produtos e serviços farmacêuticos através da determinação, medida e comparação de seus custos e desfechos. Redação alternativa: Fármaco-economia

pharmacoeconomics

A discipline for economic evaluation of pharmaceutical products and services through determination, measurement and comparison of their costs and outcomes. Alternate spelling: pharmaco-economics

\begin{center}\rule{0.5\linewidth}{0.5pt}\end{center}

fator de risco

Um fator associado com a ocorrência de um problema de saúde. Nota: pode ser um comportamento pessoal, estilo de vida, exposição a perigos ambientais, características congênitas ou hereditárias, etc. Por exemplo, fumar é um fator de risco para câncer no pulmão.

risk factor

A factor associated with the occurrence of a health problem. Note: It may be personal behaviour, lifestyle, exposure to environmental hazards, congenital or hereditary characteristics, etc. For example, smoking is a risk factor for lung cancer.

\begin{center}\rule{0.5\linewidth}{0.5pt}\end{center}

folheto para informação do paciente

Um documento escrito para pacientes e para o público em geral que apresenta com clareza informação-chave sobre um tratamento ou procedimento, incluindo um explicação clara de termos técnicos.

patient information leaflet

A written document for patients and the general public that clearly presents the key information about a treatment or procedure, including a clear explanation of any technical terms.

\begin{center}\rule{0.5\linewidth}{0.5pt}\end{center}

formulário

Uma nomenclatura das drogas cobertas no todo ou em parte num plano de seguro.

formulary

A nomenclature of the drugs covered in whole or in part under an insurance plan.

\begin{center}\rule{0.5\linewidth}{0.5pt}\end{center}

fronteira eficiente

Uma curva formada pelos incrementos das razões de custo-efetividade ou custo-utilidade numa representação gráfica de comparadores não dominantes.

efficient frontier

A curve formed by the incremental cost-effectiveness or cost-utility ratios in a graphical representation of the non-dominated comparators.

\begin{center}\rule{0.5\linewidth}{0.5pt}\end{center}

garantia da qualidade

Um método combinando todas as ações pré-determinadas e sistemáticas necessárias para prover uma confiança apropriada de que o produto ou serviço atendem aos requisitos da qualidade. Nota: a garantia da qualidade inclui a definição e aplicação de padrões de cuidados com a saúde baseada no melhor conhecimento disponível, e isso envolve a avaliação da qualidade e atividades para corrigir e reduzir variações nas práticas de cuidados com a saúde relativas a esses padrões, ou para melhorar tais práticas.

quality assurance

A method combining all the pre-determined and systematic actions needed to provide appropriate confidence that a product or service will meet given quality requirements. Note: Quality assurance includes setting and applying health care standards based on the best available knowledge, and this involves quality assessment and activities to correct and reduce variations in health care practices relative to these standards, or to improve those practices.

\begin{center}\rule{0.5\linewidth}{0.5pt}\end{center}

generalizabilidade

Veja validade externa. Outra forma de escrita: generalizabilidade.

generalizability

See external validity. Alternate spelling: generalisability.

\begin{center}\rule{0.5\linewidth}{0.5pt}\end{center}

gerenciamento de doenças

\begin{enumerate}
\def\labelenumi{\arabic{enumi})}
\tightlist
\item
  Um processo sistemático de gerenciamento do cuidado de pacientes com doenças específicas ou condições (particularmente condições crônicas), incluindo serviços ambulatoriais, de internamento e auxiliares, com o objetivo de, entre outros, reduzir episódios agudos, hospitalizações e variações no cuidado da saúde e em práticas de serviços, melhorando o desfecho em saúde e reduzindo custos. Nota: o gerenciamento de doenças pode envolver a melhoria contínua da qualidade e outros paradigmas de gerenciamento. Pode envolver um processo cíclico ou seguir protocolos da prática profissional, medindo o resultado dos desfechos, retroalimentando esses resultados para os médicos e revisando protocolos se adequado. 2) Todos os cuidados e serviços proporcionados a um paciente.
\end{enumerate}

disease management

\begin{enumerate}
\def\labelenumi{\arabic{enumi})}
\tightlist
\item
  A systematic process of managing care of patients with specific diseases or conditions (particularly chronic conditions), including outpatient, inpatient, and ancillary services, and aimed at, among other things, reducing acute episodes, hospitalizations and variations in health care and services practices, improving health outcomes, and reducing costs. Note: Disease management may involve continuous quality improvement or other management paradigms. It may involve a cyclical process of following professional practice protocols, measuring the resulting outcomes, feeding those results back to clinicians, and revising protocols as appropriate. 2) All the care and services with which a patient is provided.
\end{enumerate}

\begin{center}\rule{0.5\linewidth}{0.5pt}\end{center}

gráfico de funil

Um diagrama de dispersão no qual cada ponto representa um estudo, relatando o efeito estimado e o tamanho da amostra (ou outro indicador de precisão do estudo). Nota: na ausência de viés de publicação, o diagrama de dispersão formará um funil invertido. Um buraco no canto esquerdo do funil indicará a presença de um viés de publicação.

funnel plot

A scatter diagram in which each point represents a study, relating the estimate of effect and the sample size (or another indicator of the precision of the study). Note: In the absence of publication bias, the scatter diagram will form an inverted funnel. A hole in the lower left corner of the funnel indicates the presence of publication bias.

\hypertarget{g}{%
\chapter*{G}\label{g}}
\addcontentsline{toc}{chapter}{G}

grupo de controle

Um grupo de de pessoas que serve como base de comparação para se estabelecer os efeitos da intervenção aplicada aos elementos do grupo de tratamento. Nota: dependendo das circunstâncias do ensaio, o grupo de controle pode não receber nenhum tratamento, ou receber o tratamento padrão, ou receber um placebo. Para tornar a comparação válida, a composição do grupo de controle deve se assemelhar tanto quanto possível ao grupo de tratamento.

control group

A group of persons that serves as the basis of comparison for assessing the effects of the intervention applied to the subjects in the treatment group. Note: Depending on the circumstances of the trial, a control group may receive no treatment, the usual or standard treatment, or a placebo. To make the comparison valid, the composition of the control group must resemble that of the treatment group as closely as possible.

\begin{center}\rule{0.5\linewidth}{0.5pt}\end{center}

grupo de enfoque

Um processo para colher informação através de entrevistas em grupo, com foco na comunicação e interações entre os participantes. Nota: os participantes constituem um grupo relativamente homogêneo (8 a 12 pessoas) solicitados a comentar sobre certas questões propostas pelo moderador. Eles são encorajados a comentar os pontos de vista dos demais elementos do grupo, ao invés de esperar que o moderador pergunte. Este processo ajuda a explorar e esclarecer pontos de vista, que poderiam ser mais difíceis de obter numa entrevista cara a cara.

focus group

A process for gathering information through group interviews, with a focus on communication and interactions among the participants. Note: The participants are a relatively homogeneous group (8 to 12 people) asked to comment on certain questions posed by the moderator. They are encouraged to comment on each other's points of view, rather than to wait for the moderator to ask them to answer questions in turn. This process helps explore and clarify points of view, which would be more difficult to do in one-on-one meetings.

\begin{center}\rule{0.5\linewidth}{0.5pt}\end{center}

guia

Veja guia clínica.

guideline

See clinical guideline.

\hypertarget{h}{%
\chapter*{H}\label{h}}
\addcontentsline{toc}{chapter}{H}

hermenêutica

A interpretação de um todo em relação a suas partes e vice-versa. Nota: originalmente a interpretação de textos bíblicos.

hermeneutics

The interpretation of a whole in relation to its parts, and vice versa. Note: Originally, the interpretation of biblical texts.

\begin{center}\rule{0.5\linewidth}{0.5pt}\end{center}

heterogeneidade

Numa revisão sistemática corresponde à variabilidade ou às diferenças presentes nos estudos selecionados. Nota: às vezes uma distinção é feita entre ``heterogeneidade estatística'' ( diferenças nos efeitos reportados) e ``heterogeneidade metodológica'' (diferenças no projeto do estudo em relação às características chave dos indivíduos, intervenções ou dos critéros de avaliação de desfechos). Testes estatísticos de heterogeneidade são usados para se determinar se a variabilidade observada nos resultados do estudo é maior do que a variabilidae que pode ser esperada ao acaso. No entanto, esses testes tem baixo poder estatístico. Antônimo: homogeneidade

heterogeneity

In a systematic review, the variability of or differences in the selected studies.

Note: A distinction is sometimes made between ``statistical heterogeneity'' (differences in the reported effects) and ``methodological heterogeneity'' (differences in study design with regard to the key characteristics of the subjects, interventions or outcome evaluation criteria). Statistical tests of heterogeneity are used to determine whether the observed variability in study results effect size is greater than the variability that can be expected to occur by chance. However, these tests have low statistical power. Antonym: homogeneity.

\begin{center}\rule{0.5\linewidth}{0.5pt}\end{center}

hierarquia de evidência

Uma classificação de estudos com base na sua validade. Nota: esse método é usado para se determinar o peso que deveria ser atribuido a um estudo. Várias hierarquias de evidência são usadas na avaliação de tecnologia em saúde.

hierarchy of evidence

A classification of studies on the basis of their validity. Note: This method is used to determine the weight that should be given to a study. Various hierarchies of evidence are used in health technology assessment.

\begin{center}\rule{0.5\linewidth}{0.5pt}\end{center}

hipótese nula

No teste de hipóteses, uma proposta de que não há associação entre duas variáveis ou que não existe diferença entre dois valores. Nota: por exemplo, a hipótese nula pode indicar que uma intervenção não tem efeito, isto é, que não há uma diferença verdadeira entre os resultados obtidos no grupo experimental e no grupo de controle.

null hypothesis

In hypothesis testing, a proposal that there is no association between two variables or no difference between two values. Note: For example, the null hypothesis could indicate that an intervention has no effect, i.e.~that there is no true difference between the results obtained in the experimental group and the control group.

\begin{center}\rule{0.5\linewidth}{0.5pt}\end{center}

história natural

O curso de uma doença desde seu começo até sua resolução. Nota: muitas doenças possuem estágios bem definidos, tais como o início patológico e os estágios pré-sintomáticos e de manifestação clínica da doença.

natural history

The course of a disease from onset to resolution. Note: Many diseases have well-defined stages, such as the pathological onset, presymptomatic and clinically manifest disease stages.

\begin{center}\rule{0.5\linewidth}{0.5pt}\end{center}

homogeneidade

A similaridade de estudos selecionados numa revisão sistemática. Nota: estudos são considerados ``estatisticamente homogêneos'' se seus resultados não variarem mais do que seria esperado pelo acaso. Numa revisão sistemática, a ``homogeneidade clínica'' significa que os indivíduos, intervenções e instrumentos de medição são similares ou comparáveis. Antônimo: heterogeneidade.

Homogeneity

In a systematic review, the similarity of the studies selected. Note: Studies are considered ``statistically homogeneous'' if their results vary no more than they might be expected to by chance. In a systematic review, ``clinical homogeneity'' means that the subjects, interventions and measuring instruments are similar or comparable. Antonym: heterogeneity.

\hypertarget{i}{%
\chapter*{I}\label{i}}
\addcontentsline{toc}{chapter}{I}

impacto diagnóstico

Uma medida de utilidade de resultados de testes na elaboração de um diagnóstico. Nota: este conceito é usado relativamente pouco.

diagnostic impact

A measurement of the usefulness of test results for making a diagnosis. Note: This concept is used relatively little.

\begin{center}\rule{0.5\linewidth}{0.5pt}\end{center}

impacto terapêutico

Uma característica de eficácia do teste diagnóstico que descreve o efeito de um teste diagnóstico em opções de tratamento.

therapeutic impact

A characteristic of diagnostic test efficacy that describes the effect of a diagnostic test on treatment choices.

\begin{center}\rule{0.5\linewidth}{0.5pt}\end{center}

imperativo tecnológico

A inclinação em usar uma tecnologia que tenha potencial para proporcionar algum benefício, seja marginal ou infundado, baseada numa permanente fascinação pela tecnologia, pela expectativa de que o novo é melhor, e por outros incentivos financeiros ou profissionais.

technological imperative

The inclination to use a technology that has potential for some benefit, however marginal or unsubstantiated, based on an abiding fascination with technology, the expectation that new is better, and financial and other professional incentives.

\begin{center}\rule{0.5\linewidth}{0.5pt}\end{center}

INAHTA

International Network of Agencies for Health Technology (INAHTA); Rede Internacional de Agências para Avaliação de Tecnologias em Saúde; Uma associação internacional sem fins lucrativos de agências de avaliação de tecnologia em saúde, representando mais de 40 agências no mundo todo. (\url{http://www.inahta.net}).

INAHTA

International Network of Agencies for Health Technology (INAHTA); An international non-profit association of health technology assessment agencies, representing over 40 agencies worldwide (\url{http://www.inahta.net} ).

\begin{center}\rule{0.5\linewidth}{0.5pt}\end{center}

incidência

O número de novos casos de uma doença ou condição numa população em risco num dado período, normalmente de um ano.

incidence

The number of new cases of a disease or condition in a population at risk over a given period, usually one year.

\begin{center}\rule{0.5\linewidth}{0.5pt}\end{center}

Index Medicus

Uma lista bibliográfica de referências a artigos biomédicos de periódicos do mundo todo. O Index Medicus deixou de existir em dezembro de 2004 e foi o equivalente impresso da base de dados MEDLINE.

Index Medicus

A bibliographic listing of references to articles from biomedical journals worldwide. Note: Index Medicus ceased to exist in December 2004. It was the print equivalent of the MEDLINE database.

\begin{center}\rule{0.5\linewidth}{0.5pt}\end{center}

indicação

Um sintoma clínico, fator de risco ou circunstância para os quais o uso de uma particular intervenção poderia ser adequada, como recomendado numa guia ou protocolo de prática médica de cuidados, ou por um corpo regulatório ou outra fonte de autoridade.

indication

A clinical symptom, risk factor or circumstance for which the use of a particular intervention would be appropriate, as recommended in a clinical practice guideline or protocol of care, or by a regulatory body or other authoritative source.

\begin{center}\rule{0.5\linewidth}{0.5pt}\end{center}

Índice de Utilidades na Saúde

(\url{http://www.fhs.mcmaster.ca/hug/}). Um sistema genérico e completo de pontuação, baseado nas preferências de vários usuários e projetado para medir o estado da saúde e a qualidade de vida, e para produzir pontuações sobre utilidade.

Health Utilities Index (HUI)

(\url{http://www.fhs.mcmaster.ca/hug/}). A generic and complete scoring system based on various users' preferences and designed to measure health status and quality of life and to produce utility scores.

\begin{center}\rule{0.5\linewidth}{0.5pt}\end{center}

intermediação de conhecimentos

Uma atividade de transferência de conhecimento na qual uma terceira parte faz a ligação entre pesquisadores e os usuários dos pesquisadores, para facilitar a formação de redes de relacionamento e a troca de informações, difusão, inovação e suporte mútuo.

knowledge brokering

A knowledge-transfer activity in which a third party liaises between researchers and research users to facilitate networking and information exchange, diffusion of innovation and mutual support.

\begin{center}\rule{0.5\linewidth}{0.5pt}\end{center}

intervalo de confiança (IC)

intervalo de valores acima e abaixo do ponto estimado com uma dada probabilidade de incluir o valor verdadeiro de um dado parâmetro, tal como o efeito de um tratamento. Nota: o intervalo de confiança é a área de incerteza na estimativa de um parâmetro. O uso desse intervalo reflete o fato de que o estudo provê uma estimativa do parâmetro, fora das muitas estimativas que seriam possíveis se o estudo fosse repetido diversas vezes. Se um IC de x\% for construido a cada repetição, x\% dos intervaos conterão o valor verdadeiro do parâmetro. Investigadores frequentemente usam intervalos de confiança de 90\%, 95\% ou 99\%. Assim, um IC de 95\% indica que há 95\% de probabilidade do intervalo de confiança , calculado num estudo em particular, incluir o valor verdadeiro do parâmetro. Se o intervalo incluir o valor nulo (uma diferença de zero, em média, um risco relativo ou razão de chances de 1, ou um coeficiente de correlação de zero, por exemplo), a hipótese nula não poderá ser rejeitada. Um intervalo de confiança mais estreito ao redor de um ponto estimado indica uma estimativa mais precisa do que um intervalo de confiança mais amplo.

confidence interval (CI)

A range of values below and above the point estimate that has a given probability of including the true value of a given parameter, such as a treatment effect. Note: The confidence interval is the area of uncertainty for the estimating of a parameter. The use of this interval reflects the fact that a study provides one estimate of a parameter, out of the many estimates that would be possible if the study were repeated several times. If an X\% CI is constructed for each repetition, X\% of the intervals will contain the true value of the parameter. Investigators typically use confidence intervals of 90\%, 95\% or 99\%. Thus, a 95\% confidence interval indicates that there is a 95\% probability that the confidence interval calculated from a particular study includes the true value of the parameter. If the interval includes a null value (a difference in means of 0, an odds ratio or a relative risk of 1, or a correlation coefficient of 0, for example), the null hypothesis cannot be rejected. A narrow confidence interval around a point estimate indicates a more precise estimate than a wide confidence interval.

\begin{center}\rule{0.5\linewidth}{0.5pt}\end{center}

intervalo de credibilidade

Faixa de valores ao redor de uma estimativa central de um parâmetro, construida usando métodos Bayesianos. Nota: o intervalo de credibilidade é o equivalente ao intervalo de confiança Bayesiano, mas a interpretação é ligeriamente diferente: a probabilidade que um parâmetro esteja em um intervalo de credibilidade de X\% é X/100. Por exemplo, um intervalo de credibilidade de 95\% (0.83-1.36) para uma razão de chances para mortalidade, significa que existe uma probabilidade de 0.95 de que a razão de chance da popilação esteja entre 0.82 e 1.36.

credible interval

A range of values around the central estimate of a parameter, constructed using Bayesian methods. Note: The credible interval is the Bayesian equivalent of a confidence interval, but the interpretation of it is slightly different: the probability that a parameter is in an X\% credible interval is X/100. For example, a 95\% credible interval (0.82-1.36) for an odds ratio for mortality means that there is a 0.95 probability that the odds ratio in the population is between 0.82 and 1.36.

\begin{center}\rule{0.5\linewidth}{0.5pt}\end{center}

ISTAHC - International Society of Technology Assessment in Health Care (ISTAHC)

A primeira associação dos envolvidos em avaliação de tecnologia em saúde. A sociedade funcionou por cerca de vinte anos até ser dissolvida ao redor de 2002. Nota: ISTAHC foi substituida pela HTAi.

ISTAHC - International Society of Technology Assessment in Health Care (ISTAHC)
The first international association for those involved in health technology assessment. The society operated for nearly twenty years, until it was disbanded around 2002. Note: ISTAHC was replaced by HTAi.

\hypertarget{j}{%
\chapter*{J}\label{j}}
\addcontentsline{toc}{chapter}{J}

Justiça

O princípio de equidade que estabelece que as pessoas devem ser tratadas igualmente. Outra definição: a retidão no entendimento, reconhecimento e respeito aos direitos de cada pessoa e o requerimento de dar a cada um aquilo que a mesma está intitulada, de acordo com as adequadas regras de alocação. Justiça distributiva: o princípio de justiça aplicado à distribuição de bens comuns que estabelece o compartilhamento proporcional e um adequado uso dos recursos. (Conceito relacionado: equidade.)

Justice

The principle of fairness that states that people must be treated equally. Other definition: A rectitude in understanding, recognising and respecting the rights and merit of each person and a requirement to give to each that to which he or she is entitled, in accordance with appropriate allocation rules. Distributive justice: The principle of justice applied to the distribution of common goods that establishes a proportional sharing and an appropriate use of resources. (Related concept: equity).

\hypertarget{l}{%
\chapter*{L}\label{l}}
\addcontentsline{toc}{chapter}{L}

licenciamento

Uma autorização de marketing para medicamentos que aderem a padrões de segurança, qualidade ou eficácia.

licensing

A marketing authorisation for medicines which meet standards of safety, quality and efficacy.

\begin{center}\rule{0.5\linewidth}{0.5pt}\end{center}

lista de referência

Uma lista de documentos citados num trabalho, e usualmente encontrados ao final do texto. Nota: a citação de cada documento inclui o(s) nome(s) do(s) autor(es), o título, a fonte e o ano de publicação, e segue um estilo de apresentação específico (por exemplo, o Estilo Vancouver).

reference list

A list of the documents cited in a work, and usually found at the end of the text. Note: The entry for each document includes the name(s) of the author(s), the title, the source and the year of publication, and follows a specific presentation style (the Vancouver Style, for example).

\begin{center}\rule{0.5\linewidth}{0.5pt}\end{center}

literatura cinza

Documentos publicados para uma audiência limitada, fora dos principais canais de distribuição, e difícil de encontrar nas bases de dados usuais, tais como: apresentações em conferências, avaliações de tecnologia da saúde feitas por hospitais e certos documentos governamentais. Nota: existe uma falta de unanimidade sobre a definição de literatura cinza, pois as fronteiras dessa literatura se deslocam continuamente e evoluem à medida em que novos meios técnicos são desenvolvidos.

grey literature

The documents published for a limited audience, outside the major distribution channels, and difficult to find in the usual databases, such as presentations at conferences, health technology assessments done by hospitals, and certain government documents. Note: There is a lack of unanimity on the definition of grey literature, as the boundaries of this literature continually shift and evolve as new technical means are developed.

\begin{center}\rule{0.5\linewidth}{0.5pt}\end{center}

logaritmo da razão de chances

Um logaritmo natural usado em modelos logísticos e em representações gráficas de razões de chances

log-odds ratio

A natural logarithm used in logistic models and graphical representations of odds ratios.

\hypertarget{m}{%
\chapter*{M}\label{m}}
\addcontentsline{toc}{chapter}{M}

magnitude do efeito do tratamento

Uma avaliação subjetiva do tamanho do efeito estimado na parametrização de um ensaio clínico.

magnitude of treatment effect

A subjective evaluation of the size of the estimated effect in a clinical setting.

\begin{center}\rule{0.5\linewidth}{0.5pt}\end{center}

mascaramento

Veja cegamento.

masking

See blinding.

\begin{center}\rule{0.5\linewidth}{0.5pt}\end{center}

média
Uma medida de tendência central calculada pela divisão da soma dos valores observados pelo número de observações. Nota: a palavra média usada isoladamente geralmente se refere à média aritmética. Sinônimo: média aritmética

mean

A measure of central tendency calculated by dividing the sum of all the observed values by the number of observations. Note: The word mean used alone generally refers to the arithmetic mean. Syn.: arithmetic mean

\begin{center}\rule{0.5\linewidth}{0.5pt}\end{center}

média aritmética

veja média

arithmetic mean

see mean

\begin{center}\rule{0.5\linewidth}{0.5pt}\end{center}

mediana

Uma medida de tendência central correspondente ao valor abaixo do qual 50\% são encontradas 50\% das observações. Nota: a mediana é o ponto médio das observações classificadas em ordem ascendente. Ela pode proporcionar uma melhor estimativa da média quando valores extremos causam assimetria na distribuição das observações.

median

A measure of central tendency corresponding to the value below which 50\% of the observations are found. Note: The median is the midpoint of observations ranked in ascending order. It can provide a better estimate of the mean when extreme values cause asymmetry in the distribution of the observations.

\begin{center}\rule{0.5\linewidth}{0.5pt}\end{center}

medicina baseada em evidências

Um tipo de medicina baseado no uso da melhor evidência proveniente de pesquisa médica e científica na tomada de decisões sobre o cuidado com pacientes individuais. Nota: este enfoque envolve a formulação de uma questão clara e específica baseada num problema específico de saúde, sistematicamente pesquisando-se a literatura científica e médica sobre o problema, determinando-se os indicadores de interesse, apreciando criticamente os resultados relevantes e usando as conclusões resultantes na rotina da prática médica.

evidence-based medicine

A type of medicine based on using the best evidence from scientific and medical research to make decisions about the care of individual patients. Note: This approach involves formulating a clear and specific question based on a specific health problem, systematically searching the scientific and medical literature on the problem, determining the indicators of interest, critically appraising the relevant results and using the resulting conclusions in routine medical practice.

\begin{center}\rule{0.5\linewidth}{0.5pt}\end{center}

MEDLARS (Medical Literature Analysis and Retrieval System) (sistema de arquivamento e análise de Literatura Médica)

Um sistema computadorizado compreendendo cerca de 40 bases de dados médicas gerenciadas pela Biblioteca Nacional de Medicina dos Estados Unidos.

MEDLARS (Medical Literature Analysis and Retrieval System)

A computerized system comprising about 40 medical databases managed by the United States National Library of Medicine.

\begin{center}\rule{0.5\linewidth}{0.5pt}\end{center}

MEDLINE (MEDlars onLINE)

Uma base de dados eletrônica de artigos médicos produzida pela Biblioteca Nacional de Medicina dos Estados Unidos. Nota: em 2008, haviam indexados acima de 17 milhões de referências de cerca de 5 mil periódicos biomédicos publicados desde 1966. (Conceito relacionado: PubMed)

MEDLINE (MEDlars onLINE)

An electronic database of medical articles produced by the United States National Library of Medicine. Note: In 2008, there were over 17 million references from about 5,000 biomedical journals published since 1966 indexed in it. (Related concept: PubMed)

\begin{center}\rule{0.5\linewidth}{0.5pt}\end{center}

MeSH (Medical Subject Headings) (descritores de referências médicas)

Descritores médicos que comõe o vocabulário controlado de cerca de 19 mil termos usados na MEDLINE e em outras determinadas bases de dados da United States National Library of Medicine (Biblioteca Nacional de Medicina dos Estados Unidos).

Medical Subject Headings (MeSH)

The medical descriptors that are part of a controlled vocabulary of about 19,000 terms used in MEDLINE and certain other United States National Library of Medicine MEDLARS databases.

\begin{center}\rule{0.5\linewidth}{0.5pt}\end{center}

meta-análise

Um método estatístico que consiste em combinar sistematicamente resultados de diferentes estudos para obter uma estimativa quantitativa do efeito geral de uma determinada intervenção ou variável.

meta-analysis

A statistical method that consists of systematically combining results from different studies to obtain a quantitative estimate of the overall effect of a particular intervention or variable.

\begin{center}\rule{0.5\linewidth}{0.5pt}\end{center}

meta-análise cumulativa

Meta-análise apresentando o acúmulo gradual dos resultados de estudos, à medida em que os estudos são adicionados. Nota: no gráfico de uma meta-análise cumulativa, cada linha horizontal representa a integração dos estudos precedentes, ao invés de apresentar os resultados de um único estudo. Os estudos são integrados um a cada vez, numa ordem especificada (por exemplo, de acordo com a data da publicação ou com a qualidade).

cumulative meta-analysis

A meta-analysis presenting the gradual accumulation of the results of studies as the studies are added. Note: In a cumulative meta-analysis graph, each horizontal line represents the integration of the preceding studies, rather than the results of a single study. The studies are integrated one at a time and in a specified order (e.g.~according to the date of publication or quality).

\begin{center}\rule{0.5\linewidth}{0.5pt}\end{center}

meta-regressão

Em uma meta-análise, um modelo de regressão para estudar a relação entre as diferentes características do estudo e a estimação do efeito observado nesses estudos. Nota: As características do estudo são, por exemplo, o método de distribuição nos grupos, o tipo de cegamento, o risco inicial e o cronograma de administração de intervenção.

meta-regression

In a meta-analysis, a regression model for studying the relationship between the different study characteristics and the estimation of the effect observed in these studies. Note: The study characteristics are, for example, the method of distribution into the groups, the type of blinding, the initial risk and the intervention administration schedule.

\begin{center}\rule{0.5\linewidth}{0.5pt}\end{center}

método de Peto

Em meta análise, um método usado para combinar razões de chances. Nota: os cálculos são diretos e de fácil interpretação, mas o resultado pode às vezes diferir substancialmente dos obtidos por outros métodos. O método de Peto é baseado no modelo de efeito fixo.

Peto method

The ratio (quotient) of the exposure odds among the cases (subjects who are ill or have died, for example) to that among the controls (subjects who are not ill or who are alive, for example).

\begin{center}\rule{0.5\linewidth}{0.5pt}\end{center}

método do custo-fricção

Um método para se estimar o custo de produção perdido, resultante da ausência de um empregado durante o período necessitado pela organização para substituir o empregado e promover o retorno ao nível inicial de produtividade.

friction-cost method

A method of estimating the cost of the production loss resulting from the absence of an employee during the period needed by the organization to replace the employee and return to the initial level of productivity.

\begin{center}\rule{0.5\linewidth}{0.5pt}\end{center}

método do perfil de confiança

meta-análise baseada em estatística Bayesiana, para combinar resultados de múltiplos estudos que tem diferentes projetos de pesquisa, tais como ensaios randomizados controlados e estudos observacionais, pelo ajuste dos resultados de cada estudo individual em função de seus viéses metodológicos antes de se combinar os resultados numa distribuição de probabilidade para o(s) parâmetro(s) de interesse.

confidence profile method

A meta-analysis based on Bayesian statistics for combining results of multiple studies that have different research designs, such as randomized controlled trials and observational studies, by adjusting the results of each individual study for its methodological biases before combining the results into a probability distribution for the parameter(s) of interest.

\begin{center}\rule{0.5\linewidth}{0.5pt}\end{center}

mini-HTA

Relatório que inclui uma revisão compreensiva e sistemática de literatura, ou uma revisão sistemática de evidências de alto nível, avaliando a segurança e a efetividade de uma tecnologia, porém sem incluir uma análise de custo-efetividade da tecnologia. Nota 1: o relatório sempre: descreve as características do uso atual da tecnologia; criticamente avalia a qualidade da base de evidência; provê informação sobre custos e impacto financeiro. Nota 2: o relatório pode opcionalmente discutir considerações organizacionais. Nota 3: um ``mini-HTA'' pode ser tão rigoroso quanto um relatório HTA, mas tipicamente tem um escopo restrito para ser produzido rapidamente. Nota 4: termos relacionados incluem Avaliação de tecnologia em saúde, relatório HTA, revisão rápida, nível de evidência.

mini-HTA

A report that includes a comprehensive systematic literature review, or a systematic review of high level evidence, evaluating the safety and effectiveness of a technology, but that does not include an analysis of the cost-effectiveness of the technology. Note 1: The report always:
• describes the characteristics and current use of the technology;
• critically appraises the quality of the evidence base; and
• provides information on costs and financial impact. Note 2: The report may optionally discuss organisational considerations Note 3: A ``mini-HTA'' may be as rigorous as an HTA report but typically has a restricted scope and so is quicker to produce. Note 4: Related terms include health technology assessment, HTA report, rapid review, level of evidence.

\begin{center}\rule{0.5\linewidth}{0.5pt}\end{center}

minimização

Um método de alocação da intervenção para obter grupos de estudo que sejam semelhantes para várias variáveis. Nota: este método pode ser usado sem um componente de randomização.

minimization

A method of intervention allocation for obtaining study groups that are very similar for several variables. Note: This method can be used with or without a randomization component. Alternate spelling: minimisation.

\begin{center}\rule{0.5\linewidth}{0.5pt}\end{center}

mistura de casos

todas as características de um estudo populacional que podem influenciar o efeito ou a escolha do tratamento e que devem ser levados em conta ao se estabelecer o desfecho do tratamento. Nota: exemplos incluem severidade de doenças e condições coexistentes

case-mix

All the characteristics of a study population that may influence the effect or the choice of a treatment and that must be taken into account when assessing treatment outcomes. Note: Examples include disease severity and coexisting conditions.

\begin{center}\rule{0.5\linewidth}{0.5pt}\end{center}

modelo aditivo

um modelo no qual o efeito combinado de diversos fatores independentes é a soma dos efeitos isolados de cada fator. Nota: por exemplo, se um fator X aumenta um risco por a na ausência de Y, e um segundo fator Y aumenta o riso por b na ausência de X, o efeito combinado dos dois fatores é a+b.

additive model

A model in which the combined effect of several independent factors is the sum of the isolated effects of each factor. Note: For example, if a factor X increases a risk by a in the absence of Y, and if a second factor Y increases the risk by b in the absence of X, the combined effect of the two factors is a + b.

\begin{center}\rule{0.5\linewidth}{0.5pt}\end{center}

modelo de efeito fixo

Em meta-análise, um modelo estatístico que leva em conta apenas a variabilidade interna ao estudo, na avaliação do grau de incerteza (intervalo de confiança) do efeito combinado obtido em estudos. Nota: o modelo de efeito fixo assume que as unidades analisadas são unidades de interesse e, consequentemente, que compõe a população em unidades. Nesse modelo, a variação entre as estimativas do efeito de cada estudo (heterogeneidade) não afeta o intervalo de confiança. Conceitos relacionados: modelo de efeito randômico, método de Peto.

fixed-effect model

In a meta-analysis, a statistical model that takes only within-study variability into account in assessment of the degree of uncertainty (confidence interval) of the combined effect from the studies. Note: The fixed-effect model assumes that the units analysed are the units of interest and, consequently, that they make up the population of units. In such a model, variation between the estimates of effect from each study (heterogeneity) does not affect the confidence interval. (Related concepts: random effects model, Peto method)

\begin{center}\rule{0.5\linewidth}{0.5pt}\end{center}

modelo de efeito randômico

Numa meta análise, um modelo estatístico no qual tanto a variância interna e as variâncias entre estudos são incluidas na avaliação do grau de incerteza (intervalo de confiança) do efeito combinado dos estudos. Nota: se há uma heterogeneidade significativa entre os resultados dos estudos incluidos na meta análise, um modelo randomizado de efeito poderá proporcionar intervalos de confiança mais largo do que no modelo de efeito fixo. (Conceito relacionado: modelo de efeito fixo).

random effects model

In a meta-analysis, a statistical model in which both the within-study variance and the between-studies variances are included in the assessment of the degree of uncertainty (confidence interval) of the combined effect of the studies. Note: If there is significant heterogeneity among the results of the studies included in a meta-analysis, a random effects model will give wider confidence intervals than a fixed-effect model. (Related concepts: fixed-effect model)

\begin{center}\rule{0.5\linewidth}{0.5pt}\end{center}

modelo de Markov

Um tipo de modelagem quantitativa que envolve um conjunto especificado de exaustivos estados de saúde, mutuamente exclusivos, para os quais existem probabilidades transacionais de movimentação de um estado a outro, incluindo a probabilidade de permanecer no mesmo estado. Nota: tipicamente, estados possuem um período de tempo uniforme e as probabilidades transacionais permanecem constantes ao longo do tempo.

Markov model

A type of quantitative modelling that involves a specified set of mutually exclusive and exhaustive health states for which there are transitional probabilities of moving from one state to another, including the probability of remaining in the same state. Note: Typically, states have a uniform time period, and transitional probabilities remain constant over time.

\begin{center}\rule{0.5\linewidth}{0.5pt}\end{center}

modelo econômico

No campo da saúde, um modelo matemático dos passos clínicos que representam as escolhas-chave e as consequências das opções sendo estudadas. Nota: um modelo econômico pode ser usado para extrapolar, à partir de desfechos intermediários, os desfechos de longo prazo importantes para os pacientes.

economic model

In the health field, a mathematical model of the clinical pathway that represents the key choices and the consequences of the options studied. Note: An economic model can be used to extrapolate, from intermediate outcomes, the long-term outcomes of importance to patients.

\begin{center}\rule{0.5\linewidth}{0.5pt}\end{center}

modelo logístico

Um modelo de regressão estatística que estima a probabilidade de um valor de uma variável dicotômica com base em múltiplas variáveis preditoras.

logistic model

A statistical regression model that estimates the probability of a value of a dichotomous variable on the basis of multiple predictor variables.

\begin{center}\rule{0.5\linewidth}{0.5pt}\end{center}

modelo multiplicativo

Um modelo no qual o efeito combinado de dois ou mais fatores é o produto dos efeitos isolados de cada fator. Nota: por exemplo, se um fator X multiplica um risco por a na ausência de Y, e um segundo fator Y multiplica o risco por b na ausência de X, o efeito combinado desses dois fatores será a x b. (conceito relacionado: modelo aditivo).

multiplicative model

A model in which the combined effect of two or more factors is the product of the isolated effects of each factor. Note: For example, if a factor X multiplies a risk by a in the absence of Y, and if a second factor Y multiplies the risk by b in the absence of X, the combined effect of the two factors will be a x b. (Related concept: additive model)

\hypertarget{n}{%
\chapter*{N}\label{n}}
\addcontentsline{toc}{chapter}{N}

nível de evidência

Veja hierarquia de evidência.

level of evidence

See hierarchy of evidence.

\begin{center}\rule{0.5\linewidth}{0.5pt}\end{center}

nova tecnologia em saúde

Uma tecnologia em saúde que está em fase de lançamento, logo após os estágios de pós-marketing ou dos estágios iniciais de difusão. Nota 1: termos relacionados incluem tecnologia emergente em saúde.

new health technology

A health technology that is in the launch, early post-marketing or early diffusion stages. Note 1: Related terms include emerging health technology.

\begin{center}\rule{0.5\linewidth}{0.5pt}\end{center}

número necessário para produzir um dano (NNH)

O número de indivíduos tratados por indivíduo experimentando um efeito adverso. Nota: este conceito é o inverso da diferença de risco absoluto. O termo foi criado como uma variante do número necessário para tratamento.

number needed to harm (NNH)

The number of subjects treated per subject experiencing an adverse effect. Note: This concept is the inverse of the absolute risk difference. The term was created as a variant of number needed to treat.

\begin{center}\rule{0.5\linewidth}{0.5pt}\end{center}

número necessário para tratar (NNT)

O número de indivíduos tratados por indivíduo com um desfecho favorável. Nota: é o inverso da redução do risco absoluto (1 ÷ redução do risco absoluto). Assim, se os resultados de um estudo indicarem que a probabilidade de morte num grupo controle for 25\% e a probabilidade de morte num grupo tratamento dor 10\%, o número necessário para tratar deveria ser 1 ÷ (0,25 - 0,10) = 6,7, ou seja, 7 indivíduos. Conceitos relacionados: diferença de risco, redução do risco absoluto, redução do risco relativo e proporcão de chances.

number needed to treat (NNT)

The number of subjects treated per subject with a favourable outcome. Note: It is the inverse of absolute risk reduction (1 ÷ absolute risk reduction). Thus, if the results of a study indicate that the probability of death in a control group is 25\% and the probability of death in a treatment group is 10\%, the number needed to treat would be 1.0 ÷ (0.25 - 0.10) = 6.7, therefore 7 subjects. Related concepts: risk difference, absolute risk reduction, relative risk reduction and odds ratio

\hypertarget{o}{%
\chapter*{O}\label{o}}
\addcontentsline{toc}{chapter}{O}

ocultação da distribuição

é o processo usado para se distribuir grupos de pacientes num RCT (ensaio controlado e randomizado) de forma a evitar a previsão (pela pessoa gerenciando a alocação bem como os pacientes dessa distribuição. Números de registros médicos, números de identificação de pessoas ou datas de aniversário não são adequados para a ocultação da distribuição. Certos esquemas centralizados de randomização e números sequenciais selados, envelopes opacos estão entre os métodos adequados de ocultação da distribuição (veja também cegamento).

concealment of allocation

The process used to assign patients to alternative groups in an RCT in a manner that prevents foreknowledge (by the person managing the allocation as well as the patients) of this assignment. Medical record numbers, personal identification numbers, or birthdays are not adequate for concealment of allocation. Certain centralized randomization schemes and sequentially numbered sealed, opaque envelopes are among adequate methods of allocation concealment. (See also blinding)

\hypertarget{p}{%
\chapter*{P}\label{p}}
\addcontentsline{toc}{chapter}{P}

paciente

Uma pessoa, apresentando sinais clínicos ou não, que consulta um médico. Veja consumidor.

patient

A person, presenting with clinical signs or not, who consults a physician. See consumer.

\begin{center}\rule{0.5\linewidth}{0.5pt}\end{center}

padrão do melhor interesse

O princípio da beneficência em que deve basear-se em qualquer decisão relativa a um paciente que não consegue expressar seus desejos.

best interest standard

The beneficence principle on which must be based any decision relating to a patient who is unable to express his or her wishes.

\begin{center}\rule{0.5\linewidth}{0.5pt}\end{center}

padrão-ouro

O método, procedimento ou medida que é largamente aceita como sendo a melhor, e contra a qual novos procedimentos devem ser comparados. Nota: é particularmente importante usar um padrão-ouro em estudos de acurácia de testes diagnósticos.

gold standard

The method, procedure or measurement that is widely accepted as being the best and that against which new ones should be compared. Note: It is particularly important to use such a gold standard in studies of the accuracy of diagnostic tests.

\begin{center}\rule{0.5\linewidth}{0.5pt}\end{center}

perspectiva analítica

o ponto de vista escolhido para uma dada análise, incluindo o ponto de vista da sociedade, do governo, do sistema de saúde ou do contribuinte.

analytic perspective

The viewpoint chosen for a given analysis, including a viewpoint of society, the government, the health care system or the payer.

\begin{center}\rule{0.5\linewidth}{0.5pt}\end{center}

pesquisa com ação participativa

veja pesquisa de ação

participatory action research

See action research

\begin{center}\rule{0.5\linewidth}{0.5pt}\end{center}

pesquisa de serviços de saúde

Um campo interdisciplinar de pesquisa que examina o impacto: da organização, do financiamento e da gestão de serviços de cuidados com a saúde administrando: entrega, qualidade, custo, acessibilidade e desfechos de tais serviços.

health services research

An interdisciplinary field of inquiry that examines the impact of the organization, financing and management of health care services on the delivery, quality, cost, accessibility and outcomes of such services.

\begin{center}\rule{0.5\linewidth}{0.5pt}\end{center}

pesquisa manual

A busca metodológica manual, página a página, do conteúdo de um jornal, incluindo artigos, editoriais, cartas dos leitores, por exemplo, com a finalidade de identificar estudos relevantes e completar as buscas não indexadas à base de dados. Nota: normalmente, se começa à partir do ano corrente e então se retrocede no tempo, até a produção se tornar desprezível ou até se encontrar o primeiro volume.

handsearching

The methodological searching, page by page (by hand) of the content of a journal, including articles, editorials, letters from readers, etc., to identify the relevant studies and complete the non-indexed searching in the databases. Note: Normally, one starts with the issues from the current year, and then works backwards until the yield becomes negligible or the first volume is reached.

\begin{center}\rule{0.5\linewidth}{0.5pt}\end{center}

pesquisa primária

Investigação na qual os dados são coletados pela primeira vez diretamente dos pacientes (ensaio randomizado controlado, estudo observacional, séries de casos, etc.) Nota: O termo ``pesquisa primária'' é frequentemente utilizado para distinguir esse tipo de pesquisa da ``pesquisa secundária'', que é uma pesquisa sintética (re-análise de dados previamente coletados), meta-análise e outros tipos de estudos combinados (tais como análises econômicas e análises de decisão). Conceito realacionado: pesquisa secundária. Sinônimos: estudo primário, estudo original.

primary research

An investigation in which the data are collected for the first time directly from patients (randomized controlled trial, observational study, case series, etc.). Note: The term ``primary research'' is often used to distinguish this type of research from ``secondary research,'' which is synthesis research (re-analysis of previously collected data), meta-analyses and other ways of combining studies (such as economic analysis and decision analysis). Related concept: secondary research. Syn.: primary study, original study.

\begin{center}\rule{0.5\linewidth}{0.5pt}\end{center}

pesquisa secundária

Um tipo de pesquisa que não produz dados originais, mas envolve a síntese qualitativa ou quantitativa de informação a partir de múltiplos estudos originais. Nota: revisões de literatura, meta análises, análises de decisão e relatórios de consenso são exemplos de pesquisa sintética. Sinônimo: estudo sintético.

secondary research

A type of research that does not produce original data, but that involves the qualitative or quantitative synthesis of information from multiple original studies. Note: Literature reviews, meta-analyses, decision analyses and consensus reports are examples of synthesis research. Syn. : synthetic study.

\begin{center}\rule{0.5\linewidth}{0.5pt}\end{center}

placebo

Num cenário clínico, um tratamento ou intervenção que é inativo, mas que pode agir através de um mecanismo psicológico (efeito placebo) se o paciente acredita estar recebendo um tratamento ou uma intervenção ativos. Em pesquisa, um tratamento fictício dado a um grupo controle, impossível de ser diferenciado do tratamento real dado ao grupo de tratamento, usado para isolar o efeito do tratamento sendo estudado.

placebo

The effect on a patient's health status (improvement or worsening) that may occur due to the patient's expectation that a particular intervention will have an effect.

\begin{center}\rule{0.5\linewidth}{0.5pt}\end{center}

poder

Em teste de hipóteses, a probabilidade de rejeição de uma falsa hipótese nula, isto é, concluindo que um tratamento é efetivo quando o efeito do tratamento de uma dada magnitude realmente existe, por exemplo. Nota: na presença de uma verdadeiramente falsa hipótese nula (quando há um efeito verdadeiro do tratamento, de uma dada magnitude, por exemplo), o poder é a probabilidade de se evitar um erro do Tipo II, e assim corresponde a (1 - β).

power

In hypothesis testing, the probability of rejecting a false null hypothesis, i.e.~of concluding that a treatment is effective when a treatment effect of a given magnitude truly exists, for example. Note: In the presence of a truly false null hypothesis (when there is a true treatment effect of a given magnitude, for example), the power is the probability of avoiding a Type II error and therefore corresponds to (1 - β).

\begin{center}\rule{0.5\linewidth}{0.5pt}\end{center}

poder da evidência

A significância clínica e estatística dos resultados de um estudo, determinados pela magnitude, precisão e reprodutibilidade do efeito da intervenção sendo estudada. Nota: em estudos não randomizados, fatores adicionais tais como a plausibilidade biológica, gradente biológico e consistência longitudinal de associações, podem ser considerados.

strength of evidence

The clinical and statistical significance of the results of a study, determined by the magnitude, precision and reproducibility of the effect of the intervention being studied. Note: In nonrandomised studies, additional factors, such as biological plausibility, biological gradient and longitudinal consistency of associations, may be considered.

\begin{center}\rule{0.5\linewidth}{0.5pt}\end{center}

poder estatístico

Veja poder.

Statistical power

See power.

\begin{center}\rule{0.5\linewidth}{0.5pt}\end{center}

ponto final

Um indicador escolhido para determinar o efeito de uma intervenção.

endpoint

An indicator chosen for determining the effect of an intervention.

\begin{center}\rule{0.5\linewidth}{0.5pt}\end{center}

prática inadequada

Uma avaliação subjetiva do tamanho do efeito estimado na parametrização de um ensaio clínico.

malpractice

A subjective evaluation of the size of the estimated effect in a clinical setting.

\begin{center}\rule{0.5\linewidth}{0.5pt}\end{center}

precisão

\begin{enumerate}
\def\labelenumi{\arabic{enumi})}
\tightlist
\item
  A qualidade de uma estimativa pontual obtida de um conjunto de observações tendo uma variância pequena. Nota: um intervalo de confiança estreito ao redor de um ponto estimado indica uma estimativa mais precisa do efeito, do que um intervalo de confiança largo. Note que uma estimativa precisa não é necessariamente acurada. 2) Uma medida de verossimilhança de erros randômicos nos resultados de um estudo, meta análise ou medida. Nota: em meta análise, o peso dado aos resultados de cada estudo sobre a estimativa geral do efeito de uma intervenção é frequentemente baseado na precisão de cada estudo, a qual é estimada usando-se o inverso da variância da estimativa do efeito do tamanho da amostra. 3) em pesquisa de literatura, o número de citações relevantes dividido pelo número total de citações coletadas, isto é, a proporção de estudos atendendo o critério de inclusão em registros de ensaios clínicos ou numa revisão literária. (termo relacionado: acurácia).
\end{enumerate}

precision

\begin{enumerate}
\def\labelenumi{\arabic{enumi})}
\tightlist
\item
  A quality of a point estimate obtained from a set of observations having a small variance. Note: A narrow confidence interval around a point estimate indicates a more precise estimate of effect than a wide confidence interval. Note that a precise estimate is not necessarily accurate. 2) A measure of the likelihood of random errors in the results of a study, meta-analysis or measurement. Note: In a meta-analysis, the weight given to the results of each study in the overall estimate of the effect of an intervention is often based on the precision of each study, which is estimated using the inverse of the variance of the estimate of effect or the sample size. 3) In a literature search, the number of relevant citations, divided by the total number of citations retrieved, i.e.~the proportion of studies meeting the inclusion criteria for a clinical trials register or a literature review. (Related concept: accuracy)
\end{enumerate}

\begin{center}\rule{0.5\linewidth}{0.5pt}\end{center}

preferência

No contexto da saúde, o quão desejável é um desfecho clínico. Nota: utilidade e valor são aplicações específicas deste conceito.

preference

In the context of health, the desirability of a clinical outcome. Note: Utility and value are specific applications of this concept.

\begin{center}\rule{0.5\linewidth}{0.5pt}\end{center}

preferência revelada

Uma preferência deduzida das escolhas que um indivíduo faz. Nota: preferências podem ser reveladas por escolhas feitas na definição natural ou nas respostas às questões de um investigador.

revealed preference

A preference deduced from the choices that an individual makes. Note: Preferences may be revealed by choices made in a natural setting or by responses to an investigator's questions.

\begin{center}\rule{0.5\linewidth}{0.5pt}\end{center}

prevalência

O número de pessoas de uma população com uma doença ou condição, num dado momento, usualmente expressa como a proporção do número de pessoas afetadas pelo total da população.

prevalence

The number of people in a population with a specific disease or condition at a given time, usually expressed as a proportion of the number of affected people to the total population.

\begin{center}\rule{0.5\linewidth}{0.5pt}\end{center}

previsão clínica

ferramenta para assistir a decisão clínica, podendo ser usada para determinar a probabilidade de resposta ao tratamento ou para propor um diagnóstico ou o curso de uma ação terapêutica, baseada nas váriáveis obtidas do histórico do paciente, exames físicos ou resultados de testes.

clinical prediction

A tool for assisting clinical decision making that can be used to determine the probability of response to a treatment or to propose a diagnostic or therapeutic course of action on the basis of variables obtained from the patient's history, physical exam or test results.

\begin{center}\rule{0.5\linewidth}{0.5pt}\end{center}

probabilidade de transição

A probabilidade de que a saúde de um paciente se altere de um estado a outro dentro de um dado período.

transition probability

The probability that the health of a patient changes from one state to another within a given period.

problemas de alvo móvel

Mudanças nos cuidados de saúde que podem tornar as avaliações da tecnologia da saúde desatualizadas, às vezes mesmo antes que as recomendações que delas saem possam ser implementadas. Nota: São incluídas mudanças na tecnologia mãe, em uma de suas variantes ou em uma tecnologia complementar (usada para tratar um determinado problema de saúde), emergência de novas tecnologias concorrentes e mudanças na aplicação de uma tecnologia (por exemplo, para diferentes populações de pacientes ou a diferentes problemas de saúde). Ortografia alternativa: problema de destino móvel

\begin{center}\rule{0.5\linewidth}{0.5pt}\end{center}

moving target problem

Changes in health care that can render health technology assessments out-of-date, sometimes even before the recommendations that come out of them can be implemented. Note: Included are changes in the mother technology, in one of its variants or in a complementary technology (used to treat a given health problem), emergence of new competing technologies, and changes in the application of a technology (e.g.~to different patient populations or to different health problems). Alternative spelling: moving-target problem

\begin{center}\rule{0.5\linewidth}{0.5pt}\end{center}

processo de julgamento

Um processo no qual um documento é enviado a editores e revisores experientes, especialistas no tema do documento, a especialistas em metodologia ou a usuários.

referee process

A process whereby a document is sent to expert editors and reviewers who are specialists in the field, methodological specialists or users.

\begin{center}\rule{0.5\linewidth}{0.5pt}\end{center}

prognose

Uma avaliação do curso futuro do desfecho de uma doença.

prognose

An assessment of the future course and outcome of a disease.

\begin{center}\rule{0.5\linewidth}{0.5pt}\end{center}

projeto de cruzamento

Projeto de pesquisa no qual os elementos recebem, em sequência, a intervenção experimental (ou a intervenção de controle) e então, após um tempo especificado, (fase sem intervenção), a intervenção de controle (ou a intervenção experimental). Nota: neste projeto, cada elemento atua no controle de si mesmo(a), e a randomização pode ser usada para determinar a ordem na qual o elemento recebe a intervenção experimental e a intervenção de controle.

crossover design

A research design in which the subjects receive, in sequence, the experimental intervention (or the control intervention), then, after a specified time, (phase without intervention), the control intervention (or the experimental intervention). Note: In this design, each subject serves as his/her own control, and randomisation may be used to determine the order in which a subject receives the experimental intervention and the control intervention.

\begin{center}\rule{0.5\linewidth}{0.5pt}\end{center}

projeto fatorial

Um projeto que permite a avaliação de efeitos de múltiplas intervenções e suas interações dentro de um ensaio clínico. Nota: a maioria dos ensaios consideram apenas uma interveção experimental, comparando-a com uma ou mais opções de placebo. Num ensaio utilizando um projeto fatorial 2x2, os participantes são alocados em uma de quatro possíveis combinações. Por exemplo, num ensaio randomizado controlado de terapia de substituição da nicotina e aconselhamento, os participantes poderiam ser alocados nos seguintes grupos: somente com a terapia de substituição da nicotina, somente aconselhamento, ambos ou nenhum. Assim, é possível testar o efeito independente de cada intervenção na cessação do fumo e o efeito combinado (da interação entre) as duas intervenções.

factorial design

A design that allows for evaluating the effects of multiple interventions and their interactions within a clinical trial. Note: Most trials consider only one experimental intervention, comparing it with one or more options or a placebo. In a trial using a 2x2 factorial design, participants are allocated to one of four possible combinations. For example, in a randomised controlled trial of a nicotine replacement therapy and counselling, participants would be allocated to: nicotine replacement therapy alone, counselling alone, both, or neither. In this way it is possible to test the independent effect of each intervention on smoking cessation and the combined effect of (interaction between) the two interventions.

\begin{center}\rule{0.5\linewidth}{0.5pt}\end{center}

projeto pareado

Um projeto de pesquisa no qual os indivíduos ou grupos de indivíduos são pareados (com base em fatores de prognóstico, por exemplo) e um membro de cada par é alocado no grupo experimental e o outro no grupo controle.

paired design

A research design in which the subjects or groups of subjects are matched (on the basis of prognostic factors, for example) and one member of each pair is allocated to the experimental group and the other to the control group.

\begin{center}\rule{0.5\linewidth}{0.5pt}\end{center}

proporção de risco

Veja risco relativo.

risk ratio

See relative risk.

\begin{center}\rule{0.5\linewidth}{0.5pt}\end{center}

proporção do custo-efetividade incremental

O custo adicional da intervenção mais dispendiosa comparado com o custo da intervenção menos dispendiosa, dividido pela diferença entre os efeitos das intervenções nos pacientes (o custo adicional por QALY, por exemplo).

incremental cost-effectiveness ratio (ICER)

The additional cost of the more expensive intervention compared with the less expensive intervention, divided by the difference between the effects of the interventions on the patients (the additional cost per QALY, for example).

\begin{center}\rule{0.5\linewidth}{0.5pt}\end{center}

protocolo

O plano ou conjunto de passos a serem seguidos num estudo. Nota: por exemplo, um protocolo para uma revisão sistemática deve descrever o propósito da revisão, seus objetivos e os métodos a serem usados para localizar, selecionar e avaliar criticamente os estudos incluídos, e para coletar e analisar os dados dos mesmos.

protocol

The plan or set of steps to be followed in a study. Note: For example, a protocol for a systematic review must describe the purpose of the review, its objectives and the methods to be used to locate, select and critically appraise the included studies, and to collect and analyse data from them.

\begin{center}\rule{0.5\linewidth}{0.5pt}\end{center}

protocolo per

Uma análise que exclui indivíduos que não aderiram ao regime da intervenção (fizeram tratamentos concorrentes, por exemplo). Nota: a população de interesse é assim altamente selecionada e não influenciada por fatores que possam diluir o efeito do tratamento. Esta forma de análise pode assim produzir uma maior estimativa do efeito do tratamento. (conceito relacionado: análise da intenção de tratamento). Outra definição: uma análise de um subconjunto de indivíduos de um ensaio randomizado controlado que aderiu ao protocolo suficientemente para ter uma boa probabilidade dos seus desfechos demonstrarem o feito da intervenção de interesse. Este subconjunto pode ser definido com base na exposição à intervenção, à disponibilidade das medidas e à falta de maiores resultados do protocolo. Esta estratégia de análise pode estar sujeita a viéses porque a não aderência ao protocolo pode ser devida à intervenção.

per protocol

An analysis that excludes subjects who were not compliant with the intervention regimen (took prohibited concurrent treatments, for example). Note: The population of interest is thus highly filtered and not influenced by factors that could dilute the effect of treatment. This form of analysis may therefore produce a larger estimate of treatment effect. (Related concept: intention-to-treat analysis). Other definition: An analysis of the subset of subjects of a randomised controlled trial who have complied with the protocol sufficiently for there to be a good chance that their outcomes demonstrate the effect of the intervention of interest. This subset may be defined on the basis of exposure to the intervention, the availability of measurements and the lack of major departures from the protocol. This analysis strategy may be subject to bias because the non-compliance with the protocol may be due to the intervention.

\begin{center}\rule{0.5\linewidth}{0.5pt}\end{center}

pseudo-randomização

Veja randomização e exclusão de alocação.

pseudorandomization

See randomization and concealment of allocation. Alternate spellings: pseudorandomisation, pseudo-randomisation, pseudo-randomization.

\begin{center}\rule{0.5\linewidth}{0.5pt}\end{center}

PubMed

Um banco de dados da Biblioteca Nacional de Medicina dos Estados Unidos que inclui cerca de 18 milhões de citações à partir do MEDLINE e de outros periódicos de artigos biomédicos publicados desde o início do século vinte. Nota: Pubmed inclui links para sites provendo o texto completo dos artigos. Inclui não somente a base de dados MEDLINE mas também a base OLDMEDLINE, artigo em processo e registros de artigos fornecidos por editoras. Está disponível, sem custo, na Internet em \url{http://www.ncbi.nlm.nih.gov/sites/entrez?db=pubmed}.

PubMed

An electronic database of medical articles produced by the United States National Library of Medicine.

\hypertarget{q}{%
\chapter*{Q}\label{q}}
\addcontentsline{toc}{chapter}{Q}

qualidade metodológica

A medida em que o projeto, a conduta e a análise de um estudo provavelmente terão impedido um erro sistemático e ter produzido resultados confiáveis e válidos. Nota 1: Estudos de melhor qualidade são mais propensos a produzir resultados que são menos tendenciosos e mais próximos do valor verdadeiro para um resultado do que estudos de qualidade mais fraca.

methodological quality

The extent to which the design, conduct and analysis of a study are likely to have preventedsystematic error and to have produced results that are reliable and valid. Note 1: Better quality studies are more likely to yield results that are less biased and closer to the true value for an outcome than poorer quality studies.

\begin{center}\rule{0.5\linewidth}{0.5pt}\end{center}

qualidade

Todas as características de um bem ou serviço que vem de encontro às necessidades e expectativas do usuário, continuamente. Nota: qualidade pode ser definida usando-se diversos critérios: confiabilidade, disponibilidade, sustentabilidade, segurança no uso, custo do uso, etc. Qualidade é medida usando-se indicadores de estrutura, processo e desfecho.

quality

All the characteristics of a good or service that enable it to meet users' needs and expectations on an ongoing basis. Note: Quality can be defined using several criteria: reliability, availability, sustainability, operating safety, cost of use, etc. Quality is measured using structure, process and outcome indicators.

\begin{center}\rule{0.5\linewidth}{0.5pt}\end{center}

qualidade da evidência

O grau para o qual o projeto de pesquisa e a condução do estudo nos quais a evidência está baseada, tem feito isso possível para se obter resultados válidos.

quality of evidence

The degree to which the research design and the conducting of the study on which the evidence is based have made it possible to obtain valid results.

\begin{center}\rule{0.5\linewidth}{0.5pt}\end{center}

qualidade de vida

veja qualidade de vida relacionada à saúde.

quality of life (QOL)
See health-related quality of life.

\begin{center}\rule{0.5\linewidth}{0.5pt}\end{center}

qualidade de vida relacionada à saúde

As medidas do impacto de uma intervenção no estado de saúde de pacientes, indo além das tradicionais medidas de mortalidade e morbidade para incluir dimensões tais como: fisiologia, função, vida social, cognição, emoções, dormir e descansar, energia e vitalidade, percepção de saúde e satisfação com a vida em geral. Nota: alguns desses elementos são também chamados de estado da saúde, estado funcional ou medidas de qualidade de vida.

health-related quality of life (HRQL,HRQoL,QOL)

The measures of the impact of an intervention on patients' health status, extending beyond the traditional measures of mortality and morbidity to include dimensions such as physiology, function, social life, cognition, emotions, sleep and rest, energy and vitality, health perception and general life satisfaction. Note: Some of these elements are also called health status, functional status or quality-of-life measures.

\begin{center}\rule{0.5\linewidth}{0.5pt}\end{center}

qualidade do cuidado

O grau para o qual o cuidado com a saúde aumenta a probabilidade de se obter o desfecho desejado e atender os padrões de cuidados. (Conceitos relacionados: avaliação da qualidade e garantia da qualidade).

quality of care

The degree to which health care increases the probability of achieving the desired health outcomes and meets the standards of care. (Related concepts: quality assessment and quality assurance).

\hypertarget{r}{%
\chapter*{R}\label{r}}
\addcontentsline{toc}{chapter}{R}

randomização

A alocação randômica de indivíduos num ensaio clínico para os grupos de intervenção e de controle, usando mecanismos tais como uma tabela de números aleatórios ou números aleatórios gerados por um computador. Nota: este tipo de alocação reduz potenciais viéses na atribuição de um indivíduo porque tende a neutralizar fatores conhecidos e desconhecidos de prognóstico, pela dispersão destes entre os grupos de controle e de intervençõ. A randomização é uma condição básica para a validação de muitos testes estatísticos. A alocação não randômica sistemática, ou não estritamente randômica, que aloca os indivíduos com base em elementos tais como o dia da semana, nome, data de nascimento, etc. não é equivalente à randomização e pode levar a sérios viéses. (Conceito relacionado: cegamento).

randomization

The random allocation of subjects of a clinical trial to the intervention and control groups using mechanisms such as a random number table or a computer-generated random number list. Note: This type of allocation reduces potential bias in subject assignment because it tends to neutralise known and unknown prognostic factors by spreading them evenly between the intervention and control groups. Randomisation is a basic condition to the validity of many statistical tests. Non-random systematic allocation, or quasi randomisation, which allocates the subjects on the basis of elements such as day of the week, name, date of birth, etc., is not equivalent to randomisation and can lead to serious biases. (Related concept: blinding). Alternate spelling: randomisation.

\begin{center}\rule{0.5\linewidth}{0.5pt}\end{center}

randomização estratificada

Um método para assegurar que um número igual de participantes com uma característica que possa influenciar o prognóstico ou a resposta à intervenção possa ser alocada a cada um dos grupos comparados, para reduzir o potencial viés de confusão. Por exemplo, num estudo de câncer no seio, pode ser importante ter a mesma proporção de mulheres pré e pós a menopausa em cada grupo. Para essa finalidade, a randomização é feita separadamente para cada estrato de estado da menopausa (frequentemente usando a randomização permutada em blocos) ou por meio de minimização.

stratified randomization

A method for ensuring that an equal number of participants with a characteristic that may influence the prognosis or the response to the intervention will be allocated to each of the compared groups, to reduce the potential for confusion bias. Note: For example, in a breast cancer study, it may be important to have the same proportion for pre-menopausal and post-menopausal women in each group. To that end, the randomisation is done separately for each menopausal status stratum (often using random permuted blocks) or by means of minimization. Alternate spelling: stratified randomisation.

\begin{center}\rule{0.5\linewidth}{0.5pt}\end{center}

randomização permutada em blocos

Um método de randomização que garante que, em qualquer ponto de um ensaio, um número igual de indivíduos foram alocados em cada grupo sendo comparado. Nota: permutação em blocos é frequentemente utilizada em combinação com a randomização estratificada.

random permuted block

A method of randomization that ensures that, at any point in a trial, an equal number of subjects have been allocated to each comparison group. Note: Permuted blocks are often used in combination with stratified randomization.

\begin{center}\rule{0.5\linewidth}{0.5pt}\end{center}

razão de chances (OR)

A proporção (o quociente) da probabilidade da exposição entre os casos (de indivíduos que estão doentes ou que morreram, por exemplo) com relação os casos de controle (indivíduos que não estão doentes ou que estão vivos, por exemplo). Nota: a razão de chances é a estimativa do risco relativo usado num tipo de estudo (caso controle) no qual é impossível calcular diretamente o risco relativo. É uma boa estimativa do risco relativo nos casos nos quais a doença é rara. Assim, se os resultados de um ensaio são de que a probabilidade de morte é de 25\% no grupo controle e de 10\% no grupo de intervenção, a razão de chances de sobrevivência é {[}0,10 ÷ (1 - 0,10){]} ÷ {[}(0,25 ÷ (1 - 0,25){]}=0,33. (conceitos relacionados: redução do risco absoluto, número necessário para tratar e risco relativo)

odds ratio (OR)

The ratio (quotient) of the exposure odds among the cases (subjects who are ill or have died, for example) to that among the controls (subjects who are not ill or who are alive, for example). Note: The odds ratio is the estimate of relative risk used in a type of study (case-control) in which it is impossible to calculate relative risk directly. It is a good estimate of relative risk in cases in which the disease is rare. Thus, if the results of a trial are that the probability of death is 25\% in the control group and 10\% in the experimental group, the odds ratio of survival is {[}0.10 ÷ (1.0 - 0.10){]} ÷ {[}(0.25 ÷ (1.0 - 0.25){]} = 0.33. (Related concepts: absolute risk reduction, number needed to treat and relative risk)

\begin{center}\rule{0.5\linewidth}{0.5pt}\end{center}

razão de verossimilhança
 
1) Em epidemiologia clínica, a razão da comparação de probabilidades de se obter um teste positivo em pessoas com uma doença contra as que não a tem. Nota: a razão de verossimilhança é estimada como a sensibilidade ÷ (1-especificidade). É também possível se estimar a razão de verossimilhança de um resultado negativo: (1-sensibilidade) ÷ especificidade. 2) em modelagem estatística, um indicador da comparação entre a aderência aos dados de dois modelos aninhados.

likehood ratio

\begin{enumerate}
\def\labelenumi{\arabic{enumi})}
\tightlist
\item
  In clinical epidemiology, a ratio comparing the probabilities of obtaining a positive test in persons with the disease with those of doing so in persons without the disease. Note: The likelihood ratio is estimated as sensitivity ÷ (1 minus specificity). It is also possible to estimate the likelihood ratio of a negative result: (1 minus sensitivity) ÷ specificity. 2) In statistical modelling, an indicator comparing the fit of two nested models to the data.
\end{enumerate}

\begin{center}\rule{0.5\linewidth}{0.5pt}\end{center}

reação adversa

veja efeito adverso

adverse reaction

see adverse effect

\begin{center}\rule{0.5\linewidth}{0.5pt}\end{center}

redução absoluta do risco

é a diferença entre a probabilidade que um evento ocorrerá no grupo exposto a um dado fator e a probabilidade que esse evento ocorrerá no grupo não exposto a esse fator.

absolute risk reduction

The value of the difference between the probability that an event will occur in the group exposed to a given factor and the probability that this event will occur in the group not exposed to this factor.

\begin{center}\rule{0.5\linewidth}{0.5pt}\end{center}

redução relativa do risco

A proporção (quociente) da diferença de risco para a ocorrência de um evento no grupo exposto a um dado fator e a probabilidade que esse evento ocorrerá no grupo não exposto a esse fator. Nota: assim, se os resultados de um ensaio mostrarem que a probabilidade de morte é de 25\% no grupo controle e de 10\% no grupo experimental, a redução relativa do risco é (0,25 - 0,10) ÷ 0,25 = 0,6. (conceitos relacionados: redução absoluta do risco, número necessário para tratar, e razão de chances).

\begin{center}\rule{0.5\linewidth}{0.5pt}\end{center}

registro

Veja banco de dados

register

See database.

\begin{center}\rule{0.5\linewidth}{0.5pt}\end{center}

registro

Um arquivo no qual é registrada, continuamente, informação relativa a todos os casos de uma dada doença ou outro problema numa população definida, de modo que os casos possam ser relacionados ao banco de dados referentes a essa população.

registry

A file in which is registered, on an ongoing basis, information concerning all cases of a given disease or other problem in a defined population, such that the cases may be related to the database concerning that population.

\begin{center}\rule{0.5\linewidth}{0.5pt}\end{center}

registros de ensaios

Um banco de dados de citações sobre ensaios randomizados controlados e ensaios clínicos controlados na Cochrane Collaboration. Nota: este banco de dados inclui informação sobre ensaios em andamento e não publicados.

trials register

A database of citations on randomised controlled trials and controlled clinical trials in the Cochrane Collaboration. Note: This database includes information about unpublished and ongoing trials.

\begin{center}\rule{0.5\linewidth}{0.5pt}\end{center}

regressão múltipla

Veja análise de regressão

multiple regression

See regression analysis.

\begin{center}\rule{0.5\linewidth}{0.5pt}\end{center}

regressão ponderada por mínimos quadrados

Uma extensão da regressão de mínimos quadrados usada quando há uma heterogeneidade nas variâncias ou dependência entre as observações. Nota: Este método é frequentemente usado em meta-regressões, e os pesos correspondem então à precisão de cada estimativa de efeito do estudo.

weighted least squares regression

An extension of least squares regression used when there is heterogeneity in the variances or dependence between the observations. Note: This method is often used in meta-regression, and the weights then correspond to the precision of each study's estimate of effect.

\begin{center}\rule{0.5\linewidth}{0.5pt}\end{center}

relatório de ATS

Um relatório que inclui uma revisão compreensiva e sistemática da literatura, ou uma revisão sistemática de evidências de alto nível, avaliando a segurança e a efetividade de uma tecnologia, bem como uma análise do custo-efetividade de uma tecnologia. Nota 1: o relatório sempre: descreve as caracteríticas e o uso atual da tecnologia; avalia criticamente a qualidade da base da evidência; provê informação dos custos do impacto financeiro; e discute considerações organizacionais. Nota 2: o relatório pode opcionalmente endereçar considerações éticas, sociais e legais derivadas do uso da tecnologia. Nota 3: o custo-efetividade da tecnologia pode ser endereçado através de modelagem econômica, quando apropriado. Nota 4: termos relacionados incluem ATS (avaliação de tecnologia em saúde), mini-ATS, revisão rápida, nível de evidência.

HTA report

A report that includes a comprehensive systematic literature review, or a systematic review of high level evidence, evaluating the safety and effectiveness of a technology, as well as an analysis of the cost-effectiveness of the technology. Note 1: The report always: • describes the characteristics and current use of the technology; • critically appraises the quality of the evidence base; • provides information on costs and financial impact; and • discusses organisational considerations. Note 2: The report may optionally address ethical, social and legal considerations arising from use of the technology. Note 3: The cost-effectiveness of the technology may be addressed through economic modelling, when it is appropriate. Note 4: Related terms include health technology assessment, mini-HTA, rapid review, level of evidence.

\begin{center}\rule{0.5\linewidth}{0.5pt}\end{center}

relatório de caso

em pesquisa clínica: um estudo observacional não controlado de uma intervenção ou exposição a um dado fator, e o seu desfecho num único indivíduo. Em pesquisa social: um método empírico que estuda um fenômeno contemporâneo no seu contexto natural, quano a relação entre o fenômeno e seu contexto é complexa, e quando são necessárias várias fontes de dados para o entendimento da relação.

case report

In clinical research: an uncontrolled observational study of an intervention or of exposure to a given factor, and of its outcome in a single individual. In social research: an empirical research method that studies a contemporary phenomenon in its natural context, when the relationship between the phenomenon and its context is complex, and when a number of data sources are needed in order to understand the relationship.

\begin{center}\rule{0.5\linewidth}{0.5pt}\end{center}

relatório de consenso

Recomendações ou diretrizes práticas produzidas através de consenso especializado

consensus report

The recommendations or practice guideline produced through expert consensus.

\begin{center}\rule{0.5\linewidth}{0.5pt}\end{center}

representação esquemática causal

uma representação, esquemática, por exemplo, de relações diretas e indiretas entre intervenções e desfechos. Nota: para um problema clínico, este método tipicamente inclui um elemento populacional, uma ou mais modalidades de intervenção (Por exemplo: triagens, diagnósticos ou tratamentos), desfechos intermediários (i.e.~marcadores biológicos) e desfechos em saúde. Representações causais são feitas para clarificar e explicitar as questões a serem respondidas por uma avaliação, e são úteis na identificação de relações cruciais para os quais dados científicos possam esta faltando.

causal pathway

A depiction, in a schematic, for example, of the direct and indirect linkages between interventions and outcomes. Note: For a clinical problem, this method typically includes a subject population, one or more intervention modalities (e.g.~screening, diagnoses or treatment), intermediate outcomes (e.g.~biological markers) and health outcomes. Causal pathways are intended to clearly and explicitly define the questions to be addressed in an assessment, and they are useful in identifying pivotal linkages for which scientific data may be lacking.

\begin{center}\rule{0.5\linewidth}{0.5pt}\end{center}

representação esquemática clínica

uma representação esquemática do cuidado clínico sistemático de um paciente que propõe, para um período de tratamento relativo a um diagnóstico simples, uma coordenação otimizada na entrega do cuidado proporcionado por médicos e outros profissionais da saúde, desde o momento em que o diagnóstico é confirmado até o momento em que o paciente é admitido, mais o estágio de reabilitação e os cuidados na residência (home care)

clinical pathway

A graphical depiction of the systematic clinical care of a patient that proposes, for a treatment period related to a single diagnosis, an optimal co-ordination of the delivery of care provided by physicians and other health professionals, from when the diagnosis is confirmed to when the patient is admitted, and at the stage of rehabilitation and home care.

\begin{center}\rule{0.5\linewidth}{0.5pt}\end{center}

representante do paciente

veja advogado do consumidor
patient representative

See consumer advocate.

\begin{center}\rule{0.5\linewidth}{0.5pt}\end{center}

resolução

A capacidade de um dispositivo de imagem distinguir dois objetos separados tanto por uma distância física (resolução espacial) ou por sua composição (resolução por contraste).

resolution

The ability of an imaging device to distinguish two objects that are separate in either physical distance (spatial resolution) or in composition (contrast resolution).

\begin{center}\rule{0.5\linewidth}{0.5pt}\end{center}

resumo da característica de operação do receptor (SROC)

Numa revisão sistemática de estudos em testes diagnósticos ou de triagem, uma representação gráfica da relação entre a proporção de verdadeiros positivos (sensibilidade) e os falsos positivos (1 - especificidade), em todos os estudos individuais de um dado teste. Nota: uma linha de regressão pode ser ajustada através desses pontos.

summary receiver operating curve (SROC)

In a systematic review of studies on diagnostic or screening tests, a graphical depiction of the relationship between the true positive ratio (sensitivity) and the false positive ratio (1 - specificity) in all of the individual studies on a given test. Note: A regression line can be fitted through these points.

\begin{center}\rule{0.5\linewidth}{0.5pt}\end{center}

revisão

Uma síntese de informação sobre um tópico em particular. Nota 1: uma revisão pode ou não ser compreensiva ou tomar outras medidas para reduzir viés, como no caso de revisões sistemáticas. Nota 2: termos relacionados incluem revisão sistemática; revisão narrativa; revisão de escopo; revisão rápida; meta análise; revisão Cochrane; revisão por pares. Nota 3: sinônimos incluem revisão de literatura; visão geral.

review

A synthesis of information on a particular topic. Note 1: A review may or may not be comprehensive or take other measures to reduce bias, as is the case with systematic reviews. Note 2: Related terms include systematic review; narrative review; scoping review; rapid review; meta-analysis; Cochrane review; peer review. Note 3: Synonyms include literature review; overview.

\begin{center}\rule{0.5\linewidth}{0.5pt}\end{center}

Revisão Cochrane

resumo metodológico atualizado de evidências científicas confiáveis dos riscos e benefícios associados com vários tipos de intervenção em cuidados. Nota: as revisões sistemáticas da Cochrane pretendem facilitar a tomada de decisões em matérias da saúde. Para ser denominada uma ``revisão Cochrane'', a revisão deve estar no banco de dados ``pai'' mantido pela colaboração Cochrane. Este banco de dados é composto de módulos de revisões sistemáticas. As revisões são colocadas num dos módulos segundo a recomendação da equipe editorial de um dos grupos de revisão colaborativa (CRGs) registrado na Colaboração Cochrane. Os autores da revisão sistemática devem aderir às diretrizes publicadas no Manual de Revisores Cochrane. A metodologia da revisão sistemática Cochrane é preparada usando o software Review Manager (RevMan) produzido pela Colaboração e aderente ao formato da estrutura definido no Manual.

Cochrane review

A methodical and up-to-date summary of reliable scientific evidence of the risks and benefits associated with various care interventions. Note: Systematic Cochrane reviews are intended to facilitate decision making in health matters. To be called a ``Cochrane review,'' a review must be in the Parent Database maintained by the Cochrane Collaboration. This database is composed of modules of systemic reviews. The reviews are placed in one of the modules on the recommendation of the editorial team of one of the collaborative review groups (CRGs) registered with the Cochrane Collaboration. The authors of the systemic reviews must adhere to the guidelines published in the Cochrane Reviewers' Handbook. The methodology of a systematic Cochrane review is prepared using the Review Manager (RevMan) software produced by the Collaboration and adheres to a structured format defined in the Handbook.

\begin{center}\rule{0.5\linewidth}{0.5pt}\end{center}

revisão de literatura

Um sumário, com ou sem interpretação, dos resultados de uma pesquisa reportados na literatura científica. Nota: uma revisão de literatura pode incluir pesquisas qualitativas, bem como várias análises sistemáticas e quantitativas, tais como meta-análises.

literature review

A summary, with or without interpretation, of the research findings reported in the scientific literature. Note: A literature review may include qualitative research, as well as various systematic and quantitative analyses, such as meta-analyses.

\begin{center}\rule{0.5\linewidth}{0.5pt}\end{center}

revisão de narrativa

Uma visão geral de estudos originais que não tenham sido identificados ou analisados de uma forma sistemática (ou seja, padronizada e objetiva).

narrative review

An overview of original studies which have not been identified or analysed in a systematic (i.e.~standardised and objective) way.

\begin{center}\rule{0.5\linewidth}{0.5pt}\end{center}

revisão rápida

Um relatório ue usualmente inclui uma revisão do nível de evidência mais alto ou de evidência recente, e que pode restringir a literatura em uma ou duas bases de dados. Nota 1: o relatório sempre: • descreve as características e o uso atual da tecnologia e • avalia as questões de segurança e efetividade. Nota 2: o relatório nem sempre criticamente avalia a qualidade da base da evidência ou provê informação sobre custos e impacto financeiro. Nota 3: uma revisão rápida tipicamente não é tão rigorosa quanto um relatório mini-HTA ou um relatório HTA mas é de confecção mais rápida. Nota 4: termos relacionados incluem avaliação de tecnologia em saúde, relatório HTA, mini-HTA, nível de evidência.

rapid review

A report that usually includes a review of the highest level of evidence or of recent evidence and that may restrict the literature to one or two databases. Note 1: The report always:
• describes the characteristics and current use of the technology, and
• evaluates safety and effectiveness issues Note 2: The report does not always critically appraise the quality of the evidence base or provide information on costs and financial impact. Note 3: A rapid review is typically not as rigorous as a mini-HTA or a HTA report but is quicker to produce. Note 4: Related terms include health technology assessment, HTA report, mini-HTA, level of evidence.

\begin{center}\rule{0.5\linewidth}{0.5pt}\end{center}

revisão sistemática

Uma síntese que coleta todas as evidências aderentes ao critério de elegibilidade pré-especificado de modo a responder a uma questão específica de pesquisa. Nota 1: revisões sistemáticas são conduzidas de acordo com um protocolo pré-especificado. Os métodos usados são selecionados com o propósito de reduzir viéses, promovendo assim descobertas mais confiáveis das quais as conclusões poderão ser obtidas e as decisões tomadas. Nota 2: muitas revisões sistemáticas contem meta análises. Uma meta análise consiste no uso de métodos estatísticos para resumir os resultados de estudos independentes.

systematic review

A synthesis that collates all empirical evidence fitting pre-specified eligibility criteria in order to answer a specific research question. Note 1: Systematic reviews are conducted according to a pre-specified protocol. The methods used are selected with a view to minimizing bias, thus providing more reliable findings from which conclusions can be drawn and decisions made. Note 2: Many systematic reviews contain meta-analyses. A meta-analysis is the use of statistical methods to summarize the results of independent studies.

\begin{center}\rule{0.5\linewidth}{0.5pt}\end{center}

revisor

No contexto da revisão Cochranem uma pessoa responsável por preparar ou atualizar a revisão.

reviewer

In the context of a Cochrane review, a person responsible for preparing or updating the review.

\begin{center}\rule{0.5\linewidth}{0.5pt}\end{center}

risco

A probabilidade de que um evento, geralmente adverso, possa ocorrer. Nota: esta probabilidade pode ser estimada como a proporção de pessoas num grupo com as quais um evento ocorrer. (conceito relacionado: fator de risco).

risk

The probability that an event, generally adverse, will occur. Note: This probability can be estimated as the proportion of persons in a group for whom an event occurs. (Related concept: risk factor).

\begin{center}\rule{0.5\linewidth}{0.5pt}\end{center}

risco relativo

A proporção (quociente) do risco que um evento ocorrerá entre os indivíduos expostos a um dado fator e o risco de que esse evento ocorrerá entre os indivíduos não expostos a esse fator. Nota: o risco relativo (RR) = 1 indica que o risco é igual nos grupos comparados; RR \textgreater{} 1 indica que o fator aumenta o risco e RR\textless1 indica que o fator diminui o risco. Sinônimo: proporção de risco.

risco

A probabilidade de que um evento, geralmente adverso, possa ocorrer. Nota: esta probabilidade pode ser estimada como a proporção de pessoas num grupo com as quais um evento ocorrer. (conceito relacionado: fator de risco).

\begin{center}\rule{0.5\linewidth}{0.5pt}\end{center}

roteiro do paciente

veja roteiro clínico

patient pathway

See clinical pathway.

\hypertarget{s}{%
\chapter*{S}\label{s}}
\addcontentsline{toc}{chapter}{S}

seguimento

Uma observação, durante um dado período, de uma pessoa, grupo ou população definida no início, de certas características daqueles ou daquilo que está sendo avaliado, para identificar as alterações no seu estado de saúde ou em certas variáveis relacionadas à saúde. Nota: se dados dos efeitos nos sujeitos da intervenção sendo estudados forem perdidos (devido aos sujeitos terem, se mudado ou saído por alguma razão), os resultados do estudo podem ser influenciados, especialmente se alguns tipos de sujeitos forem sistematicamente perdidos do estudo pelas mesmas razões. Os investigadores devem reportar o número e tipo de sujeitos que não puderam ser avaliados, para que a possibilidade de viés possa ser considerada. Conceito relacionado: análise de intenção de tratamento.

follow-up

An observation, over a given period, of a person, group or population defined at the outset, certain characteristics of whom or of which are evaluated to identify the changes in their health status or in certain health-related variables. Note: If data on the effects on subjects of the intervention being studied are lost (because subjects have moved or have dropped out for some other reason), the study's results may be influenced, especially if some types of subjects are systematically lost from the study for the same reasons. The investigators must report the number and type of subjects who could not be evaluated, in order that the possibility of bias may be taken into consideration. Related concept: intention-to-treat analysis

\begin{center}\rule{0.5\linewidth}{0.5pt}\end{center}

segurança

Um julgamento concernente à aceitação do risco (uma medida da probabilidade de um desfecho adverso e sua severidade) associado ao uso de uma tecnologia numa dada situação (por exemplo para um paciente com um particular problema de saúde) por um clínico com certo treinamento, ou na definição de um tratamento especificado.

safety

A judgment concerning the acceptability of the risk (a measure of the probability of an adverse outcome and its severity) associated with using a technology in a given situation (e.g.~for a patient with a particular health problem) by a clinician with certain training, or in a specified treatment setting.

\begin{center}\rule{0.5\linewidth}{0.5pt}\end{center}

sensibilidade

Uma característica de performance de um teste diagnóstico, definido como a proporção de pessoas que tiveram um resultado positivo no teste e as pessoas que tiveram a doença; é calculada assim: verdadeiros positivos ÷ (verdadeiros positivos + falsos negativos). Conceitos relacionados: acurácia, curva ROC e especificidade.

sensitivity
A characteristic of the performance of a diagnostic test, defined as the proportion of persons who have a positive test result among the persons who have the disease; it is calculated as follows: true positives ÷ (true positives + false negatives). Related concepts: accuracy, ROC curve and specificity.

\begin{center}\rule{0.5\linewidth}{0.5pt}\end{center}

série

Veja série de casos

series

See case series.

\begin{center}\rule{0.5\linewidth}{0.5pt}\end{center}

série temporal interrompida

Um estudo do efeito terapêutico de uma tecnologia através da comparação dos resultados de medições múltiplas feitas antes e depois da tecnologia ser utilizada num dado grupo de indivíduos. Nota: o valor desse tipo de estudo pode ser fortalecido pela adição de um grupo de controle no qual a tecnologia ainda não tenha sido tentada e pela observação desse grupo nos mesmos tempos. Além disso, séries temporais múltiplas podem ser estudadas pelo escalonamento na introdução da tecnologia.

interrupted time series

A study of the therapeutic effect of a technology through comparison of the results of multiple measurements taken before and after the technology is used in a given group of subjects. Note: The value of this type of study can be strengthened by the addition of a control group in which the technology has not been tried and which is observed at the same points in time. Also, multiple time series can be studied by staggering the introduction of the technology.

\begin{center}\rule{0.5\linewidth}{0.5pt}\end{center}

séries de casos

estudo observacional não controlado de uma intervenção ou exposição a um dado fator, e seu desfecho numa série de elementos

case series

An uncontrolled observational study of an intervention or of exposure to a given factor, and of its outcome in a series of subjects.

\begin{center}\rule{0.5\linewidth}{0.5pt}\end{center}

séries temporais

Uma série de medidas das mesmas variáveis realizada ao longo do tempo. Nota: uma série temporal é dita a ser interrompida quando um conjunto de medidas, feito antes de uma intervenção, é realizado ou uma alteração é feita num dado sistema, e outras medidas são feitas ao longo do tempo após a alteração.

time series

A series of measurements of the same variables taken over time. Note: A time series is said to be interrupted when a set of measurements is taken before an intervention is carried out or a change is made in a given system, and other measurements are taken over time after the change.

\begin{center}\rule{0.5\linewidth}{0.5pt}\end{center}

significância clínica

diferença no tamanho do efeito que os especialistas consideram significante para decisões clínicas ou estratégicas, independentemente de seu significado estatístico. Nota: num ensaio clínico com grande número de elementos, uma pequena diferença entre o grupo de intervençaõ e o grupo de controle, pode ser estatisticamente significante, mas clinicamente não importante. Num ensaio com um número pequeno de elementos, uma diferença clinicamente significativa pode ser observada, podendo não ser estatisticamente significante. Um ensaio maior pode ser necessário para confirmar diferenças estatisticamente significantes.

clinical significance

A difference in effect size that the experts consider significant for clinical or strategic decisions, regardless of statistical significance. Note: In a trial with a large number of subjects, a small difference between the treatment group and the control group may be statistically significant, but clinically unimportant. In a trial with a small number of subjects, an important clinical difference may be observed that is not statistically significant. A larger trial may be needed to confirm that the difference is statistically significant.

\begin{center}\rule{0.5\linewidth}{0.5pt}\end{center}

significância estatística

Em teste de hipóteses, uma conclusão obtida quando a hipótese nula é rejeitada, isto é, quando o valor p está abaixo do valor pré-determinado do nível alfa. (Conceito relacionado: valor p).

statistical significance

In hypothesis testing, a conclusion drawn when the null hypothesis is rejected, i.e.~when the p-value is below the pre-determined alpha level. (Related concept: p-value)

\begin{center}\rule{0.5\linewidth}{0.5pt}\end{center}

Simulação de Monte Carlo

A classificação de uma pessoa ou uma característica em uma categoria diferente daquela a qual deveria ter sido atribuída, levando a um viés de medição. Nota: Por exemplo, uma simulação de Monte Carlo pode ser usada para representar ou modelar muitos pacientes individuais em uma população com intervalos de valores para determinadas características de saúde ou resultados clínicos. Em alguns casos, os componentes aleatórios são adicionados aos valores de uma variável de entrada conhecida para determinar os efeitos das flutuações dessa variável nos valores da variável de saída.

Monte Carlo simulation

The classification of a person or a characteristic as being in a category other than that to which it should have been assigned, leading to a measurement bias. Note: For example, a Monte Carlo simulation can be used to represent or model many individual patients in a population with ranges of values for certain health characteristics or clinical outcomes. In some cases, the random components are added to the values of a known input variable to determine the effects of fluctuations of this variable on the values of the output variable.

\begin{center}\rule{0.5\linewidth}{0.5pt}\end{center}

sistema precoce de alerta

Uma unidade estável com conexões e fontes confiáveis, com o objetivo de: identificar novas tecnologias que tenham potencial para provocar grande impacto nos serviços de saúde; filtrar e priorizar essas tecnologias para selecionar aquelas que provavelmente terão impacto na saúde, serviços e orçamentos; e estabelecer esse impacto.

early warning system

A stable unit with reliable connections and sources which aims to: identify new technologies that have the potential to make a large impact on health services; filter and prioritize these technologies to select those most likely to have a impact on health, services and budgets; and assess that impact.

\begin{center}\rule{0.5\linewidth}{0.5pt}\end{center}

sistema precoce de alerta e sistema de alerta

Um sistema que pretende identificar, filtrar e priorizar tecnologias em saúde, novas e emergentes, ou novos usos de intervenções existentes; para estabelecer ou predizer seu impacto na saúde, serviços de saúde e/ou na sociedade; e para disseminar informação. Nota 1: ``filtrar'' é um processo para remover, da lista de tecnologias originadas do processo de identificação, as tecnologias que não são relevantes para o sistema precoce de alerta e para o sistema de alerta. Nota 2: ``priorizar'' é um processo para determinar a significância das tecnologias filtradas ou a ordem com que devem ser analisadas, de acordo com a sua importância relativa para os objetivos do sistema precoce de alerta e do sistema de alerta. Nota 3: sinônimos incluem sistema precoce de aviso; sistema de prospecção do horizonte; sistema EAA.

early awareness and alert system

A system that aims to identify, filter and prioritize new and emerging health technologies, or new uses of existing interventions; to assess or predict their impact on health, health services and/or society; and to disseminate information. Note 1: `Filter' is a process to remove technologies that are not relevant to the early awareness and alert system from the list of technologies originating from the identification process. Note 2: `Prioritize' is a process to determine the significance of, or order for dealing with, filtered technologies according to their relative importance to the aims of the early awareness and alert system. Note 3: Synonyms include early warning system; horizon scanning system; EAA system

\hypertarget{t}{%
\chapter*{T}\label{t}}
\addcontentsline{toc}{chapter}{T}

tabela de contingência

Tabela apresentando uma classificação cruzada de dados categóricos, de modo que essas categorias de uma característica seja, indicadas horizontalmente (em linhas) e as categorias da outra característica sejam indicadas verticalmente (em colunas). Nota: cada célula da tabela indicará o número ou proporção de indivíduos que tenham ambas as características mencionadas na linha e na coluna. A tabela de contingência mais simples é a quádrupla, ou tabela 2x2, que é usada em ensaios clínicos, por exemplo, para comparar desfechos dicotômicos em dois grupos.

contingency table

A table presenting a cross-classification of categorical data such that the categories of one characteristic are indicated horizontally (in rows) and the categories of another characteristic are indicated vertically (in columns). Note: Each cell of the table will indicate the number or proportion of individuals who have both the characteristic on its row and in its column. The simplest contingency table is the fourfold, or 2x2 table, which is used in clinical trials, for example, to compare dichotomous outcomes for two groups.

\begin{center}\rule{0.5\linewidth}{0.5pt}\end{center}

tabela de evidências

Uma tabela sumarizando características de estudos de uma particular intervenção ou problema de saúde. Nota: este tipo de tabela inclui caracteristicas tais como o plano metodológico, os sujeitos e os desfechos.

evidence table

A summary table of selected characteristics of the studies of a particular intervention or health problem. Note: This type of table includes characteristics such as the methodological plan, the subjects and the outcomes.

\begin{center}\rule{0.5\linewidth}{0.5pt}\end{center}

tamanho da amostra

O número de indivíduos estudados, incluindo os do grupo de intervenção e os do grupo controle, quando aplicável. Nota: em geral, um tamanho maior de amostra diminui a probabilidade de erros Tipo I (α) (falso positivo) e aumenta o poder estatístico do ensaio, isto é, diminui a probabilidade de erros Tipo II (β) (falso-negativo). Um maior tamanho de amostra diminui o efeito da variação randômica sobre a estimativa do efeito do tratamento. No projeto de um estudo, o tamanho desejado da amostra pode ser calculado usando-se uma fórmula estatística baseada nos níveis aceitáveis de erros α e β, na menor diferença entre os grupos de intervenção considerados clinicamente relevantes e na medida da variância. Conceitos relacionados: alfa e beta.

sample size

The number of subjects studied, including those in the intervention group and the control group, where applicable. Note: In general, a large sample size decreases the probability of Type I (α) (false-positive) errors and increases the statistical power of the trial, i.e.~decreases the probability of Type II (β) (false-negative) errors. A large sample size decreases the effect of random variation on the estimate of a treatment effect. In designing a study, the desired sample size can be calculated using a statistical formula based on the acceptable levels of α and β errors, the smallest difference between intervention groups considered clinically relevant, and the measurement variance. Related concepts: alpha and beta.

\begin{center}\rule{0.5\linewidth}{0.5pt}\end{center}

tamanho do efeito

Uma medida adimensional do grau de presença de um fenômeno na população de interesse. Nota: no caso de variáveis contínuas, é geralmente definida como a diferença das médias entre os grupos experimental e o de controle, divididas pelo desvio padrão do grupo de controle ou de ambos os grupos. Quando escalas diferentes (avaliação de dor, por exemplo) são usadas para medir o desfecho, elas se tornam comparáveis. Nota: esta medida é usada para determinar o tamanho de uma amostra.

effect size

A dimensionless measure of the degree of presence of a phenomenon in the population of interest. Note: In the case of continuous variables, it is generally defined as the difference in means between the experimental and control groups, divided by the standard deviation of the control group or both groups. When different scales (pain assessment, for example) are used to measure an outcome, they become comparable. Note: This measure is used to determine the size of a sample.

\begin{center}\rule{0.5\linewidth}{0.5pt}\end{center}

taxa de desconto

A taxa de juros usada para determinar o valor presente de custos e benefícos futuros. Nota: A taxa de desconto (por exemplo 3\% ou 5\%) é usualmente baseada em taxas de empréstimo governamentais (por exemplo, em notas de três meses do Tesouro) ou em taxas de juros de mercado onde o vencimento é equivalente ao período no qual a intervenção ou programa é avaliado.

discount rate

The interest rate used to determine the present value of future costs and benefits. Note: The discount rate (e.g.~3\% or 5\%) is usually based on government borrowing rates (on three-month Treasury bills, for example) or market interest rates where the maturity is equivalent to the period in which the intervention or program is evaluated.

\begin{center}\rule{0.5\linewidth}{0.5pt}\end{center}

taxa de evento

A proporção de membros de um grupo no qual um evento é observado duarante um peródo de tempo especificado. Nota: se um evento (por exemplo, um acidente vascular) é observado em 32 pacientes de um grupo de 100 sendo monitorados, por um ano, a taxa de acidente vascular é de 0,31 num ano.

event rate

The proportion of the members of a group in whom an event is observed over a specified period of time. Note: If an event (e.g.~a stroke) is observed in 32 subjects out of a group of 100 monitored over one year, the stroke rate is 0.32 a year.

\begin{center}\rule{0.5\linewidth}{0.5pt}\end{center}

taxa de mortalidade

A proporção de óbitos em uma determinada população dentro de um período específico (geralmente um ano). Nota: A taxa é frequentemente expressa como um número por 100.000, para facilitar a interpretação (por exemplo, 18.3 mortes por ano por 100.000 pessoas).

mortality rate

The proportion of deaths in a given population within a specified period (usually one year). Note: The rate is often expressed as a number per 100,000, to facilitate interpretation (e.g.~18.3 deaths per year per 100,000 persons).

\begin{center}\rule{0.5\linewidth}{0.5pt}\end{center}

técnica de grupo nominal

Uma técnica estruturada para coletar informação de peritos (de 9 a 12 pessoas), consistindo na avaliação de questões e uma discussão. Nota: a técnica de grupo nominal é composta dos seguintes passos: 1. Antes da reunião, uma revisão da literatura é enviada aos peritos; 2. Os peritos investem alguns minutos para escrever seus comentários sobre as questões sendo estudadas; 3. Cada participante contribui com uma idéia, a qual é registrada pelo moderador num quadro; 4. Sugestões similares são então agrupadas; 5. Cada participante isoladamente atribui uma nota a cada idéia; 6. As notas individuais são tabuladas, e a tabulação é apresentada aos participantes; 7. A posição geral é discutida e é feito novo ranqueamento; 8. Os resultados são apresentados aos participantes para aprovação. (Conceito relacionado: técnica Delphi).

nominal group technique

A structured technique for gathering information from experts (9 to 12 people), consisting of an evaluation of questions and a discussion. Note: The nominal group technique is composed of the following steps:
1. Before the meeting, a literature review is sent to the experts.
2. The experts take a few minutes to write down their comments on the questions being studied.
3. In turn, each participant contributes an idea, which the moderator records on a flip chart.
4. Similar suggestions are then grouped together.
5. Each participant privately ranks each idea.
6. The individual ranking is tabulated, and the tabulation is presented to the participants.
7. The overall ranking is discussed and re-ranked.
8. The results are presented to the participants for approval. (Related concept: Delphi technique)

\begin{center}\rule{0.5\linewidth}{0.5pt}\end{center}

tecnologia em saúde emergente

Uma tecnologia que ainda não foi adotada pelo sistema de saúde. Nota 1: Drogas estão na fase II ou III do ensaio clínico, ou num estágio pré-lançamento; dispositivos médicos estão no estágio pré-marketing. Nota 2: termos relacionados incluem nova tecnologia em saúde.

emerging health technology

A health technology that has not yet been adopted within the healthcare system. Note 1: Pharmaceuticals are in the Phase II or III clinical trial, or pre-launch stage; medical devices are in the pre-marketing stage. Note 2: Related terms include new health technology.

\begin{center}\rule{0.5\linewidth}{0.5pt}\end{center}

tecnologia

A aplicação de conhecimento científico ou outro, com relação a ferramentas, técnicas, produtos, processos, métodos, modos de organização,ou sistemas - para tarefas práticas. Nota: em cuidados com a saúde, a tecnologia engloba: drogas; testes diagnósticos, incluindo indicadores e reagentes; dispositivos; equipamentos e suprimentos; procedimentos médicos e cirúrgicos; sistemas de suporte; sistemas organizacionais e de gestão utilizados na prevenção, rastreamento, diagnóstico, tratamento e reabilitação. (conceito relacionado: tecnologia em saúde).

technology

The application of scientific or other organized knowledge -- relating to tools, techniques, products, processes, methods, modes of organisation, or systems -- to practical tasks. Note: In health care, technology encompasses: drugs; diagnostic tests, including indicators and reagents; devices; equipment and supplies; medical and surgical procedures; support systems; and organizational and managerial systems used in prevention, screening, diagnosis, treatment and rehabilitation. (Related concept:health technology).

\begin{center}\rule{0.5\linewidth}{0.5pt}\end{center}

tecnologia em saúde

tecnologia em saúde é definida como uma intervenção que pode ser usada para promover a saúde, prevenir, diagnosticar ou tratar males agudos ou crônicos, ou para a reabilitação. Tecnologias em saúde incluem drogas, dispositivos, procedimentos e sistemas de saúde utilizados nos cuidados da saúde.

health technology

health technology is defined as an intervention that may be used to promote health, to prevent, diagnose or treat acute or chronic disease, or for rehabilitation. Health technologies include pharmaceuticals, devices, procedures and organizational systems used in health care.

\begin{center}\rule{0.5\linewidth}{0.5pt}\end{center}

tecnologia em saúde - obsoleta

Uma tecnologia em saúde que não seja mais o padrão de cuidados ou cujo benefício clínico, segurança ou eficiência em custo tenha sido superada por tecnologias alternativas disponíveis. Nota 1: termos relacionados incluem desinvestimento

obsolete health technology

A health technology that is no longer the standard of care or whose clinical benefit, safety or cost-effectiveness has been superseded by available alternative technologies. Note 1: Related terms include disinvestment

\begin{center}\rule{0.5\linewidth}{0.5pt}\end{center}

teoria fundamentada

Uma teoria resultante da coleta sistemática e da análise de informação sobre o processo de pesquisa com o objetivo de construir uma teoria de base empírica em relação a fenômenos sociais nos quais pouca análise tenha sido feita.

grounded theory

A theory resulting from the systematic gathering and analysis of information through the research process, with the objective of building empirically based theories in relation to social phenomena regarding which little analysis has been done.

\begin{center}\rule{0.5\linewidth}{0.5pt}\end{center}

teste de hipóteses

Um método de inferência estatística para avaliação da plausibilidade da hipótese nula à luz dos dados observados. Nota: a hipótese nula é assumida como verdadeira no início, mas, se a mesma provar estar errada quando os dados observados forem observados, ela será rejeitada em favor da hipótese alternativa (negativa para a hipótese nula).

hypothesis testing

A method of statistical inference for evaluating the plausibility of the null hypothesis in the light of the observed data. Note: The null hypothesis is assumed to be true at the outset, but, if that assumption proves wrong when the observed data are examined, it is rejected in favour of the alternative hypothesis (negative of the null hypothesis).

\begin{center}\rule{0.5\linewidth}{0.5pt}\end{center}

teste de Mante-Haentzel

Um teste qui-quadrado para levar em conta variávei categóricas confundidoras, ao se fazer a soma ponderada das associações em cada uma das categorias, denominadas estratos (strata). Nota: em meta-análise, os estratos consistem nos vários estudos incluidos.

Mante-Haentzel test

A chi-square test for taking categorical confounding variables into account when doing a weighted sum of the associations in each of the categories, called strata. Note: In meta-analyses, the strata are the various studies included.

\begin{center}\rule{0.5\linewidth}{0.5pt}\end{center}

teste estatístico

Veja teste de hipóteses.

statistical test

See hypothesis testing.

\begin{center}\rule{0.5\linewidth}{0.5pt}\end{center}

teste múltiplo

Diversos testes estatísticos realizados sobre as mesmas observações. Nota: num estudo, o tamanho da amostra e o nível alfa são estabelecidos de acordo com o princípio de que uma hipótese nula única será testada ao final do estudo, uma vez que os dados tenham sido colhidos. Se mais do que um teste estatístico for realizado, isto é, se forem feitas comparações de várias intervenções, ou testes em variáveis distintas, em momentos distintos ou em sub-grupos diferentes, haverá um aumentoda probabilidade de ocorrência do erro de Tipo 1 (alfa) planejado. Alguns métodos estatísticos se propõe a considerar testes múltiplos, mas são controversos.

multiple testing

Several statistical tests done on the same observations. Note: In a study, the sample size and the alpha level are set according to the principle that a single null hypothesis will be tested at the end of the study, once all the data have been gathered. If more than one statistical test is undertaken, e.g.~pairwise comparisons of several interventions, or tests on different variables, at different timepoints or on different sub-groups, there will be an increase in the planned overall Type I error (alpha) probability. Some statistical methods are proposed to take multiple testing into account, but they are controversial.

\begin{center}\rule{0.5\linewidth}{0.5pt}\end{center}

tomada de decisão baseada em evidências

A consideração dos achados de pesquisas médicas e científicas, no processo de tomada de decisão. Sinônimo: tomada de decisão por evidência informada.

evidence-based decision making

The consideration of the findings from scientific and medical research in the decision-making process. Syn.: evidence-informed decision making

\begin{center}\rule{0.5\linewidth}{0.5pt}\end{center}

triagem

Um serviço de saúde pública no qual os membros de uma população de risco de uma doença ou de sua complicação é submetido a perguntas ou oferecido um teste, para identificar quais indivíduos estão num risco suficientemente alto que justifique testes diagnósticos ou tratamentos adicionais.

screening

A public health service in which members of a population at risk for a disease or its complications are asked questions or offered a test, to identify those individuals who are at sufficiently high risk that further diagnostic tests or treatments are justified.

\hypertarget{b-1}{%
\chapter*{B}\label{b-1}}
\addcontentsline{toc}{chapter}{B}

banco de dados bibliográfico

fonte computadorizada e indexada de citações descrevendo documentos tais como artigos de jornais e relatórios científicos. Nota: as citações tipicamente incluem o(s) nome(s) do autor(es), o titulo, a fonte, um resumo e, quando aplicavel, informações relacionadas (incluindo o texto completo em alguns casos). Exemplos: MEDLINE e EMBASE.

bibliographic database

An indexed computerized or printed source of citations, describing documents such as journal articles and scientific reports. Note: The citations typically include the name(s) of the author(s), the title, the source, an abstract and, where applicable, related information (including the full text in some cases). Example: MEDLINE and EMBASE.

\begin{center}\rule{0.5\linewidth}{0.5pt}\end{center}

base de dados de ATS

Uma base de dados contendo relatórios sobre avaliações de tecnologia na saúde em desenvolvimento e completadas. Nota: essa base de dados é gerenciada pelo CRD (Centre for Reviews and Dissemination - centro de revisões e disseminação) em colaboração com o secretariado da INAHTA. Está disponível na Livraria Cochrane no site do CRD (\url{http://www.york.ac.uk/inst/crd/}).

HTA database

A database containing the reports on ongoing and completed health technology assessments. Note: This database is managed by the Centre for Reviews and Dissemination (CRD) in collaboration with the INAHTA Secretariat. It is available in the Cochrane Library and on the CRD site (\url{http://www.york.ac.uk/inst/crd/}).

\begin{center}\rule{0.5\linewidth}{0.5pt}\end{center}

base de dados factual

Uma fonte indexada e computadorizada ou impressa que provê acesso direto a informações tais como guias para diagnóstico e tratamento, estatísticas e diretórios de instituições. Outra definição: uma coleção extensa e considerada completa, de fatos e dados sobre materiais num campo especializado, disponíveis para análise e aplicação. Nota: a coleção pode ser computadorizada através de vários métodos modernos de pesquisa bibliográfica. Este conceito difere do banco de dados bibliográfico, que se limita a coleções de citações.

factual database

An indexed computerized or printed source that provides direct access to information such as guidelines for diagnosis and treatment, statistics, and directories of institutions. Other definition: An extensive collection, deemed complete, of facts and data from materials in a specialized field, made available for analysis and application. Note: The collection may be computerized through various modern literature search methods. This concept differs from that of a bibliographic database, which is limited to collections of citations.

\begin{center}\rule{0.5\linewidth}{0.5pt}\end{center}

base do estudo

Uma população estudada numa particular localização por um período de tempo especificado.

study base

A population studied in a particular location over a specified period of time.

\begin{center}\rule{0.5\linewidth}{0.5pt}\end{center}

benefício líquido

O benefício de uma opção em unidades monetárias, menos o seu custo total em unidades monetárias. Nota: o benefício líquido é um critério básico de decisão na análise de custo-benefício (ACB).

net benefit

The benefit of an option in monetary units, minus its total cost in monetary units. Note: Net benefit is a basic decision criterion in cost-benefit analysis (CBA).

\begin{center}\rule{0.5\linewidth}{0.5pt}\end{center}

benefício marginal

O benefício adicional produzido pelo uso de recursos adicionais, expressos, por exemplo, em unidades de desfecho em saúde.

marginal benefit

The additional benefit produced by the use of additional resources, expressed, for example, in units of health outcome.

\begin{center}\rule{0.5\linewidth}{0.5pt}\end{center}

beta

em teste de hipóteses, a probabilidade de erro do Tipo II, i.e.~a probabilidade de se concluir incorretamente que a hipótese nula é verdadeira. Nota: por exemplo, beta poderia ser a probabilidade de se concluir que uma intervenção não é efetiva quando ela tem um efeito verdadeiro. (1-beta) é o poder estatístico de um teste permitindo a rejeição de uma hipótese nula que seja verdadeiramente falsa (i.e.~detectando o efeito de uma intervenção que relamente existe)

beta

In hypothesis testing, the probability of a Type II error, i.e.~the probability of concluding incorrectly that a null hypothesis is true. Note: For example, β could be the probability of concluding that an intervention is not effective when it has a true effect. (1-β) is the statistical power of a test allowing for rejection of a null hypothesis that is truly false (e.g.~detecting the effect of an intervention that truly exists).

\begin{center}\rule{0.5\linewidth}{0.5pt}\end{center}

Biblioteca Cochrane

conjunto de bases de dados atualizadas trimestralmente, projetadas para prover informação científica e evidências para suportar a tomada de decisões no campo da saúde, publicadas em CD-ROM e na Internet. Nota: as bases de dados da Biblioteca Cochrane são: o banco de dados Cochrane de revisões sistemáticas; a central de registro de ensaios controlados da Cochrane; o banco de dados de resumos de revisões de efeitos da Cochrane; o banco de dados de Avaliação de Tecnologias em Saúde e o banco de dados NHS de avaliação econômica.

Cochrane library

A collection of databases updated quarterly, designed to provide scientific information and evidence to support decision making in the health field, published on CD-ROM and on the Internet. Note: The Cochrane Library databases are: the Cochrane Database of Systematic Reviews, the Cochrane Central Register of Controlled Trials, the Database of Abstracts of Reviews of Effects, the Cochrane Methodology Register, the Health Technology Assessment database and the NHS Economic Evaluation database.

\begin{center}\rule{0.5\linewidth}{0.5pt}\end{center}

Bioética

campo de investigação compreendendo a análise de questões éticas e dilemas morais pertencentes ao estudo, classificação e manipulação de vida biológica, especialmente na saúde humana.

bioethics

A field of inquiry comprising the analysis of ethical issues and moral dilemmas pertaining to the study, classification, and manipulation of biological life, especially human health.

\hypertarget{u}{%
\chapter*{U}\label{u}}
\addcontentsline{toc}{chapter}{U}

unidade de alocação

Num ensaio, a unidade atribuida a uma das intervenções de interesse. Nota: a unidade escolhida pode ser uma pessoa, mas, em alguns ensaios, um grupo pode ser escolhido para reduzir o risco de contaminação ou por conveniência; uma unidade também pode ser um hospital ou comunidade. Em outros ensaios, diferentes partes de uma pessoa (tais como o olho esquerdo ou direito) podem ser estabelecidas para receber intervenções diferentes.

unit of allocation

In a trial, the unit assigned to one of the interventions of interest. Note: The unit chosen may be a person, but, in some trials, a group may be chosen to reduce the risk of contamination or for convenience; a unit may also be a hospital or community. In other trials, different parts of a person (such as the left or right eye) may be assigned to receive different interventions.

\begin{center}\rule{0.5\linewidth}{0.5pt}\end{center}

usuário

Uma pessoa que está sob cuidados médicos ou tratamento. (O mesmo que consumidor e paciente)

user

See consumer and patient.

\begin{center}\rule{0.5\linewidth}{0.5pt}\end{center}

usuário de revisões

Um paciente, profissional da saúde ou tomador de decisão usando uma revisão de literatura para tomar decisões práticas sobre cuidados, ou um pesquisador conduzindo ou considerando uma nova pesquisa.

user of reviews

A patient, health professional or decision maker using a literature review to make practical decisions about care, or a researcher conducting or considering new research.

\begin{center}\rule{0.5\linewidth}{0.5pt}\end{center}

utilidade

\begin{enumerate}
\def\labelenumi{\arabic{enumi})}
\tightlist
\item
  Em análise econômica e em análise de decisão, o desejo de um específico estado de saúde ou desfecho em saúde, usualmente expresso como sendo um contínuo de zero a um (morte tendo um valor de utilidade zero e uma vida totalmente saudável tendo um valor de utilidade igual a um). Nota: este termo é frequentemente usado como sinônimo de preferência. 2) O desejo relativo ou a preferência para que ocorra um desfecho em saúde específico, ou um estado de saúde (usualmente pela perspectiva do paciente).
\end{enumerate}

utility

\begin{enumerate}
\def\labelenumi{\arabic{enumi})}
\tightlist
\item
  In economic and decision analysis, the desirability of a specific health status or health outcome, usually expressed as being on a continuum from zero to one (death having a utility value of zero, and a full healthy life having a utility value of one). Note: This term is often used as a synonym for preference. 2) The relative desirability of, or preference for, a specific health outcome or health status (usually from the perspective of the patient).
\end{enumerate}

\hypertarget{v}{%
\chapter*{V}\label{v}}
\addcontentsline{toc}{chapter}{V}

validade

A habilidade do resultado de uma medida ou estudo ser e livre de viéses (erros sitemáticos).

validity

The ability of the result of a measurement or study to be accurate and free of bias (systematic errors). Note: In the context of a study, two types of validity are recognized: internal validity and external validity. In the context of measurement, there are several types of validity: construct validity, content validity, face validity, etc. Related concept: methodological quality

\begin{center}\rule{0.5\linewidth}{0.5pt}\end{center}

validade do estudo

A capacidade de um estudo em permitir conclusões corretas quando forem levados em conta: os métodos de pesquisa, a representatividade da validade interna da amostra do estudo e a natureza da população da qual foi feito o recrutamento com validade externa.

study validity

The ability of a study to allow for the drawing of correct conclusions when the research methods, the representativeness of the study sample internal validity and the nature of the population from which it is recruited external validity are taken into account.

\begin{center}\rule{0.5\linewidth}{0.5pt}\end{center}

validade externa

A habilidade de um projeto de pesquisa prover descobertas que possam ser generalizadas para outras populações, contextos e períodos. Sinônimo: generalizabilidade. Conceito relacionado: aplicabilidade.

external validity

The ability of a research design to provide findings that can be generalised to other populations, contexts and periods. Syn.: generalizability. Related concept: applicability

\begin{center}\rule{0.5\linewidth}{0.5pt}\end{center}

validade interna

A habilidade de um projeto de pesquisa de representar a relação verdadeiramente causal entre uma intervenção e um efeito em circunstâncias particulares da pesquisa. Nota: se a organização interna de um estudo conduz a uma distorção na medida da estimativa, tal medida perde sua validade interna. Isso é o que ocorre quando, por exemplo, uma pessoa é pesada com uma escala defeituosa. Com uma escala bem calibrada um estudo etiológico deve estar livre de defeitos internos devidos à distorção nas medidas. Esses defeitos podem ser encontrados em três estágios do estudo: na seleção do indivíduo, na coleta de informações e no estágio de análise.

internal validity

The ability of a research design to represent the true causal relationship between an intervention and an effect in the particular circumstances of the research.Note: If the internal organization of a study leads to distortion in the estimation measurement, that measurement lacks internal validity. This is what happens when, for example, a person is weighed using a defective scale. As with a well-calibrated scale, an etiological study must be free of measurement-distorting internal defects. These defects may be encountered at three stages in the study: the subject selection, information-gathering and analysis stages.

\begin{center}\rule{0.5\linewidth}{0.5pt}\end{center}

valor

Uma medida cardinal da preferência por um estado de saúde específico ou por um desfecho de saude, medido sob certa incerteza. Nota: esta variável também é usada para dar uma descrição geral e balanceada dos (custos) aspectos positivos e negativos de uma situação (por exemplo: o melhor valor para o investimento). (conceitos relacionados: utilidade e preferencia). Nota: de acordo com Drummond, valor é uma preferência livre de riscos (situação de certeza) e utilidade é a preferência medida sob incerteza.

value

A cardinal measure of the preference for a specific health status or health outcome, measured under uncertainty. Note: This variable is also used to give an overall and balanced description of the positive and negative (costs) aspects of a situation (e.g.~best value for money). (Related concepts: utility and preference.) Note: According to Drummond, value is risk-free preference (situation of certainty) and utility is preference measured under uncertainty.

\begin{center}\rule{0.5\linewidth}{0.5pt}\end{center}

valor atípico

Um valor observado que por ser tão diferente dos demais valores observados leva a uma suspeita de que foi cometido um erro ou que ocorreu um acaso.

outlier

An observed value that is so different from the other observation values as to lead one to suspect that an error has been made or that chance has played a part.

\begin{center}\rule{0.5\linewidth}{0.5pt}\end{center}

valor p

Em teste de hipóteses, a probabilidade que um parâmetro a ser testado tenha um valor tão ou mais extremo que o valor que seria observado se a hipótese nula fosse verdadeira. Nota: se o valor de p associado com o teste estatístico for maior ou igual do que o nível alfa que foi determinado (0,01 ou 0,05, por exemplo), isso significa que a associação ou a diferença observada pode ser devida ao acaso e ue a hipótese nula não pode ser rejeitada. No entanto, se o valor for menor do que o nível alfa que foi determinado, a associação ou diferença é estatísticamente significativa e a hipótese nula é rejeitada.

p value

In hypothesis testing, the probability that a parameter to be tested has a value as extreme or more extreme than the value that would be observed if the null hypothesis were true. Note: If the p value associated with the statistical test is equal to or greater than the alpha level that was determined (0.01 or 0.05, for example), this means that the association or difference observed may be due to chance and that the null hypothesis cannot be rejected. However, if the p value is less than the alpha level that was determined, the association or difference is statistically significant and the null hypothesis is rejected. Alternate spelling: p-value

\begin{center}\rule{0.5\linewidth}{0.5pt}\end{center}

valor preditivo negativo (VPN)

Uma característica de desempenho de teste, definida como a proporção de pessoas que não tem a doença entre aquelas que tem um resultado negativo do teste diagnóstico; é calculada da seguinte forma: verdadeiros negativos ÷ (verdadeiros negativos + falsos negativos). Nota: o VPN varia com a prevalência da doença na população de interesse. Conceito relacionado: valor preditivo positivo.

negative predictive value (NPV)

A test performance characteristic, defined as the proportion of persons who do not have the disease among those who have a negative diagnostic test result; it is calculated as follows: true negatives ÷ (true negatives + false negatives). Note: The NPV varies with the prevalence of the disease in the population of interest. Related concept: positive predictive value.

\begin{center}\rule{0.5\linewidth}{0.5pt}\end{center}

valor preditivo positivo (VPP)

Uma característica de um teste de performance, definida como a proporção de pessoas com um resultado positivo num teste diagnóstico que possui a doença; é calculada da seguinte forma: verdadeiros positivos ÷ (verdadeiros positivos + falsos positivos). Nota: o VPP varia com a prevalência da doença na população de interesse. Conceito relacionado: valor preditivo negativo).

positive predictive value (PPV)

A test performance characteristic, defined as the proportion of persons with a positive result in a diagnostic test who have the disease; it is calculated as follows: true positives ÷ (true positives + false positives). Note: The PPV varies with the prevalence of the disease in the population of interest. Related concept: negative predictive value.

\begin{center}\rule{0.5\linewidth}{0.5pt}\end{center}

valoração da contingência

Método para determinação do valor monetário de uma opção ou de um desfecho clínico, por meio de pesquisas nas quais as pessoas sejam expostas a situações e questionadas sobre o preço máximo que aceitariam pagat por essa opção ou por esse desfecho clínico

contingent valuation

A method for determining the monetary value of an option or clinical outcome by means of surveys in which people are presented with situations and asked what the maximum price is that they would be willing to pay for this option or clinical outcome.

\begin{center}\rule{0.5\linewidth}{0.5pt}\end{center}

variância

Uma medida da variação apresentada por um conjunto de observações, definida pela soma dos quadrados dos desvios da média, dividida pelo número de graus de liberdade no conjunto de observações.

variance

A measure of the variation shown by a set of observations, defined by the sum of the squares of deviations from the mean, divided by the number of degrees of freedom in the set of observations.

\begin{center}\rule{0.5\linewidth}{0.5pt}\end{center}

variável

Um fato que pode ter diferentes valores. Nota: os valores medidos num estudo são dados.

variable

A factor that can have different values. Note: The values measured in a study are data.

\begin{center}\rule{0.5\linewidth}{0.5pt}\end{center}

viés

em geral, qualquer fator que distorce a natureza verdadeira de um evento ou observação. Em investigações clínicas, um viés é qualquer fator sistemático distinto daquele da intervenção de interesse que afeta a magnitude (i.e.~tende a aumentar ou diminuir) uma diferença observada no desfecho de grupo em tratamento ou de um grupo de controle. O viés diminui a acurácia (embora não necessáriamente a precisão) de uma observação. Randomização é uma técnica usada para reduzir esse tipo de viés. O viés também se refere a um ponto de vista prejudicado ou parcial que poderia afetar a interpretação de um problema. Cegagem dupla é uma técnica usada para reduzir esse tipo de viés.

bias

In general, any factor that distorts the true nature of an event or observation. In clinical investigations, a bias is any systematic factor other than the intervention of interest that affects the magnitude of (i.e.~tends to increase or decrease) an observed difference in the outcomes of a treatment group and a control group. Bias diminishes the accuracy (though not necessarily the precision) of an observation. Randomization is a technique used to decrease this form of bias. Bias also refers to a prejudiced or partial viewpoint that would affect someone's interpretation of a problem. Double blinding is a technique used to decrease this type of bias.

\begin{center}\rule{0.5\linewidth}{0.5pt}\end{center}

viés de atrito

um viés associado com a saída ou com a exclusão diferencial de objetos de um estudo. Nota: se os objetos sairem por alguma razão, sua exclusão da análise poderá causar uma estimativa exagerada da efetividade da intervenção. A saída diferencial dos objetos pode ser devida a um viés de seleção num estudo experimental comparativo ou num estudo observacional de coorte.

attrition bias

A bias associated with the dropping out or the differential exclusion of subjects from a study. Note: If subjects drop out of the study for some reason, their exclusion from the analysis could lead to an overestimate of the effectiveness of the intervention. The differential dropping out of subjects may be due to a selection bias in a comparative experimental study or cohort observation study.

\begin{center}\rule{0.5\linewidth}{0.5pt}\end{center}

viés de contexto

Um viés de medida devido à influência do contexto em estudo, ou na interpretação dos resultados do estudo. Nota: por exemplo, quando testes são feitos em grupos com alta prevalência de uma doença, os leitores tendem a interpretar os resultados como anomalias. Um fator de contexto também pode ser um fator de prognóstico em alguns casos,e, quando isso ocorre, pode levar a um viés de confusão. Além disso, a avaliação da efetividade de programas de suporte para aderência de tratamento antiretroviral pode levar a resultados questionáveis se o determinantes sociais da saúde, tais como o acesso a água potável no sub-Saara africano ou a pobreza de pessoas de rua no Canadá não foram levadas em consideração.

context bias

A measurement bias due to the influence of the study's context on the interpretation of the study's results. Note: For example, when tests are done in groups with a high prevalence of a disease, readers are more likely to interpret the results as anomalies. A contextual factor may also be a prognostic factor in some cases, and, if that happens, it may lead to a confusion bias. Also, evaluation of the effectiveness of support programs for antiretroviral treatment compliance may lead to questionable results if social determinants of health, such as access to drinking water in sub-Saharan Africa or the poverty of homeless persons in Canada have not been taken into account.

\begin{center}\rule{0.5\linewidth}{0.5pt}\end{center}

viés de cruzamento

Viés de confusão que ocorre quando os elementos incluidos no grupo experimental não recebem a intervenção ou recebem outra intervenção, ou quando alguns elementos do grupo de controle recebem a intervenção. Nota: se os resultados desses elementos forem analisados com aqueles do seu grupo original, esse tipo de viés tenderá a ``diluir'' (diminuir) o efeito observado do tratamento.

crossover bias

A confusion bias that occurs when the subjects assigned to the experimental group do not receive the intervention or receive another intervention, or when some subjects in the control group receive the intervention. Note: If these subjects' results are analysed with those of their original group, this type of bias will tend to ``dilute'' (diminish) the observed treatment effect.

\begin{center}\rule{0.5\linewidth}{0.5pt}\end{center}

viés de defasamento temporal

Um tipo de viés devido a limites temporais de um método de pesquisa de literatura com possibilidade de afetar a inclusão de estudos numa revisão sistemática ou em qualquer outro tipo de síntese de literatura. Nota: tal viés pode ocorrer quando a direção (positiva ou negativa) ou o poder (significância estatística) de um estudo resulta no atraso de sua publicação.

time lag bias

A form of bias due to the temporal limits of a literature search method likely to affect the inclusion of studies in a systematic review or any other form of literature synthesis. Note: Such bias may occur when the direction (positive or negative) or the strength (statistical significance) of a study's results delays its publication. Alternate spelling: time-lag bias

\begin{center}\rule{0.5\linewidth}{0.5pt}\end{center}

Viés de detecção

Um viés de medida devido a diferenças sistemáticas entre os grupos sendo comparados e o modo como os desfechos foram apurados.

detection bias

A measurement bias due to systematic differences between the compared groups in how the outcomes are ascertained.

\begin{center}\rule{0.5\linewidth}{0.5pt}\end{center}

viés de execução

Um viés de confusão devido à diferenças sistemáticas nos cuidados recebidos fora da intervenção sendo avaliada. Nota: por exemplo, se os indivíduos sabem a que grupo pertencem, eles tendem a usar outras formas de cuidados. Se os que proporcionam os cuidados tiverem a mesma informação, eles podem trtar indivíduos diferentemente, de acordo com o grupo ao qual pertençam. O cegamento (tanto para os recebedores quanto para os provedores de cuidados) é utiliziado, entre outras medidas, como proteção contra o viés de execução.

performance bias

A confusion bias due to systematic differences in care received outside the intervention being evaluated. Note: For example, if subjects know which group they are in, they are more likely to use other forms of care. If care providers have the same information, they may treat the subjects differently, according to what group they are in. Blinding (both the recipients and the providers of care) is used to, among other things, protect against performance bias.

\begin{center}\rule{0.5\linewidth}{0.5pt}\end{center}

viés de lingua

Um tipo de viés relacionado a uma pesquisa restrita de literatura que exclui certos estudos, em função da língua em que foi publicada, numa revisão sistemática ou qualquer outra forma de síntese literária.

language bias

A form of bias linked to a restrictive literature search that excludes certain studies, because of the language in which they are published, from a systematic review or any other form of literature synthesis.

\begin{center}\rule{0.5\linewidth}{0.5pt}\end{center}

viés de medida

Um viés introduzido pela medida incorreta da exposição de um desfecho aos participantes de um estudo. Nota 1: o viés de medida pode ser introduzido pelo observador, pelos participantes do estudo ou pelas ferramentas de medida. Nota 2: sinônimos: incluem viés de informação e viés de observação.

measurement bias

A bias introduced by incorrect measurement of the exposure or outcome status of participants in a study. Note 1: Measurement bias can be introduced by the observer, the study participants or the measurement tools. Note 2: Synonyms include information bias and observation bias.

\begin{center}\rule{0.5\linewidth}{0.5pt}\end{center}

viés de publicação

Um viés devido a estudos sendo publicados, baseados na natureza e direção dos seus resultados. Exemplo: um estudo com resultados estatisticamente significantes favorecendo a intervenção de interesse pode ter maior possibilidade de ser publicado.

publication bias

A bias due to studies being published based on the nature and direction of their results. Example: A study with statistically significant results favouring the intervention of interest may have a greater likelihood of being published.

\begin{center}\rule{0.5\linewidth}{0.5pt}\end{center}

viés de seleção

Um tipo de viés devido a um erro na estimativa do efeito de uma intervenção pelo modo como os indivíduos foram selecionados na população de interesse. Nota: existe viés de selecão em estudos de coorte se o recrutamento dos indivíduos expostos (não expostos) for relacionado à presença da doença. Nos estudos de caso controle, o viés de seleção ocorre se o recrutamento dos casos (controle) for relacionado à exposição ao fator de interesse.

selection bias

A type of bias due to an error in the estimation of the effect of an intervention because of the way in which the subjects were selected in the population of interest. Note: There is selection bias in cohort studies if the recruitment of the subjects who are exposed (not exposed) is linked to the presence of the disease. In case-control studies, selection bias occurs if the recruitment of the cases (controls) is linked to exposure to the factor of interest.

\begin{center}\rule{0.5\linewidth}{0.5pt}\end{center}

viés de seleção de paciente

veja viés de seleção

patient selection bias

See selection bias.

\begin{center}\rule{0.5\linewidth}{0.5pt}\end{center}

viés de trabalho

Veja víés de verificação.

workup bias

See verification bias.

\begin{center}\rule{0.5\linewidth}{0.5pt}\end{center}

viés de verificação

É um tipo de viés de medição em que os resultados de uma prova de diagnóstico afetam se o procedimento padrão-ouro é usado para verificar o resultado do teste. Este tipo de viés também é conhecido como ``viés de referência'' ou ``viés de trabalho''

verification bias

See: workup bias.

\begin{center}\rule{0.5\linewidth}{0.5pt}\end{center}

viés do entrevistador

Um viés que resulta da ação de um entrevistador ao extrair, registrar ou interpretar informações de uma entrevista.

interviewer bias

A bias that results from the action of an interviewer when eliciting, recording or interpreting information from an interview.

\begin{center}\rule{0.5\linewidth}{0.5pt}\end{center}

vocabulário controlado

Um sistema de termos que possuem, entre outras características, definições, uma estrutura hierarquica e referências cruzadas, usadas para indexar e pesquisar o conjunto de dados numa base de dados bibliográfica, fatual ou em qualquer outra. Nota: por exemplo, o vocabulário controlado MeSH usado no MEDLINE e em outras bases de dados MEDLARS da United States National Library of Medicine

controlled vocabulary

A system of terms with, among other things, definitions, a hierarchical structure and cross-references, used to index and search a body of data in a bibliographic, factual or other database. Note: For example, the MeSH controlled vocabulary used in MEDLINE and the other MEDLARS databases of the United States National Library of Medicine.

\hypertarget{w}{%
\chapter*{W}\label{w}}
\addcontentsline{toc}{chapter}{W}

(vazio)

\begin{center}\rule{0.5\linewidth}{0.5pt}\end{center}

\hypertarget{x}{%
\chapter*{X}\label{x}}
\addcontentsline{toc}{chapter}{X}

(vazio)

\begin{center}\rule{0.5\linewidth}{0.5pt}\end{center}

\hypertarget{y}{%
\chapter*{Y}\label{y}}
\addcontentsline{toc}{chapter}{Y}

(vazio)

\begin{center}\rule{0.5\linewidth}{0.5pt}\end{center}

\hypertarget{z}{%
\chapter*{Z}\label{z}}
\addcontentsline{toc}{chapter}{Z}

(vazio)

\hypertarget{referuxeancias}{%
\chapter*{Referências}\label{referuxeancias}}
\addcontentsline{toc}{chapter}{Referências}

INAHTA\textsuperscript{1}
\url{https://www.inahta.org/}

Glossário da INAHTA\textsuperscript{2}
\url{http://htaglossary.net/List+all+terms}

Carlos Lederman\textsuperscript{3}
\url{http://lattes.cnpq.br/4968413275599200}

Wanderley Marques Bernardo\textsuperscript{4}
\url{http://lattes.cnpq.br/6173673567296321}

\hypertarget{licenciamento}{%
\chapter*{Licenciamento}\label{licenciamento}}
\addcontentsline{toc}{chapter}{Licenciamento}

\begin{center}\includegraphics[width=height=4in,width=3in]{licenca_isbn} \end{center}

\hypertarget{isbn}{%
\chapter*{ISBN}\label{isbn}}
\addcontentsline{toc}{chapter}{ISBN}

\begin{center}\includegraphics[width=height=4in,width=3in]{isbnv2} \end{center}

\hypertarget{ficha-catalogruxe1fica}{%
\chapter*{Ficha Catalográfica}\label{ficha-catalogruxe1fica}}
\addcontentsline{toc}{chapter}{Ficha Catalográfica}

  \bibliography{book.bib,packages.bib}

\end{document}
